\section{共同编年：征服怎样穿过不同区域（1206--1279）}

蒙古的扩张不是一条从草原直线抵达临安的箭头。每一次接收都要面对旧有的城防、粮区、道路、翻译者、工匠和地方组织；这些条件也使西夏、金、南宋、大理和海港所经历的战争并不相同。《元史》本纪和《元朝秘史》提供蒙古方面的重要年序，却有编纂层次、译名与宫廷立场；金、宋、西夏材料以及高丽、波斯文、碑刻、文书和考古材料并非旁注，而是辨认空间与代价所必需的对读者。

\subsection{先是西北与华北，不是预定的中国王朝（1206--1234）}

\begin{event}{\eventref{YUAN-1206-EVENT-151}{铁木真获成吉思汗称号}（一二〇六年；蒙古；漠北草原；称号可核，诸部整合的范围不宜夸定）}
《元史》相关本纪、《元朝秘史》以及波斯文史籍和《高丽史》的对应叙事，可交叉辨认称号与联盟重组的年序。不同语言、编纂层次与宫廷立场并不允许我们把所有诸部的服从方式、军事人数或随后征服的方向写成已经预定。
\end{event}

一二〇六年铁木真获成吉思汗称号。\hyperlink{evt:YUAN-1206-EVENT-151}{成吉思汗称号（1206）}联盟的重组使草原上的亲属、俘虏、商路与军事组织有了新的枢纽；它没有自动提供跨区域征收的制度。一二〇九年蒙古攻西夏，迫使西夏接受和议与贡赋安排。对河西而言，绿洲水源、城镇、寺院和穿越沙漠边缘的道路决定军队能停留多久，因而不能以一份和议概括实际控制。

\label{evt:mongol-jin-frontier}
一二一一年蒙古大举攻金，北方防线受到冲击；一二一八年又消灭西辽，进入中亚政治网络。金方《金史》、蒙古材料和中亚叙述所使用的地名与功绩语言并不相同，不能拼成一份无缝的战场记录。可以把握的是，金廷须在北方城防、仓储和道路日益受压时维持中都及南方边界；蒙古方面也要将被征服地区的工匠、粮食和投附者纳入不断移动的军政网络。

一二二七年西夏覆亡，同年成吉思汗去世；旧国家留下的河西通道和多语文书世界并未消散。一二三四年蒙古与南宋攻破蔡州，金朝终结。短暂协作发生在城池、粮道与河南故地的利益已经重叠之时，不能倒写成双方共有长期目标。金、宋、蒙三方对蔡州进军与责任的叙述各有位置，尤不宜将围城中的死亡、投降或破坏规模化为未经核验的数字。

\subsection{长江以前的多条前线（1235--1264）}

一二三五年蒙古贵族会议后，进攻南宋的行动在不同方向展开；这不是一条单一前线。四川山道、两淮水网、荆襄水陆交通和江南财赋的节奏各不相同。一二五三年大理被征服，南宋西南的战略环境随之改变；旧有地方权力、山地通道和西南的征发能力都要重新安排，不能把“大理已下”写成行政已经整齐接通。

一二五八年蒙哥发动大规模攻宋计划，次年钓鱼城战事期间去世，战争节奏与继承问题随之改变。宋、蒙史料对战事与个人死因的解释并不相同，较稳妥的说法是：四川前线的城防、山地运输与统帅层的变化共同限制了既有部署。随后一二六〇年忽必烈即位并同阿里不哥争位，到一二六四年主要继承战争结束；这几年首先是蒙古政治中心与资源调度重新组合，不能预先当作“元朝”已拥有稳定疆域。

\subsection{襄樊、临安与海上的终局（1268--1279）}

\label{evt:mongol-yuan-xiangfan}
一二六八年蒙古军围攻襄阳、樊城，围城延续至一二七三年。两城的意义在于汉水、荆襄道路、粮运、造船和江南财政在此相接；它不是一座孤立城池。宋、元材料都留下战事框架，却会以各自的攻守功绩组织叙述。攻城技术和水路封锁很重要，但沿线运粮、城内外居民与地方协作的材料并不均衡，不能把它们写成一种无摩擦的军事机器。

一二七一年忽必烈定国号元，是政治命名与机构调整的节点，不是对所有地区已经完成的描述。一二七四年元军分路南下，一二七六年进入临安、宋恭帝降元；朝廷降服、地方接收和遗民政权的存续随地点而不同。东南仍有宗室、将领、海商与地方军民在不同条件下行动。至一二七九年崖山海战后，南宋政权终结。宋末遗民诗文、元方军令和港口遗存能照出不同的失落、征服与迁徙经验，却没有任何一类材料可以替所有人在海上与沿岸所作的选择作证。

征服的制度得失在这一阶段已经显露：新政权获得了城镇、税源、船只、技术与行政人员，却必须以站赤、运输、赋役和地方中介使它们能够运作；沿线家庭、船户、军户与被迁徙者承担的成本，不会因为帝号更替而自动消失。下一章由此转入元朝如何试图把这些差异接入行省、路府、河运和海港的时间线。


\begin{event}{\eventanchor{YUAN-1218-EVENT-155}{西辽覆亡：蒙古进入中亚并不等于一条道路已经到手}}{一二一八年}{蒙古与西辽 · 七河、河中及通向中亚的城镇道路}{蒙古消灭西辽的主干可核；行军路径、城镇接收和不同群体的选择不能由征服叙事概括}
蒙古军消灭西辽时，获得的不是一条已经铺平的中亚走廊。七河、绿洲城镇、牧地和商道各有守备者、翻译者、税源与旧的政治承诺；军队要继续前进，必须把投附者、马匹、粮食和信息接入不断移动的网络。对西辽遗留的官员、商人和地方首领而言，改换保护者并不必然意味着同一种服从：有人可能在新军到来前转移，有人要在城门、市场和水源的控制改变后才调整关系。蒙古的行动把草原政治推入中亚，同时也使征服者面对远离本部的补给和行政选择，不能只用“进入中亚”四字盖过这一过程。

《元史》本纪提供基本年序，《元朝秘史》、波斯文史籍与《高丽史》的相关叙事可用来比较人名、路线和政治关系；它们皆带有宫廷、译名或后出编纂的层次，不能拼成无争议的行军图。多语碑刻和文书能显示局部城镇与身份安排，也不能替整个区域结算赋役和人口。西辽覆亡的直接后果是蒙古取得继续向中亚活动的条件；中期代价则是必须在新的城市、商路和地方中介中维持秩序。这条链也解释了为何后来的元朝不能被简单当作蒙古征服的自然终点。
\end{event}

% 逐条呈现账本事件；不把共享年序坐标当作具体行动。
\section{蒙古兴起与征服转换的逐年入口}
每一项均以其出处、行动者、空间线索和可追踪后果展开。材料没有说明地点、人数或地方反应时，正文保留这一限度。
\begin{textwindow}{本节材料怎样阅读}
《元朝秘史》以“成吉思合罕”书写统治者名号；它与《元史》保存蒙古统治者、诸部联盟和征服次序的不同叙述层。它们可校核行动者与相对先后，不能直接换算为军数、民众去留或控制深度。
\end{textwindow}

\subsection{1206年（元中统-53年；公元换算据《元史》本纪纪年）}
\noindent\eventanchor{SYD-1206-MONGOL-001}{阿剌忽思即以是謀報帝，居無何，舉部來歸}\textbf{蒙古。}
蒙古 在 1206年（元中统-53年；公元换算据《元史》本纪纪年） 的材料痕迹是〔阿剌忽思即以是謀報帝，居無何，舉部來歸〕，出处为 《元史》卷001〈本纪第1卷本年条〉。材料未将地点细化到城址，不能用中枢所在地替它补足。

前一层压力可由以下线索辨认：本纪年条提供日期和中枢可见行动；前因不据单条补定。 这一处置留下的后续条件是：可回查同年及后续条目，地方影响仍须另证。

《元史》为明初仓促官修，本纪保存大量实录条目但有编次与纪年问题；行省执行、数字、族称及对外关系须以条格、多语文书和对方史料互校。 因此，蒙古征服、草原政治与跨区域文献比较研究 只能扩展问题，不能代替未保存的细节。

\noindent\eventanchor{SYD-1206-MONGOL-053}{別里古台亦曰：「乃蠻欲奪我弧矢，是小我也，我輩義當同死}\textbf{蒙古。}
从 《元史》卷001〈本纪第1卷本年条〉 的 1206年（元中统-53年；公元换算据《元史》本纪纪年） 条目看，〔別里古台亦曰：「乃蠻欲奪我弧矢，是小我也，我輩義當同死〕使 蒙古 的一个具体行动进入叙事；材料未将地点细化到城址，不能用中枢所在地替它补足。

本纪年条提供日期和中枢可见行动；前因不据单条补定。 记录所见的结果则是：可回查同年及后续条目，地方影响仍须另证。 日期相邻不等于因果已经证实。

关于可说范围，须保留：《元史》为明初仓促官修，本纪保存大量实录条目但有编次与纪年问题；行省执行、数字、族称及对外关系须以条格、多语文书和对方史料互校。 进一步讨论可参照 蒙古征服、草原政治与跨区域文献比较研究

\noindent\eventanchor{SYD-1206-MONGOL-103}{不聽，各持馬乳橦疾鬭，奪忽兒真、火里真二哈敦以歸}\textbf{蒙古。}
1206年（元中统-53年；公元换算据《元史》本纪纪年） 的记载将 蒙古 与〔不聽，各持馬乳橦疾鬭，奪忽兒真、火里真二哈敦以歸〕相连。此处可确认的是这项被记录的行动；材料未将地点细化到城址，不能用中枢所在地替它补足。

把本项放回事件链，能够看到的前提为：本纪年条提供日期和中枢可见行动；前因不据单条补定。 而它所打开或收紧的下一步是：可回查同年及后续条目，地方影响仍须另证。

证据的边界是：《元史》为明初仓促官修，本纪保存大量实录条目但有编次与纪年问题；行省执行、数字、族称及对外关系须以条格、多语文书和对方史料互校。 蒙古征服、草原政治与跨区域文献比较研究 可用于继续检验这一结论。

\noindent\eventanchor{SYD-1206-MONGOL-150}{帝悅，曰：「以此眾戰，何憂不勝}\textbf{蒙古。}
〔帝悅，曰：「以此眾戰，何憂不勝〕见于 《元史》卷001〈本纪第1卷本年条〉 的 1206年（元中统-53年；公元换算据《元史》本纪纪年） 记录。蒙古 的选择因此有了可回查的时间位置；材料未将地点细化到城址，不能用中枢所在地替它补足。

它承接的条件是：本纪从朝廷视角记录军政行动；出兵与战果需分辨。 随后可追踪的变化为：后续路线、控制与损失应与对方记录和遗址材料复核。

《元史》为明初仓促官修，本纪保存大量实录条目但有编次与纪年问题；行省执行、数字、族称及对外关系须以条格、多语文书和对方史料互校。 因此，蒙古征服、草原政治与跨区域文献比较研究 只能扩展问题，不能代替未保存的细节。

\noindent\eventanchor{SYD-1206-MONGOL-192}{既而果勝，族人按彈、火察兒、答力台三人背約，帝怒，盡奪其所獲，分之軍中}\textbf{蒙古。}
1206年（元中统-53年；公元换算据《元史》本纪纪年），《元史》卷001〈本纪第1卷本年条〉 所见的〔既而果勝，族人按彈、火察兒、答力台三人背約，帝怒，盡奪其所獲，分之軍中〕把 蒙古 的处置留在可复核的年序中；材料未将地点细化到城址，不能用中枢所在地替它补足。

前一层压力可由以下线索辨认：本纪年条提供日期和中枢可见行动；前因不据单条补定。 这一处置留下的后续条件是：可回查同年及后续条目，地方影响仍须另证。

关于可说范围，须保留：《元史》为明初仓促官修，本纪保存大量实录条目但有编次与纪年问题；行省执行、数字、族称及对外关系须以条格、多语文书和对方史料互校。 进一步讨论可参照 蒙古征服、草原政治与跨区域文献比较研究

\noindent\eventanchor{SYD-1206-MONGOL-232}{其將火力速八赤對曰：「先王戰伐，勇進不回，馬尾人背，不使敵人見之}\textbf{蒙古。}
蒙古 在 1206年（元中统-53年；公元换算据《元史》本纪纪年） 的材料痕迹是〔其將火力速八赤對曰：「先王戰伐，勇進不回，馬尾人背，不使敵人見之〕，出处为 《元史》卷001〈本纪第1卷本年条〉。材料未将地点细化到城址，不能用中枢所在地替它补足。

本纪从朝廷视角记录军政行动；出兵与战果需分辨。 记录所见的结果则是：后续路线、控制与损失应与对方记录和遗址材料复核。 日期相邻不等于因果已经证实。

证据的边界是：《元史》为明初仓促官修，本纪保存大量实录条目但有编次与纪年问题；行省执行、数字、族称及对外关系须以条格、多语文书和对方史料互校。 蒙古征服、草原政治与跨区域文献比较研究 可用于继续检验这一结论。

\noindent\eventanchor{SYD-1206-MONGOL-265}{驅車徑出，輾傷諸兒，有至死者}\textbf{蒙古。}
从 《元史》卷001〈本纪第1卷本年条〉 的 1206年（元中统-53年；公元换算据《元史》本纪纪年） 条目看，〔驅車徑出，輾傷諸兒，有至死者〕使 蒙古 的一个具体行动进入叙事；材料未将地点细化到城址，不能用中枢所在地替它补足。

把本项放回事件链，能够看到的前提为：本纪年条提供日期和中枢可见行动；前因不据单条补定。 而它所打开或收紧的下一步是：可回查同年及后续条目，地方影响仍须另证。

《元史》为明初仓促官修，本纪保存大量实录条目但有编次与纪年问题；行省执行、数字、族称及对外关系须以条格、多语文书和对方史料互校。 因此，蒙古征服、草原政治与跨区域文献比较研究 只能扩展问题，不能代替未保存的细节。

\subsection{1206年}
\noindent\eventanchor{YUAN-1206-EVENT-151}{铁木真获成吉思汗称号，蒙古诸部统一}\textbf{蒙古。}
1206年 的记载将 蒙古 与〔铁木真获成吉思汗称号，蒙古诸部统一〕相连。此处可确认的是这项被记录的行动；材料未将地点细化到城址，不能用中枢所在地替它补足。

它承接的条件是：“铁木真获成吉思汗称号，蒙古诸部统一”发生在蒙古诸部整合及对西夏、金的战争的具体压力与机会之中，相关行动者的选择不能脱离此前事件链解释。 随后可追踪的变化为：“铁木真获成吉思汗称号，蒙古诸部统一”为后续政治、边防或资源配置留下条件；其社会影响与实际覆盖范围仍须以异类材料限定。

关于可说范围，须保留：《元史》明初速修，纪年与人名有异同；蒙古大汗政权不等同所有汗国，征服、赋役和身份分类须按地区及材料语言分列。 进一步讨论可参照 蒙古帝国、元朝行省财政、多语文书与区域社会研究

\subsection{1207年（元中统-52年；公元换算据《元史》本纪纪年）}
\noindent\eventanchor{SYD-1207-MONGOL-002}{既而野牒亦納里部、阿里替也兒部，皆遣使來獻名鷹}\textbf{蒙古。}
〔既而野牒亦納里部、阿里替也兒部，皆遣使來獻名鷹〕见于 《元史》卷001〈本纪第1卷本年条〉 的 1207年（元中统-52年；公元换算据《元史》本纪纪年） 记录。蒙古 的选择因此有了可回查的时间位置；材料未将地点细化到城址，不能用中枢所在地替它补足。

前一层压力可由以下线索辨认：本纪记录使节、贡赐或往来，不等于稳定统属关系。 这一处置留下的后续条件是：可与对方编年和交通、文书材料核对其后续落实。

证据的边界是：《元史》为明初仓促官修，本纪保存大量实录条目但有编次与纪年问题；行省执行、数字、族称及对外关系须以条格、多语文书和对方史料互校。 蒙古征服、草原政治与跨区域文献比较研究 可用于继续检验这一结论。

\noindent\eventanchor{SYD-1207-MONGOL-054}{是歲，遣按彈、不兀剌二人使乞力吉思}\textbf{蒙古。}
1207年（元中统-52年；公元换算据《元史》本纪纪年），《元史》卷001〈本纪第1卷本年条〉 所见的〔是歲，遣按彈、不兀剌二人使乞力吉思〕把 蒙古 的处置留在可复核的年序中；材料未将地点细化到城址，不能用中枢所在地替它补足。

本纪记录使节、贡赐或往来，不等于稳定统属关系。 记录所见的结果则是：可与对方编年和交通、文书材料核对其后续落实。 日期相邻不等于因果已经证实。

《元史》为明初仓促官修，本纪保存大量实录条目但有编次与纪年问题；行省执行、数字、族称及对外关系须以条格、多语文书和对方史料互校。 因此，蒙古征服、草原政治与跨区域文献比较研究 只能扩展问题，不能代替未保存的细节。

\subsection{1208年（元中统-51年；公元换算据《元史》本纪纪年）}
\noindent\eventanchor{SYD-1208-MONGOL-003}{時斡亦剌部等遇我前鋒，不戰而降，因用為向導}\textbf{蒙古。}
蒙古 在 1208年（元中统-51年；公元换算据《元史》本纪纪年） 的材料痕迹是〔時斡亦剌部等遇我前鋒，不戰而降，因用為向導〕，出处为 《元史》卷001〈本纪第1卷本年条〉。材料未将地点细化到城址，不能用中枢所在地替它补足。

把本项放回事件链，能够看到的前提为：本纪从朝廷视角记录军政行动；出兵与战果需分辨。 而它所打开或收紧的下一步是：后续路线、控制与损失应与对方记录和遗址材料复核。

关于可说范围，须保留：《元史》为明初仓促官修，本纪保存大量实录条目但有编次与纪年问题；行省执行、数字、族称及对外关系须以条格、多语文书和对方史料互校。 进一步讨论可参照 蒙古征服、草原政治与跨区域文献比较研究

\subsection{1209年（元中统-50年；公元换算据《元史》本纪纪年）}
\noindent\eventanchor{SYD-1209-MONGOL-004}{薄中興府，引河水灌之}\textbf{蒙古。}
从 《元史》卷001〈本纪第1卷本年条〉 的 1209年（元中统-50年；公元换算据《元史》本纪纪年） 条目看，〔薄中興府，引河水灌之〕使 蒙古 的一个具体行动进入叙事；题名保留的城、路、水道或机构名称，是目前可以落地的空间线索。

它承接的条件是：本纪保留灾害或工程的报告地点与朝廷处置；灾情范围需另证。 随后可追踪的变化为：可与地方文书、河工和环境材料比较，不将单点记录外推为全域因果。

证据的边界是：《元史》为明初仓促官修，本纪保存大量实录条目但有编次与纪年问题；行省执行、数字、族称及对外关系须以条格、多语文书和对方史料互校。 蒙古征服、草原政治与跨区域文献比较研究 可用于继续检验这一结论。

\noindent\eventanchor{SYD-1209-MONGOL-055}{堤決，水外潰，遂撤圍還}\textbf{蒙古。}
1209年（元中统-50年；公元换算据《元史》本纪纪年） 的记载将 蒙古 与〔堤決，水外潰，遂撤圍還〕相连。此处可确认的是这项被记录的行动；材料未将地点细化到城址，不能用中枢所在地替它补足。

前一层压力可由以下线索辨认：本纪保留灾害或工程的报告地点与朝廷处置；灾情范围需另证。 这一处置留下的后续条件是：可与地方文书、河工和环境材料比较，不将单点记录外推为全域因果。

《元史》为明初仓促官修，本纪保存大量实录条目但有编次与纪年问题；行省执行、数字、族称及对外关系须以条格、多语文书和对方史料互校。 因此，蒙古征服、草原政治与跨区域文献比较研究 只能扩展问题，不能代替未保存的细节。

\noindent\eventanchor{SYD-1209-MONGOL-104}{進至克夷門，復敗夏師，獲其將嵬名令公}\textbf{蒙古。}
〔進至克夷門，復敗夏師，獲其將嵬名令公〕见于 《元史》卷001〈本纪第1卷本年条〉 的 1209年（元中统-50年；公元换算据《元史》本纪纪年） 记录。蒙古 的选择因此有了可回查的时间位置；材料未将地点细化到城址，不能用中枢所在地替它补足。

本纪年条记录政令或机构处置；实施背景须参照制度材料。 记录所见的结果则是：可在同年及后续政令中追踪；条文不单独证明地方执行。 日期相邻不等于因果已经证实。

关于可说范围，须保留：《元史》为明初仓促官修，本纪保存大量实录条目但有编次与纪年问题；行省执行、数字、族称及对外关系须以条格、多语文书和对方史料互校。 进一步讨论可参照 蒙古征服、草原政治与跨区域文献比较研究

\noindent\eventanchor{SYD-1209-MONGOL-151}{遣太傅訛答入中興，招諭夏主，夏主納女請和}\textbf{蒙古与西夏。}
1209年（元中统-50年；公元换算据《元史》本纪纪年），《元史》卷001〈本纪第1卷本年条〉 所见的〔遣太傅訛答入中興，招諭夏主，夏主納女請和〕把 蒙古与西夏 的处置留在可复核的年序中；材料未将地点细化到城址，不能用中枢所在地替它补足。

把本项放回事件链，能够看到的前提为：本纪记录使节、贡赐或往来，不等于稳定统属关系。 而它所打开或收紧的下一步是：可与对方编年和交通、文书材料核对其后续落实。

证据的边界是：《元史》为明初仓促官修，本纪保存大量实录条目但有编次与纪年问题；行省执行、数字、族称及对外关系须以条格、多语文书和对方史料互校。 蒙古征服、草原政治与跨区域文献比较研究 可用于继续检验这一结论。

\noindent\eventanchor{SYD-1209-MONGOL-193}{四年己巳春，畏吾兒國來歸}\textbf{蒙古。}
蒙古 在 1209年（元中统-50年；公元换算据《元史》本纪纪年） 的材料痕迹是〔四年己巳春，畏吾兒國來歸〕，出处为 《元史》卷001〈本纪第1卷本年条〉。材料未将地点细化到城址，不能用中枢所在地替它补足。

它承接的条件是：本纪年条提供日期和中枢可见行动；前因不据单条补定。 随后可追踪的变化为：可回查同年及后续条目，地方影响仍须另证。

《元史》为明初仓促官修，本纪保存大量实录条目但有编次与纪年问题；行省执行、数字、族称及对外关系须以条格、多语文书和对方史料互校。 因此，蒙古征服、草原政治与跨区域文献比较研究 只能扩展问题，不能代替未保存的细节。

\noindent\eventanchor{SYD-1209-MONGOL-233}{夏主李安全遣其世子率師來戰，敗之，獲其副元帥高令公}\textbf{蒙古与西夏。}
从 《元史》卷001〈本纪第1卷本年条〉 的 1209年（元中统-50年；公元换算据《元史》本纪纪年） 条目看，〔夏主李安全遣其世子率師來戰，敗之，獲其副元帥高令公〕使 蒙古与西夏 的一个具体行动进入叙事；材料未将地点细化到城址，不能用中枢所在地替它补足。

前一层压力可由以下线索辨认：本纪从朝廷视角记录军政行动；出兵与战果需分辨。 这一处置留下的后续条件是：后续路线、控制与损失应与对方记录和遗址材料复核。

关于可说范围，须保留：《元史》为明初仓促官修，本纪保存大量实录条目但有编次与纪年问题；行省执行、数字、族称及对外关系须以条格、多语文书和对方史料互校。 进一步讨论可参照 蒙古征服、草原政治与跨区域文献比较研究

\subsection{1209年（蒙古成吉思汗四年；西夏年号、干支与月日待核；公元换算据《元史》与夏国相关材料年号对照）}
\noindent\eventanchor{YUAN-1209-EVENT-152}{蒙古军攻西夏并迫使其接受和议与贡赋}\textbf{蒙古与西夏。}
1209年（蒙古成吉思汗四年；西夏年号、干支与月日待核；公元换算据《元史》与夏国相关材料年号对照） 的记载将 蒙古与西夏 与〔蒙古军攻西夏并迫使其接受和议与贡赋〕相连。此处可确认的是这项被记录的行动；材料未将地点细化到城址，不能用中枢所在地替它补足。

蒙古诸部整合及对西夏、金的战争里的军事消耗、边境供给与使节往来构成谈判条件，促成“蒙古军攻西夏并迫使其接受和议与贡赋”所涉的协议或名义关系。 记录所见的结果则是：“蒙古军攻西夏并迫使其接受和议与贡赋”重排了贡赐、使节或停战的制度接口；条款是否执行及边地实际控制不能由盟约文字直接推定。 日期相邻不等于因果已经证实。

证据的边界是：《元史》明初速修，纪年与人名有异同；蒙古大汗政权不等同所有汗国，征服、赋役和身份分类须按地区及材料语言分列。 蒙古帝国、元朝行省财政、多语文书与区域社会研究 可用于继续检验这一结论。

\subsection{1210 年（共享年序桥接）}
\noindent\eventanchor{SY-1210-CHRONOLOGY-BRIDGE}{【C级年序桥接】保留1210 年的多政权并行检索入口；本条不主张所列政权在当年发生直接互动，具体事件须另以单项材料入账。}\textbf{南宋、金、西夏与蒙古（年序桥接）。}
〔【C级年序桥接】保留1210 年的多政权并行检索入口；本条不主张所列政权在当年发生直接互动，具体事件须另以单项材料入账。〕见于 《宋史》；《金史》；《元史》；蒙古文、波斯文及地方材料（据事核） 的 1210 年（共享年序桥接） 记录。南宋、金、西夏与蒙古（年序桥接） 的选择因此有了可回查的时间位置；题名保留的城、路、水道或机构名称，是目前可以落地的空间线索。

把本项放回事件链，能够看到的前提为：相邻已入账事件之间缺少该年度的共享年序入口，易使跨政权材料被单一政权叙事遮蔽。 而它所打开或收紧的下一步是：保留该年度的检索与补账位置；后续仅在获得可核单项材料时拆分为具体政治、财政、军事、社会或文化事件。

本条只确认该年位于可追踪的共享年序中；正史、地方文书和外部材料的保存密度不一，不能由桥接标签推断行动、统属、疆界或社会影响。 因此，蒙古征服、区域政权转换与跨语种年序研究 只能扩展问题，不能代替未保存的细节。

\subsection{1210年（元中统-49年；公元换算据《元史》本纪纪年）}
\noindent\eventanchor{SYD-1210-MONGOL-005}{初，帝貢歲幣于金，金主使衞王允濟貢於〔淨〕州}\textbf{蒙古与金。}
1210年（元中统-49年；公元换算据《元史》本纪纪年），《元史》卷001〈本纪第1卷本年条〉 所见的〔初，帝貢歲幣于金，金主使衞王允濟貢於〔淨〕州〕把 蒙古与金 的处置留在可复核的年序中；题名保留的城、路、水道或机构名称，是目前可以落地的空间线索。

它承接的条件是：本纪年条提供日期和中枢可见行动；前因不据单条补定。 随后可追踪的变化为：可回查同年及后续条目，地方影响仍须另证。

关于可说范围，须保留：《元史》为明初仓促官修，本纪保存大量实录条目但有编次与纪年问题；行省执行、数字、族称及对外关系须以条格、多语文书和对方史料互校。 进一步讨论可参照 蒙古征服、草原政治与跨区域文献比较研究

\noindent\eventanchor{SYD-1210-MONGOL-056}{帝命遮別襲殺其眾，遂略地而東}\textbf{蒙古。}
蒙古 在 1210年（元中统-49年；公元换算据《元史》本纪纪年） 的材料痕迹是〔帝命遮別襲殺其眾，遂略地而東〕，出处为 《元史》卷001〈本纪第1卷本年条〉。材料未将地点细化到城址，不能用中枢所在地替它补足。

前一层压力可由以下线索辨认：本纪年条记录政令或机构处置；实施背景须参照制度材料。 这一处置留下的后续条件是：可在同年及后续政令中追踪；条文不单独证明地方执行。

证据的边界是：《元史》为明初仓促官修，本纪保存大量实录条目但有编次与纪年问题；行省执行、数字、族称及对外关系须以条格、多语文书和对方史料互校。 蒙古征服、草原政治与跨区域文献比较研究 可用于继续检验这一结论。

\noindent\eventanchor{SYD-1210-MONGOL-105}{會金主璟殂，允濟嗣位，有詔至國，傳言當拜受}\textbf{蒙古与金。}
从 《元史》卷001〈本纪第1卷本年条〉 的 1210年（元中统-49年；公元换算据《元史》本纪纪年） 条目看，〔會金主璟殂，允濟嗣位，有詔至國，傳言當拜受〕使 蒙古与金 的一个具体行动进入叙事；材料未将地点细化到城址，不能用中枢所在地替它补足。

本纪年条记录政令或机构处置；实施背景须参照制度材料。 记录所见的结果则是：可在同年及后续政令中追踪；条文不单独证明地方执行。 日期相邻不等于因果已经证实。

《元史》为明初仓促官修，本纪保存大量实录条目但有编次与纪年问题；行省执行、数字、族称及对外关系须以条格、多语文书和对方史料互校。 因此，蒙古征服、草原政治与跨区域文献比较研究 只能扩展问题，不能代替未保存的细节。

\noindent\eventanchor{SYD-1210-MONGOL-152}{金使還言，允濟益怒，欲俟帝再入貢，就進場害之}\textbf{蒙古。}
1210年（元中统-49年；公元换算据《元史》本纪纪年） 的记载将 蒙古 与〔金使還言，允濟益怒，欲俟帝再入貢，就進場害之〕相连。此处可确认的是这项被记录的行动；材料未将地点细化到城址，不能用中枢所在地替它补足。

把本项放回事件链，能够看到的前提为：本纪记录使节、贡赐或往来，不等于稳定统属关系。 而它所打开或收紧的下一步是：可与对方编年和交通、文书材料核对其后续落实。

关于可说范围，须保留：《元史》为明初仓促官修，本纪保存大量实录条目但有编次与纪年问题；行省执行、数字、族称及对外关系须以条格、多语文书和对方史料互校。 进一步讨论可参照 蒙古征服、草原政治与跨区域文献比较研究

\noindent\eventanchor{SYD-1210-MONGOL-194}{五年庚午春，金謀來伐，築烏沙堡}\textbf{蒙古。}
〔五年庚午春，金謀來伐，築烏沙堡〕见于 《元史》卷001〈本纪第1卷本年条〉 的 1210年（元中统-49年；公元换算据《元史》本纪纪年） 记录。蒙古 的选择因此有了可回查的时间位置；材料未将地点细化到城址，不能用中枢所在地替它补足。

它承接的条件是：本纪年条提供日期和中枢可见行动；前因不据单条补定。 随后可追踪的变化为：可回查同年及后续条目，地方影响仍须另证。

证据的边界是：《元史》为明初仓促官修，本纪保存大量实录条目但有编次与纪年问题；行省执行、数字、族称及对外关系须以条格、多语文书和对方史料互校。 蒙古征服、草原政治与跨区域文献比较研究 可用于继续检验这一结论。

\subsection{1211年（元中统-48年；公元换算据《元史》本纪纪年）}
\noindent\eventanchor{SYD-1211-MONGOL-006}{八月，帝及金師戰于宣平之會河川，敗之}\textbf{蒙古。}
1211年（元中统-48年；公元换算据《元史》本纪纪年），《元史》卷001〈本纪第1卷本年条〉 所见的〔八月，帝及金師戰于宣平之會河川，敗之〕把 蒙古 的处置留在可复核的年序中；题名保留的城、路、水道或机构名称，是目前可以落地的空间线索。

前一层压力可由以下线索辨认：本纪保留灾害或工程的报告地点与朝廷处置；灾情范围需另证。 这一处置留下的后续条件是：可与地方文书、河工和环境材料比较，不将单点记录外推为全域因果。

《元史》为明初仓促官修，本纪保存大量实录条目但有编次与纪年问题；行省执行、数字、族称及对外关系须以条格、多语文书和对方史料互校。 因此，蒙古征服、草原政治与跨区域文献比较研究 只能扩展问题，不能代替未保存的细节。

\noindent\eventanchor{SYD-1211-MONGOL-057}{二月，帝自將南伐，敗金將定薛於野狐嶺，取大水濼、豐利等縣}\textbf{蒙古。}
蒙古 在 1211年（元中统-48年；公元换算据《元史》本纪纪年） 的材料痕迹是〔二月，帝自將南伐，敗金將定薛於野狐嶺，取大水濼、豐利等縣〕，出处为 《元史》卷001〈本纪第1卷本年条〉。题名保留的城、路、水道或机构名称，是目前可以落地的空间线索。

本纪保留灾害或工程的报告地点与朝廷处置；灾情范围需另证。 记录所见的结果则是：可与地方文书、河工和环境材料比较，不将单点记录外推为全域因果。 日期相邻不等于因果已经证实。

关于可说范围，须保留：《元史》为明初仓促官修，本纪保存大量实录条目但有编次与纪年问题；行省执行、数字、族称及对外关系须以条格、多语文书和对方史料互校。 进一步讨论可参照 蒙古征服、草原政治与跨区域文献比较研究

\noindent\eventanchor{SYD-1211-MONGOL-106}{劉伯林、夾谷長哥等來降}\textbf{蒙古。}
从 《元史》卷001〈本纪第1卷本年条〉 的 1211年（元中统-48年；公元换算据《元史》本纪纪年） 条目看，〔劉伯林、夾谷長哥等來降〕使 蒙古 的一个具体行动进入叙事；材料未将地点细化到城址，不能用中枢所在地替它补足。

把本项放回事件链，能够看到的前提为：本纪从朝廷视角记录军政行动；出兵与战果需分辨。 而它所打开或收紧的下一步是：后续路线、控制与损失应与对方记录和遗址材料复核。

证据的边界是：《元史》为明初仓促官修，本纪保存大量实录条目但有编次与纪年问题；行省执行、数字、族称及对外关系须以条格、多语文书和对方史料互校。 蒙古征服、草原政治与跨区域文献比较研究 可用于继续检验这一结论。

\noindent\eventanchor{SYD-1211-MONGOL-153}{秋七月，命遮別攻烏沙堡及烏月營，拔之}\textbf{蒙古。}
1211年（元中统-48年；公元换算据《元史》本纪纪年） 的记载将 蒙古 与〔秋七月，命遮別攻烏沙堡及烏月營，拔之〕相连。此处可确认的是这项被记录的行动；材料未将地点细化到城址，不能用中枢所在地替它补足。

它承接的条件是：本纪年条记录政令或机构处置；实施背景须参照制度材料。 随后可追踪的变化为：可在同年及后续政令中追踪；条文不单独证明地方执行。

《元史》为明初仓促官修，本纪保存大量实录条目但有编次与纪年问题；行省执行、数字、族称及对外关系须以条格、多语文书和对方史料互校。 因此，蒙古征服、草原政治与跨区域文献比较研究 只能扩展问题，不能代替未保存的细节。

\noindent\eventanchor{SYD-1211-MONGOL-195}{畏吾兒國主亦都護來覲}\textbf{蒙古。}
〔畏吾兒國主亦都護來覲〕见于 《元史》卷001〈本纪第1卷本年条〉 的 1211年（元中统-48年；公元换算据《元史》本纪纪年） 记录。蒙古 的选择因此有了可回查的时间位置；题名保留的城、路、水道或机构名称，是目前可以落地的空间线索。

前一层压力可由以下线索辨认：本纪年条提供日期和中枢可见行动；前因不据单条补定。 这一处置留下的后续条件是：可回查同年及后续条目，地方影响仍须另证。

关于可说范围，须保留：《元史》为明初仓促官修，本纪保存大量实录条目但有编次与纪年问题；行省执行、数字、族称及对外关系须以条格、多语文书和对方史料互校。 进一步讨论可参照 蒙古征服、草原政治与跨区域文献比较研究

\noindent\eventanchor{SYD-1211-MONGOL-234}{西域哈剌魯部主阿昔蘭罕來降}\textbf{蒙古。}
1211年（元中统-48年；公元换算据《元史》本纪纪年），《元史》卷001〈本纪第1卷本年条〉 所见的〔西域哈剌魯部主阿昔蘭罕來降〕把 蒙古 的处置留在可复核的年序中；材料未将地点细化到城址，不能用中枢所在地替它补足。

本纪从朝廷视角记录军政行动；出兵与战果需分辨。 记录所见的结果则是：后续路线、控制与损失应与对方记录和遗址材料复核。 日期相邻不等于因果已经证实。

证据的边界是：《元史》为明初仓促官修，本纪保存大量实录条目但有编次与纪年问题；行省执行、数字、族称及对外关系须以条格、多语文书和对方史料互校。 蒙古征服、草原政治与跨区域文献比较研究 可用于继续检验这一结论。

\noindent\eventanchor{SYD-1211-MONGOL-266}{耶律阿海降，入見帝于行在所}\textbf{蒙古。}
蒙古 在 1211年（元中统-48年；公元换算据《元史》本纪纪年） 的材料痕迹是〔耶律阿海降，入見帝于行在所〕，出处为 《元史》卷001〈本纪第1卷本年条〉。题名保留的城、路、水道或机构名称，是目前可以落地的空间线索。

把本项放回事件链，能够看到的前提为：本纪从朝廷视角记录军政行动；出兵与战果需分辨。 而它所打开或收紧的下一步是：后续路线、控制与损失应与对方记录和遗址材料复核。

《元史》为明初仓促官修，本纪保存大量实录条目但有编次与纪年问题；行省执行、数字、族称及对外关系须以条格、多语文书和对方史料互校。 因此，蒙古征服、草原政治与跨区域文献比较研究 只能扩展问题，不能代替未保存的细节。

\subsection{1211年（蒙古成吉思汗六年、金大安三年、宋嘉定四年；干支、月日待核；公元换算据各方本纪年号对照）}
\noindent\eventanchor{YUAN-1211-EVENT-153}{蒙古大举进攻金朝北方边境}\textbf{蒙古与金。}
从 《元史》相关本纪；《元朝秘史》、波斯文史籍或《高丽史》对应叙事；《元典章》《通制条格》、多语碑刻与文书（据事核） 的 1211年（蒙古成吉思汗六年、金大安三年、宋嘉定四年；干支、月日待核；公元换算据各方本纪年号对照） 条目看，〔蒙古大举进攻金朝北方边境〕使 蒙古与金 的一个具体行动进入叙事；材料未将地点细化到城址，不能用中枢所在地替它补足。

它承接的条件是：蒙古诸部整合及对西夏、金的战争中既有的联盟、交通与补给压力，为“蒙古大举进攻金朝北方边境”提供了直接的军事或动员背景。 随后可追踪的变化为：“蒙古大举进攻金朝北方边境”改变了相关城镇、道路或政权间的军事选择；伤亡、俘获和控制范围仅可按各方材料分别确认。

关于可说范围，须保留：《元史》明初速修，纪年与人名有异同；蒙古大汗政权不等同所有汗国，征服、赋役和身份分类须按地区及材料语言分列。 进一步讨论可参照 蒙古帝国、元朝行省财政、多语文书与区域社会研究

\subsection{1212 年（共享年序桥接）}
\noindent\eventanchor{SY-1212-CHRONOLOGY-BRIDGE}{【C级年序桥接】保留1212 年的多政权并行检索入口；本条不主张所列政权在当年发生直接互动，具体事件须另以单项材料入账。}\textbf{南宋、金、西夏与蒙古（年序桥接）。}
1212 年（共享年序桥接） 的记载将 南宋、金、西夏与蒙古（年序桥接） 与〔【C级年序桥接】保留1212 年的多政权并行检索入口；本条不主张所列政权在当年发生直接互动，具体事件须另以单项材料入账。〕相连。此处可确认的是这项被记录的行动；题名保留的城、路、水道或机构名称，是目前可以落地的空间线索。

前一层压力可由以下线索辨认：相邻已入账事件之间缺少该年度的共享年序入口，易使跨政权材料被单一政权叙事遮蔽。 这一处置留下的后续条件是：保留该年度的检索与补账位置；后续仅在获得可核单项材料时拆分为具体政治、财政、军事、社会或文化事件。

证据的边界是：本条只确认该年位于可追踪的共享年序中；正史、地方文书和外部材料的保存密度不一，不能由桥接标签推断行动、统属、疆界或社会影响。 蒙古征服、区域政权转换与跨语种年序研究 可用于继续检验这一结论。

\subsection{1212年（元中统-47年；公元换算据《元史》本纪纪年）}
\noindent\eventanchor{SYD-1212-MONGOL-007}{冬十二月甲申，遮別攻東京不拔，即引去，夜馳還，襲克之}\textbf{蒙古。}
〔冬十二月甲申，遮別攻東京不拔，即引去，夜馳還，襲克之〕见于 《元史》卷001〈本纪第1卷本年条〉 的 1212年（元中统-47年；公元换算据《元史》本纪纪年） 记录。蒙古 的选择因此有了可回查的时间位置；题名保留的城、路、水道或机构名称，是目前可以落地的空间线索。

本纪从朝廷视角记录军政行动；出兵与战果需分辨。 记录所见的结果则是：后续路线、控制与损失应与对方记录和遗址材料复核。 日期相邻不等于因果已经证实。

《元史》为明初仓促官修，本纪保存大量实录条目但有编次与纪年问题；行省执行、数字、族称及对外关系须以条格、多语文书和对方史料互校。 因此，蒙古征服、草原政治与跨区域文献比较研究 只能扩展问题，不能代替未保存的细节。

\noindent\eventanchor{SYD-1212-MONGOL-058}{復攻西京，帝中流矢，遂撤圍}\textbf{蒙古。}
1212年（元中统-47年；公元换算据《元史》本纪纪年），《元史》卷001〈本纪第1卷本年条〉 所见的〔復攻西京，帝中流矢，遂撤圍〕把 蒙古 的处置留在可复核的年序中；题名保留的城、路、水道或机构名称，是目前可以落地的空间线索。

把本项放回事件链，能够看到的前提为：本纪年条记录政令或机构处置；实施背景须参照制度材料。 而它所打开或收紧的下一步是：可在同年及后续政令中追踪；条文不单独证明地方执行。

关于可说范围，须保留：《元史》为明初仓促官修，本纪保存大量实录条目但有编次与纪年问题；行省执行、数字、族称及对外关系须以条格、多语文书和对方史料互校。 进一步讨论可参照 蒙古征服、草原政治与跨区域文献比较研究

\noindent\eventanchor{SYD-1212-MONGOL-107}{金將紇石烈九斤等率兵三十萬來援，帝與戰于獾兒觜，大敗之}\textbf{蒙古。}
蒙古 在 1212年（元中统-47年；公元换算据《元史》本纪纪年） 的材料痕迹是〔金將紇石烈九斤等率兵三十萬來援，帝與戰于獾兒觜，大敗之〕，出处为 《元史》卷001〈本纪第1卷本年条〉。材料未将地点细化到城址，不能用中枢所在地替它补足。

它承接的条件是：本纪从朝廷视角记录军政行动；出兵与战果需分辨。 随后可追踪的变化为：后续路线、控制与损失应与对方记录和遗址材料复核。

证据的边界是：《元史》为明初仓促官修，本纪保存大量实录条目但有编次与纪年问题；行省执行、数字、族称及对外关系须以条格、多语文书和对方史料互校。 蒙古征服、草原政治与跨区域文献比较研究 可用于继续检验这一结论。

\noindent\eventanchor{SYD-1212-MONGOL-154}{金元帥左都監奧屯襄率師來援，帝遣兵誘至密谷口逆擊之，盡殪}\textbf{蒙古。}
从 《元史》卷001〈本纪第1卷本年条〉 的 1212年（元中统-47年；公元换算据《元史》本纪纪年） 条目看，〔金元帥左都監奧屯襄率師來援，帝遣兵誘至密谷口逆擊之，盡殪〕使 蒙古 的一个具体行动进入叙事；题名保留的城、路、水道或机构名称，是目前可以落地的空间线索。

前一层压力可由以下线索辨认：本纪从朝廷视角记录军政行动；出兵与战果需分辨。 这一处置留下的后续条件是：后续路线、控制与损失应与对方记录和遗址材料复核。

《元史》为明初仓促官修，本纪保存大量实录条目但有编次与纪年问题；行省执行、数字、族称及对外关系须以条格、多语文书和对方史料互校。 因此，蒙古征服、草原政治与跨区域文献比较研究 只能扩展问题，不能代替未保存的细节。

\noindent\eventanchor{SYD-1212-MONGOL-196}{七年壬申春正月，耶律留哥聚眾于隆安，自為都元帥，遣使來附}\textbf{蒙古。}
1212年（元中统-47年；公元换算据《元史》本纪纪年） 的记载将 蒙古 与〔七年壬申春正月，耶律留哥聚眾于隆安，自為都元帥，遣使來附〕相连。此处可确认的是这项被记录的行动；题名保留的城、路、水道或机构名称，是目前可以落地的空间线索。

本纪记录使节、贡赐或往来，不等于稳定统属关系。 记录所见的结果则是：可与对方编年和交通、文书材料核对其后续落实。 日期相邻不等于因果已经证实。

关于可说范围，须保留：《元史》为明初仓促官修，本纪保存大量实录条目但有编次与纪年问题；行省执行、数字、族称及对外关系须以条格、多语文书和对方史料互校。 进一步讨论可参照 蒙古征服、草原政治与跨区域文献比较研究

\subsection{1213年（元中统-46年；公元换算据《元史》本纪纪年）}
\noindent\eventanchor{SYD-1213-MONGOL-008}{八年癸酉春，耶律留哥自立為遼王，改元元統}\textbf{蒙古。}
〔八年癸酉春，耶律留哥自立為遼王，改元元統〕见于 《元史》卷001〈本纪第1卷本年条〉 的 1213年（元中统-46年；公元换算据《元史》本纪纪年） 记录。蒙古 的选择因此有了可回查的时间位置；材料未将地点细化到城址，不能用中枢所在地替它补足。

把本项放回事件链，能够看到的前提为：本纪年条记录政令或机构处置；实施背景须参照制度材料。 而它所打开或收紧的下一步是：可在同年及后续政令中追踪；条文不单独证明地方执行。

证据的边界是：《元史》为明初仓促官修，本纪保存大量实录条目但有编次与纪年问题；行省执行、数字、族称及对外关系须以条格、多语文书和对方史料互校。 蒙古征服、草原政治与跨区域文献比较研究 可用于继续检验这一结论。

\noindent\eventanchor{SYD-1213-MONGOL-059}{八月，金忽沙虎弒其主允濟，迎豐王珣立之}\textbf{蒙古。}
1213年（元中统-46年；公元换算据《元史》本纪纪年），《元史》卷001〈本纪第1卷本年条〉 所见的〔八月，金忽沙虎弒其主允濟，迎豐王珣立之〕把 蒙古 的处置留在可复核的年序中；材料未将地点细化到城址，不能用中枢所在地替它补足。

它承接的条件是：本纪年条提供日期和中枢可见行动；前因不据单条补定。 随后可追踪的变化为：可回查同年及后续条目，地方影响仍须另证。

《元史》为明初仓促官修，本纪保存大量实录条目但有编次与纪年问题；行省执行、数字、族称及对外关系须以条格、多语文书和对方史料互校。 因此，蒙古征服、草原政治与跨区域文献比较研究 只能扩展问题，不能代替未保存的细节。

\noindent\eventanchor{SYD-1213-MONGOL-108}{帝出紫荊關，敗金師于五回嶺，拔涿、易二州}\textbf{蒙古。}
蒙古 在 1213年（元中统-46年；公元换算据《元史》本纪纪年） 的材料痕迹是〔帝出紫荊關，敗金師于五回嶺，拔涿、易二州〕，出处为 《元史》卷001〈本纪第1卷本年条〉。题名保留的城、路、水道或机构名称，是目前可以落地的空间线索。

前一层压力可由以下线索辨认：本纪从朝廷视角记录军政行动；出兵与战果需分辨。 这一处置留下的后续条件是：后续路线、控制与损失应与对方记录和遗址材料复核。

关于可说范围，须保留：《元史》为明初仓促官修，本纪保存大量实录条目但有编次与纪年问题；行省执行、数字、族称及对外关系须以条格、多语文书和对方史料互校。 进一步讨论可参照 蒙古征服、草原政治与跨区域文献比较研究

\noindent\eventanchor{SYD-1213-MONGOL-155}{帝與皇子拖雷為中軍，取雄、霸、莫、安、河間、滄、景、獻、深、祁、蠡、冀、恩、濮、開、滑、博、濟、泰安、濟南、濱、棣、益都、淄、濰、登、萊、沂等郡}\textbf{蒙古。}
从 《元史》卷001〈本纪第1卷本年条〉 的 1213年（元中统-46年；公元换算据《元史》本纪纪年） 条目看，〔帝與皇子拖雷為中軍，取雄、霸、莫、安、河間、滄、景、獻、深、祁、蠡、冀、恩、濮、開、滑、博、濟、泰安、濟南、濱、棣、益都、淄、濰、登、萊、沂等郡〕使 蒙古 的一个具体行动进入叙事；题名保留的城、路、水道或机构名称，是目前可以落地的空间线索。

本纪保留灾害或工程的报告地点与朝廷处置；灾情范围需另证。 记录所见的结果则是：可与地方文书、河工和环境材料比较，不将单点记录外推为全域因果。 日期相邻不等于因果已经证实。

证据的边界是：《元史》为明初仓促官修，本纪保存大量实录条目但有编次与纪年问题；行省执行、数字、族称及对外关系须以条格、多语文书和对方史料互校。 蒙古征服、草原政治与跨区域文献比较研究 可用于继续检验这一结论。

\noindent\eventanchor{SYD-1213-MONGOL-197}{復命木華黎攻密州，屠之}\textbf{蒙古。}
1213年（元中统-46年；公元换算据《元史》本纪纪年） 的记载将 蒙古 与〔復命木華黎攻密州，屠之〕相连。此处可确认的是这项被记录的行动；题名保留的城、路、水道或机构名称，是目前可以落地的空间线索。

把本项放回事件链，能够看到的前提为：本纪年条记录政令或机构处置；实施背景须参照制度材料。 而它所打开或收紧的下一步是：可在同年及后续政令中追踪；条文不单独证明地方执行。

《元史》为明初仓促官修，本纪保存大量实录条目但有编次与纪年问题；行省执行、数字、族称及对外关系须以条格、多语文书和对方史料互校。 因此，蒙古征服、草原政治与跨区域文献比较研究 只能扩展问题，不能代替未保存的细节。

\noindent\eventanchor{SYD-1213-MONGOL-235}{及金行省完顏綱、元帥高琪戰，敗之，追至北口}\textbf{蒙古。}
〔及金行省完顏綱、元帥高琪戰，敗之，追至北口〕见于 《元史》卷001〈本纪第1卷本年条〉 的 1213年（元中统-46年；公元换算据《元史》本纪纪年） 记录。蒙古 的选择因此有了可回查的时间位置；材料未将地点细化到城址，不能用中枢所在地替它补足。

它承接的条件是：本纪年条记录政令或机构处置；实施背景须参照制度材料。 随后可追踪的变化为：可在同年及后续政令中追踪；条文不单独证明地方执行。

关于可说范围，须保留：《元史》为明初仓促官修，本纪保存大量实录条目但有编次与纪年问题；行省执行、数字、族称及对外关系须以条格、多语文书和对方史料互校。 进一步讨论可参照 蒙古征服、草原政治与跨区域文献比较研究

\noindent\eventanchor{SYD-1213-MONGOL-267}{金兵保居庸，詔可忒、薄剎守之}\textbf{蒙古与金。}
1213年（元中统-46年；公元换算据《元史》本纪纪年），《元史》卷001〈本纪第1卷本年条〉 所见的〔金兵保居庸，詔可忒、薄剎守之〕把 蒙古与金 的处置留在可复核的年序中；材料未将地点细化到城址，不能用中枢所在地替它补足。

前一层压力可由以下线索辨认：本纪年条记录政令或机构处置；实施背景须参照制度材料。 这一处置留下的后续条件是：可在同年及后续政令中追踪；条文不单独证明地方执行。

证据的边界是：《元史》为明初仓促官修，本纪保存大量实录条目但有编次与纪年问题；行省执行、数字、族称及对外关系须以条格、多语文书和对方史料互校。 蒙古征服、草原政治与跨区域文献比较研究 可用于继续检验这一结论。

\subsection{1214 年（共享年序桥接）}
\noindent\eventanchor{SY-1214-CHRONOLOGY-BRIDGE}{【C级年序桥接】保留1214 年的多政权并行检索入口；本条不主张所列政权在当年发生直接互动，具体事件须另以单项材料入账。}\textbf{南宋、金、西夏与蒙古（年序桥接）。}
南宋、金、西夏与蒙古（年序桥接） 在 1214 年（共享年序桥接） 的材料痕迹是〔【C级年序桥接】保留1214 年的多政权并行检索入口；本条不主张所列政权在当年发生直接互动，具体事件须另以单项材料入账。〕，出处为 《宋史》；《金史》；《元史》；蒙古文、波斯文及地方材料（据事核）。题名保留的城、路、水道或机构名称，是目前可以落地的空间线索。

相邻已入账事件之间缺少该年度的共享年序入口，易使跨政权材料被单一政权叙事遮蔽。 记录所见的结果则是：保留该年度的检索与补账位置；后续仅在获得可核单项材料时拆分为具体政治、财政、军事、社会或文化事件。 日期相邻不等于因果已经证实。

本条只确认该年位于可追踪的共享年序中；正史、地方文书和外部材料的保存密度不一，不能由桥接标签推断行动、统属、疆界或社会影响。 因此，蒙古征服、区域政权转换与跨语种年序研究 只能扩展问题，不能代替未保存的细节。

\subsection{1214年（元中统-45年；公元换算据《元史》本纪纪年）}
\noindent\eventanchor{SYD-1214-MONGOL-009}{冬十月，木華黎征遼東，高州盧琮、金〔朴〕等降}\textbf{蒙古。}
从 《元史》卷001〈本纪第1卷本年条〉 的 1214年（元中统-45年；公元换算据《元史》本纪纪年） 条目看，〔冬十月，木華黎征遼東，高州盧琮、金〔朴〕等降〕使 蒙古 的一个具体行动进入叙事；题名保留的城、路、水道或机构名称，是目前可以落地的空间线索。

把本项放回事件链，能够看到的前提为：本纪从朝廷视角记录军政行动；出兵与战果需分辨。 而它所打开或收紧的下一步是：后续路线、控制与损失应与对方记录和遗址材料复核。

关于可说范围，须保留：《元史》为明初仓促官修，本纪保存大量实录条目但有编次与纪年问题；行省执行、数字、族称及对外关系须以条格、多语文书和对方史料互校。 进一步讨论可参照 蒙古征服、草原政治与跨区域文献比较研究

\noindent\eventanchor{SYD-1214-MONGOL-060}{錦州張鯨殺其節度使，自立為臨海王，遣使來降}\textbf{蒙古。}
1214年（元中统-45年；公元换算据《元史》本纪纪年） 的记载将 蒙古 与〔錦州張鯨殺其節度使，自立為臨海王，遣使來降〕相连。此处可确认的是这项被记录的行动；题名保留的城、路、水道或机构名称，是目前可以落地的空间线索。

它承接的条件是：本纪从朝廷视角记录军政行动；出兵与战果需分辨。 随后可追踪的变化为：后续路线、控制与损失应与对方记录和遗址材料复核。

证据的边界是：《元史》为明初仓促官修，本纪保存大量实录条目但有编次与纪年问题；行省执行、数字、族称及对外关系须以条格、多语文书和对方史料互校。 蒙古征服、草原政治与跨区域文献比较研究 可用于继续检验这一结论。

\noindent\eventanchor{SYD-1214-MONGOL-109}{六月，金乣軍斫答等殺其主帥，率眾來降}\textbf{蒙古。}
〔六月，金乣軍斫答等殺其主帥，率眾來降〕见于 《元史》卷001〈本纪第1卷本年条〉 的 1214年（元中统-45年；公元换算据《元史》本纪纪年） 记录。蒙古 的选择因此有了可回查的时间位置；材料未将地点细化到城址，不能用中枢所在地替它补足。

前一层压力可由以下线索辨认：本纪从朝廷视角记录军政行动；出兵与战果需分辨。 这一处置留下的后续条件是：后续路线、控制与损失应与对方记录和遗址材料复核。

《元史》为明初仓促官修，本纪保存大量实录条目但有编次与纪年问题；行省执行、数字、族称及对外关系须以条格、多语文书和对方史料互校。 因此，蒙古征服、草原政治与跨区域文献比较研究 只能扩展问题，不能代替未保存的细节。

\noindent\eventanchor{SYD-1214-MONGOL-156}{乃遣使諭金主曰：「汝山東、河北郡縣悉為我有，汝所守惟燕京耳}\textbf{蒙古与金。}
1214年（元中统-45年；公元换算据《元史》本纪纪年），《元史》卷001〈本纪第1卷本年条〉 所见的〔乃遣使諭金主曰：「汝山東、河北郡縣悉為我有，汝所守惟燕京耳〕把 蒙古与金 的处置留在可复核的年序中；题名保留的城、路、水道或机构名称，是目前可以落地的空间线索。

本纪保留灾害或工程的报告地点与朝廷处置；灾情范围需另证。 记录所见的结果则是：可与地方文书、河工和环境材料比较，不将单点记录外推为全域因果。 日期相邻不等于因果已经证实。

关于可说范围，须保留：《元史》为明初仓促官修，本纪保存大量实录条目但有编次与纪年问题；行省执行、数字、族称及对外关系须以条格、多语文书和对方史料互校。 进一步讨论可参照 蒙古征服、草原政治与跨区域文献比较研究

\noindent\eventanchor{SYD-1214-MONGOL-198}{天既弱汝，我復迫汝於險，天其謂我何}\textbf{蒙古。}
蒙古 在 1214年（元中统-45年；公元换算据《元史》本纪纪年） 的材料痕迹是〔天既弱汝，我復迫汝於險，天其謂我何〕，出处为 《元史》卷001〈本纪第1卷本年条〉。材料未将地点细化到城址，不能用中枢所在地替它补足。

把本项放回事件链，能够看到的前提为：本纪年条记录政令或机构处置；实施背景须参照制度材料。 而它所打开或收紧的下一步是：可在同年及后续政令中追踪；条文不单独证明地方执行。

证据的边界是：《元史》为明初仓促官修，本纪保存大量实录条目但有编次与纪年问题；行省执行、数字、族称及对外关系须以条格、多语文书和对方史料互校。 蒙古征服、草原政治与跨区域文献比较研究 可用于继续检验这一结论。

\noindent\eventanchor{SYD-1214-MONGOL-236}{我今還軍，汝不能犒師以弭我諸將之怒耶？」金主遂遣使求和，奉衞紹王女岐國公主及金帛、童男女五百、馬三千以獻，仍遣其丞相完顏福興送帝出居庸}\textbf{蒙古与金。}
从 《元史》卷001〈本纪第1卷本年条〉 的 1214年（元中统-45年；公元换算据《元史》本纪纪年） 条目看，〔我今還軍，汝不能犒師以弭我諸將之怒耶？」金主遂遣使求和，奉衞紹王女岐國公主及金帛、童男女五百、馬三千以獻，仍遣其丞相完顏福興送帝出居庸〕使 蒙古与金 的一个具体行动进入叙事；材料未将地点细化到城址，不能用中枢所在地替它补足。

它承接的条件是：本纪记录使节、贡赐或往来，不等于稳定统属关系。 随后可追踪的变化为：可与对方编年和交通、文书材料核对其后续落实。

《元史》为明初仓促官修，本纪保存大量实录条目但有编次与纪年问题；行省执行、数字、族称及对外关系须以条格、多语文书和对方史料互校。 因此，蒙古征服、草原政治与跨区域文献比较研究 只能扩展问题，不能代替未保存的细节。

\noindent\eventanchor{SYD-1214-MONGOL-268}{詔三摸合、石抹明安與斫答等圍中都}\textbf{蒙古。}
1214年（元中统-45年；公元换算据《元史》本纪纪年） 的记载将 蒙古 与〔詔三摸合、石抹明安與斫答等圍中都〕相连。此处可确认的是这项被记录的行动；题名保留的城、路、水道或机构名称，是目前可以落地的空间线索。

前一层压力可由以下线索辨认：本纪年条记录政令或机构处置；实施背景须参照制度材料。 这一处置留下的后续条件是：可在同年及后续政令中追踪；条文不单独证明地方执行。

关于可说范围，须保留：《元史》为明初仓促官修，本纪保存大量实录条目但有编次与纪年问题；行省执行、数字、族称及对外关系须以条格、多语文书和对方史料互校。 进一步讨论可参照 蒙古征服、草原政治与跨区域文献比较研究

\subsection{1215年（元中统-44年；公元换算据《元史》本纪纪年）}
\noindent\eventanchor{SYD-1215-MONGOL-010}{八月，天倪取平州，金經略使乞住降}\textbf{蒙古。}
〔八月，天倪取平州，金經略使乞住降〕见于 《元史》卷001〈本纪第1卷本年条〉 的 1215年（元中统-44年；公元换算据《元史》本纪纪年） 记录。蒙古 的选择因此有了可回查的时间位置；题名保留的城、路、水道或机构名称，是目前可以落地的空间线索。

本纪从朝廷视角记录军政行动；出兵与战果需分辨。 记录所见的结果则是：后续路线、控制与损失应与对方记录和遗址材料复核。 日期相邻不等于因果已经证实。

证据的边界是：《元史》为明初仓促官修，本纪保存大量实录条目但有编次与纪年问题；行省执行、数字、族称及对外关系须以条格、多语文书和对方史料互校。 蒙古征服、草原政治与跨区域文献比较研究 可用于继续检验这一结论。

\noindent\eventanchor{SYD-1215-MONGOL-061}{冬十月，金宣撫蒲鮮萬奴據遼東，僭稱天王，國號大真，改元天泰}\textbf{蒙古。}
1215年（元中统-44年；公元换算据《元史》本纪纪年），《元史》卷001〈本纪第1卷本年条〉 所见的〔冬十月，金宣撫蒲鮮萬奴據遼東，僭稱天王，國號大真，改元天泰〕把 蒙古 的处置留在可复核的年序中；材料未将地点细化到城址，不能用中枢所在地替它补足。

把本项放回事件链，能够看到的前提为：本纪年条记录政令或机构处置；实施背景须参照制度材料。 而它所打开或收紧的下一步是：可在同年及后续政令中追踪；条文不单独证明地方执行。

《元史》为明初仓促官修，本纪保存大量实录条目但有编次与纪年问题；行省执行、数字、族称及对外关系须以条格、多语文书和对方史料互校。 因此，蒙古征服、草原政治与跨区域文献比较研究 只能扩展问题，不能代替未保存的细节。

\noindent\eventanchor{SYD-1215-MONGOL-110}{二月，木華黎攻北京，金元帥寅答虎、烏古倫以城降，以寅答虎為留守，吾也而權兵馬都元帥鎮之}\textbf{蒙古。}
蒙古 在 1215年（元中统-44年；公元换算据《元史》本纪纪年） 的材料痕迹是〔二月，木華黎攻北京，金元帥寅答虎、烏古倫以城降，以寅答虎為留守，吾也而權兵馬都元帥鎮之〕，出处为 《元史》卷001〈本纪第1卷本年条〉。题名保留的城、路、水道或机构名称，是目前可以落地的空间线索。

它承接的条件是：本纪从朝廷视角记录军政行动；出兵与战果需分辨。 随后可追踪的变化为：后续路线、控制与损失应与对方记录和遗址材料复核。

关于可说范围，须保留：《元史》为明初仓促官修，本纪保存大量实录条目但有编次与纪年问题；行省执行、数字、族称及对外关系须以条格、多语文书和对方史料互校。 进一步讨论可参照 蒙古征服、草原政治与跨区域文献比较研究

\noindent\eventanchor{SYD-1215-MONGOL-157}{鯨弟致遂據錦州，僭號漢興皇帝，改元興龍}\textbf{蒙古。}
从 《元史》卷001〈本纪第1卷本年条〉 的 1215年（元中统-44年；公元换算据《元史》本纪纪年） 条目看，〔鯨弟致遂據錦州，僭號漢興皇帝，改元興龍〕使 蒙古 的一个具体行动进入叙事；题名保留的城、路、水道或机构名称，是目前可以落地的空间线索。

前一层压力可由以下线索辨认：本纪年条记录政令或机构处置；实施背景须参照制度材料。 这一处置留下的后续条件是：可在同年及后续政令中追踪；条文不单独证明地方执行。

证据的边界是：《元史》为明初仓促官修，本纪保存大量实录条目但有编次与纪年问题；行省执行、数字、族称及对外关系须以条格、多语文书和对方史料互校。 蒙古征服、草原政治与跨区域文献比较研究 可用于继续检验这一结论。

\noindent\eventanchor{SYD-1215-MONGOL-199}{木華黎遣史進道等攻廣寧府，降之}\textbf{蒙古。}
1215年（元中统-44年；公元换算据《元史》本纪纪年） 的记载将 蒙古 与〔木華黎遣史進道等攻廣寧府，降之〕相连。此处可确认的是这项被记录的行动；题名保留的城、路、水道或机构名称，是目前可以落地的空间线索。

本纪从朝廷视角记录军政行动；出兵与战果需分辨。 记录所见的结果则是：后续路线、控制与损失应与对方记录和遗址材料复核。 日期相邻不等于因果已经证实。

《元史》为明初仓促官修，本纪保存大量实录条目但有编次与纪年问题；行省执行、数字、族称及对外关系须以条格、多语文书和对方史料互校。 因此，蒙古征服、草原政治与跨区域文献比较研究 只能扩展问题，不能代替未保存的细节。

\noindent\eventanchor{SYD-1215-MONGOL-237}{遣忽都忽等籍中都帑藏}\textbf{蒙古。}
〔遣忽都忽等籍中都帑藏〕见于 《元史》卷001〈本纪第1卷本年条〉 的 1215年（元中统-44年；公元换算据《元史》本纪纪年） 记录。蒙古 的选择因此有了可回查的时间位置；题名保留的城、路、水道或机构名称，是目前可以落地的空间线索。

把本项放回事件链，能够看到的前提为：本纪年条提供日期和中枢可见行动；前因不据单条补定。 而它所打开或收紧的下一步是：可回查同年及后续条目，地方影响仍须另证。

关于可说范围，须保留：《元史》为明初仓促官修，本纪保存大量实录条目但有编次与纪年问题；行省执行、数字、族称及对外关系须以条格、多语文书和对方史料互校。 进一步讨论可参照 蒙古征服、草原政治与跨区域文献比较研究

\noindent\eventanchor{SYD-1215-MONGOL-269}{遣乙職里往諭金主以河北、山東未下諸城來獻，及去帝號為河南王，當為罷兵}\textbf{蒙古与金。}
1215年（元中统-44年；公元换算据《元史》本纪纪年），《元史》卷001〈本纪第1卷本年条〉 所见的〔遣乙職里往諭金主以河北、山東未下諸城來獻，及去帝號為河南王，當為罷兵〕把 蒙古与金 的处置留在可复核的年序中；题名保留的城、路、水道或机构名称，是目前可以落地的空间线索。

它承接的条件是：本纪年条记录政令或机构处置；实施背景须参照制度材料。 随后可追踪的变化为：可在同年及后续政令中追踪；条文不单独证明地方执行。

证据的边界是：《元史》为明初仓促官修，本纪保存大量实录条目但有编次与纪年问题；行省执行、数字、族称及对外关系须以条格、多语文书和对方史料互校。 蒙古征服、草原政治与跨区域文献比较研究 可用于继续检验这一结论。

\subsection{1215年}
\noindent\eventanchor{YUAN-1215-EVENT-154}{蒙古军攻占中都}\textbf{蒙古与金。}
蒙古与金 在 1215年 的材料痕迹是〔蒙古军攻占中都〕，出处为 《元史》相关本纪；《元朝秘史》、波斯文史籍或《高丽史》对应叙事；《元典章》《通制条格》、多语碑刻与文书（据事核）。题名保留的城、路、水道或机构名称，是目前可以落地的空间线索。

前一层压力可由以下线索辨认：蒙古诸部整合及对西夏、金的战争中既有的联盟、交通与补给压力，为“蒙古军攻占中都”提供了直接的军事或动员背景。 这一处置留下的后续条件是：“蒙古军攻占中都”改变了相关城镇、道路或政权间的军事选择；伤亡、俘获和控制范围仅可按各方材料分别确认。

《元史》明初速修，纪年与人名有异同；蒙古大汗政权不等同所有汗国，征服、赋役和身份分类须按地区及材料语言分列。 因此，蒙古帝国、元朝行省财政、多语文书与区域社会研究 只能扩展问题，不能代替未保存的细节。

\subsection{1216 年（共享年序桥接）}
\noindent\eventanchor{SY-1216-CHRONOLOGY-BRIDGE}{【C级年序桥接】保留1216 年的多政权并行检索入口；本条不主张所列政权在当年发生直接互动，具体事件须另以单项材料入账。}\textbf{南宋、金、西夏与蒙古（年序桥接）。}
从 《宋史》；《金史》；《元史》；蒙古文、波斯文及地方材料（据事核） 的 1216 年（共享年序桥接） 条目看，〔【C级年序桥接】保留1216 年的多政权并行检索入口；本条不主张所列政权在当年发生直接互动，具体事件须另以单项材料入账。〕使 南宋、金、西夏与蒙古（年序桥接） 的一个具体行动进入叙事；题名保留的城、路、水道或机构名称，是目前可以落地的空间线索。

相邻已入账事件之间缺少该年度的共享年序入口，易使跨政权材料被单一政权叙事遮蔽。 记录所见的结果则是：保留该年度的检索与补账位置；后续仅在获得可核单项材料时拆分为具体政治、财政、军事、社会或文化事件。 日期相邻不等于因果已经证实。

关于可说范围，须保留：本条只确认该年位于可追踪的共享年序中；正史、地方文书和外部材料的保存密度不一，不能由桥接标签推断行动、统属、疆界或社会影响。 进一步讨论可参照 蒙古征服、区域政权转换与跨语种年序研究

\subsection{1216年（元中统-43年；公元换算据《元史》本纪纪年）}
\noindent\eventanchor{SYD-1216-MONGOL-011}{冬十月，蒲鮮萬奴降，以其子帖哥入侍}\textbf{蒙古。}
1216年（元中统-43年；公元换算据《元史》本纪纪年） 的记载将 蒙古 与〔冬十月，蒲鮮萬奴降，以其子帖哥入侍〕相连。此处可确认的是这项被记录的行动；材料未将地点细化到城址，不能用中枢所在地替它补足。

把本项放回事件链，能够看到的前提为：本纪从朝廷视角记录军政行动；出兵与战果需分辨。 而它所打开或收紧的下一步是：后续路线、控制与损失应与对方记录和遗址材料复核。

证据的边界是：《元史》为明初仓促官修，本纪保存大量实录条目但有编次与纪年问题；行省执行、数字、族称及对外关系须以条格、多语文书和对方史料互校。 蒙古征服、草原政治与跨区域文献比较研究 可用于继续检验这一结论。

\subsection{1217年（元中统-42年；公元换算据《元史》本纪纪年）}
\noindent\eventanchor{SYD-1217-MONGOL-012}{察罕破金監軍夾谷於霸州，金求和，察罕乃還}\textbf{蒙古。}
〔察罕破金監軍夾谷於霸州，金求和，察罕乃還〕见于 《元史》卷001〈本纪第1卷本年条〉 的 1217年（元中统-42年；公元换算据《元史》本纪纪年） 记录。蒙古 的选择因此有了可回查的时间位置；题名保留的城、路、水道或机构名称，是目前可以落地的空间线索。

它承接的条件是：本纪从朝廷视角记录军政行动；出兵与战果需分辨。 随后可追踪的变化为：后续路线、控制与损失应与对方记录和遗址材料复核。

《元史》为明初仓促官修，本纪保存大量实录条目但有编次与纪年问题；行省执行、数字、族称及对外关系须以条格、多语文书和对方史料互校。 因此，蒙古征服、草原政治与跨区域文献比较研究 只能扩展问题，不能代替未保存的细节。

\noindent\eventanchor{SYD-1217-MONGOL-062}{秋八月，以木華黎為太師，封國王，將蒙古、乣、漢諸軍南征，拔遂城、蠡州}\textbf{蒙古。}
1217年（元中统-42年；公元换算据《元史》本纪纪年），《元史》卷001〈本纪第1卷本年条〉 所见的〔秋八月，以木華黎為太師，封國王，將蒙古、乣、漢諸軍南征，拔遂城、蠡州〕把 蒙古 的处置留在可复核的年序中；题名保留的城、路、水道或机构名称，是目前可以落地的空间线索。

前一层压力可由以下线索辨认：本纪年条提供日期和中枢可见行动；前因不据单条补定。 这一处置留下的后续条件是：可回查同年及后续条目，地方影响仍须另证。

关于可说范围，须保留：《元史》为明初仓促官修，本纪保存大量实录条目但有编次与纪年问题；行省执行、数字、族称及对外关系须以条格、多语文书和对方史料互校。 进一步讨论可参照 蒙古征服、草原政治与跨区域文献比较研究

\noindent\eventanchor{SYD-1217-MONGOL-111}{是歲，禿滿部民叛，命鉢魯完、朵魯伯討平之}\textbf{蒙古。}
蒙古 在 1217年（元中统-42年；公元换算据《元史》本纪纪年） 的材料痕迹是〔是歲，禿滿部民叛，命鉢魯完、朵魯伯討平之〕，出处为 《元史》卷001〈本纪第1卷本年条〉。材料未将地点细化到城址，不能用中枢所在地替它补足。

本纪年条记录政令或机构处置；实施背景须参照制度材料。 记录所见的结果则是：可在同年及后续政令中追踪；条文不单独证明地方执行。 日期相邻不等于因果已经证实。

证据的边界是：《元史》为明初仓促官修，本纪保存大量实录条目但有编次与纪年问题；行省执行、数字、族称及对外关系须以条格、多语文书和对方史料互校。 蒙古征服、草原政治与跨区域文献比较研究 可用于继续检验这一结论。

\subsection{1218年（元中统-41年；公元换算据《元史》本纪纪年）}
\noindent\eventanchor{SYD-1218-MONGOL-013}{高麗王㬚遂降，請歲貢方物}\textbf{蒙古。}
从 《元史》卷001〈本纪第1卷本年条〉 的 1218年（元中统-41年；公元换算据《元史》本纪纪年） 条目看，〔高麗王㬚遂降，請歲貢方物〕使 蒙古 的一个具体行动进入叙事；材料未将地点细化到城址，不能用中枢所在地替它补足。

把本项放回事件链，能够看到的前提为：本纪从朝廷视角记录军政行动；出兵与战果需分辨。 而它所打开或收紧的下一步是：后续路线、控制与损失应与对方记录和遗址材料复核。

《元史》为明初仓促官修，本纪保存大量实录条目但有编次与纪年问题；行省执行、数字、族称及对外关系须以条格、多语文书和对方史料互校。 因此，蒙古征服、草原政治与跨区域文献比较研究 只能扩展问题，不能代替未保存的细节。

\noindent\eventanchor{SYD-1218-MONGOL-063}{金將武仙攻滿城，張柔擊敗之}\textbf{蒙古。}
1218年（元中统-41年；公元换算据《元史》本纪纪年） 的记载将 蒙古 与〔金將武仙攻滿城，張柔擊敗之〕相连。此处可确认的是这项被记录的行动；题名保留的城、路、水道或机构名称，是目前可以落地的空间线索。

它承接的条件是：本纪从朝廷视角记录军政行动；出兵与战果需分辨。 随后可追踪的变化为：后续路线、控制与损失应与对方记录和遗址材料复核。

关于可说范围，须保留：《元史》为明初仓促官修，本纪保存大量实录条目但有编次与纪年问题；行省执行、数字、族称及对外关系须以条格、多语文书和对方史料互校。 进一步讨论可参照 蒙古征服、草原政治与跨区域文献比较研究

\noindent\eventanchor{SYD-1218-MONGOL-112}{木華黎自西京入河東，克太原、平陽及忻、代、澤、潞、汾、霍等州}\textbf{蒙古。}
〔木華黎自西京入河東，克太原、平陽及忻、代、澤、潞、汾、霍等州〕见于 《元史》卷001〈本纪第1卷本年条〉 的 1218年（元中统-41年；公元换算据《元史》本纪纪年） 记录。蒙古 的选择因此有了可回查的时间位置；题名保留的城、路、水道或机构名称，是目前可以落地的空间线索。

前一层压力可由以下线索辨认：本纪保留灾害或工程的报告地点与朝廷处置；灾情范围需另证。 这一处置留下的后续条件是：可与地方文书、河工和环境材料比较，不将单点记录外推为全域因果。

证据的边界是：《元史》为明初仓促官修，本纪保存大量实录条目但有编次与纪年问题；行省执行、数字、族称及对外关系须以条格、多语文书和对方史料互校。 蒙古征服、草原政治与跨区域文献比较研究 可用于继续检验这一结论。

\noindent\eventanchor{SYD-1218-MONGOL-158}{契丹六哥據高麗江東城，命哈真、札剌率師平之}\textbf{蒙古。}
1218年（元中统-41年；公元换算据《元史》本纪纪年），《元史》卷001〈本纪第1卷本年条〉 所见的〔契丹六哥據高麗江東城，命哈真、札剌率師平之〕把 蒙古 的处置留在可复核的年序中；题名保留的城、路、水道或机构名称，是目前可以落地的空间线索。

本纪年条记录政令或机构处置；实施背景须参照制度材料。 记录所见的结果则是：可在同年及后续政令中追踪；条文不单独证明地方执行。 日期相邻不等于因果已经证实。

《元史》为明初仓促官修，本纪保存大量实录条目但有编次与纪年问题；行省执行、数字、族称及对外关系须以条格、多语文书和对方史料互校。 因此，蒙古征服、草原政治与跨区域文献比较研究 只能扩展问题，不能代替未保存的细节。

\noindent\eventanchor{SYD-1218-MONGOL-200}{十三年戊寅秋八月，兵出紫荊口，獲金行元帥事張柔，命還其舊職}\textbf{蒙古。}
蒙古 在 1218年（元中统-41年；公元换算据《元史》本纪纪年） 的材料痕迹是〔十三年戊寅秋八月，兵出紫荊口，獲金行元帥事張柔，命還其舊職〕，出处为 《元史》卷001〈本纪第1卷本年条〉。材料未将地点细化到城址，不能用中枢所在地替它补足。

把本项放回事件链，能够看到的前提为：本纪年条记录政令或机构处置；实施背景须参照制度材料。 而它所打开或收紧的下一步是：可在同年及后续政令中追踪；条文不单独证明地方执行。

关于可说范围，须保留：《元史》为明初仓促官修，本纪保存大量实录条目但有编次与纪年问题；行省执行、数字、族称及对外关系须以条格、多语文书和对方史料互校。 进一步讨论可参照 蒙古征服、草原政治与跨区域文献比较研究

\subsection{1218年}
\noindent\eventanchor{YUAN-1218-EVENT-155}{蒙古消灭西辽并进入中亚政治网络}\textbf{蒙古与西辽。}
从 《元史》相关本纪；《元朝秘史》、波斯文史籍或《高丽史》对应叙事；《元典章》《通制条格》、多语碑刻与文书（据事核） 的 1218年 条目看，〔蒙古消灭西辽并进入中亚政治网络〕使 蒙古与西辽 的一个具体行动进入叙事；材料未将地点细化到城址，不能用中枢所在地替它补足。

它承接的条件是：“蒙古消灭西辽并进入中亚政治网络”发生在蒙古诸部整合及对西夏、金的战争的具体压力与机会之中，相关行动者的选择不能脱离此前事件链解释。 随后可追踪的变化为：“蒙古消灭西辽并进入中亚政治网络”为后续政治、边防或资源配置留下条件；其社会影响与实际覆盖范围仍须以异类材料限定。

证据的边界是：《元史》明初速修，纪年与人名有异同；蒙古大汗政权不等同所有汗国，征服、赋役和身份分类须按地区及材料语言分列。 蒙古帝国、元朝行省财政、多语文书与区域社会研究 可用于继续检验这一结论。

\subsection{1219年（元中统-40年；公元换算据《元史》本纪纪年）}
\noindent\eventanchor{SYD-1219-MONGOL-014}{秋，木華黎克岢、嵐、吉、隰等州，進攻絳州，拔其城，屠之}\textbf{蒙古。}
1219年（元中统-40年；公元换算据《元史》本纪纪年） 的记载将 蒙古 与〔秋，木華黎克岢、嵐、吉、隰等州，進攻絳州，拔其城，屠之〕相连。此处可确认的是这项被记录的行动；题名保留的城、路、水道或机构名称，是目前可以落地的空间线索。

前一层压力可由以下线索辨认：本纪从朝廷视角记录军政行动；出兵与战果需分辨。 这一处置留下的后续条件是：后续路线、控制与损失应与对方记录和遗址材料复核。

《元史》为明初仓促官修，本纪保存大量实录条目但有编次与纪年问题；行省执行、数字、族称及对外关系须以条格、多语文书和对方史料互校。 因此，蒙古征服、草原政治与跨区域文献比较研究 只能扩展问题，不能代替未保存的细节。

\noindent\eventanchor{SYD-1219-MONGOL-064}{十四年己卯春，張柔敗武仙，降祁陽、曲陽、中山等城}\textbf{蒙古。}
〔十四年己卯春，張柔敗武仙，降祁陽、曲陽、中山等城〕见于 《元史》卷001〈本纪第1卷本年条〉 的 1219年（元中统-40年；公元换算据《元史》本纪纪年） 记录。蒙古 的选择因此有了可回查的时间位置；题名保留的城、路、水道或机构名称，是目前可以落地的空间线索。

本纪从朝廷视角记录军政行动；出兵与战果需分辨。 记录所见的结果则是：后续路线、控制与损失应与对方记录和遗址材料复核。 日期相邻不等于因果已经证实。

关于可说范围，须保留：《元史》为明初仓促官修，本纪保存大量实录条目但有编次与纪年问题；行省执行、数字、族称及对外关系须以条格、多语文书和对方史料互校。 进一步讨论可参照 蒙古征服、草原政治与跨区域文献比较研究

\noindent\eventanchor{SYD-1219-MONGOL-113}{夏六月，西域殺使者，帝率師親征，取訛答剌城，擒其酋哈只兒只蘭禿}\textbf{蒙古。}
1219年（元中统-40年；公元换算据《元史》本纪纪年），《元史》卷001〈本纪第1卷本年条〉 所见的〔夏六月，西域殺使者，帝率師親征，取訛答剌城，擒其酋哈只兒只蘭禿〕把 蒙古 的处置留在可复核的年序中；题名保留的城、路、水道或机构名称，是目前可以落地的空间线索。

把本项放回事件链，能够看到的前提为：本纪从朝廷视角记录军政行动；出兵与战果需分辨。 而它所打开或收紧的下一步是：后续路线、控制与损失应与对方记录和遗址材料复核。

证据的边界是：《元史》为明初仓促官修，本纪保存大量实录条目但有编次与纪年问题；行省执行、数字、族称及对外关系须以条格、多语文书和对方史料互校。 蒙古征服、草原政治与跨区域文献比较研究 可用于继续检验这一结论。

\subsection{1219年}
\noindent\eventanchor{YUAN-1219-EVENT-156}{蒙古发动对花剌子模的大规模战争}\textbf{蒙古与花剌子模。}
蒙古与花剌子模 在 1219年 的材料痕迹是〔蒙古发动对花剌子模的大规模战争〕，出处为 《元史》相关本纪；《元朝秘史》、波斯文史籍或《高丽史》对应叙事；《元典章》《通制条格》、多语碑刻与文书（据事核）。材料未将地点细化到城址，不能用中枢所在地替它补足。

它承接的条件是：蒙古诸部整合及对西夏、金的战争中既有的联盟、交通与补给压力，为“蒙古发动对花剌子模的大规模战争”提供了直接的军事或动员背景。 随后可追踪的变化为：“蒙古发动对花剌子模的大规模战争”改变了相关城镇、道路或政权间的军事选择；伤亡、俘获和控制范围仅可按各方材料分别确认。

《元史》明初速修，纪年与人名有异同；蒙古大汗政权不等同所有汗国，征服、赋役和身份分类须按地区及材料语言分列。 因此，蒙古帝国、元朝行省财政、多语文书与区域社会研究 只能扩展问题，不能代替未保存的细节。

\subsection{1220 年（共享年序桥接）}
\noindent\eventanchor{SY-1220-CHRONOLOGY-BRIDGE}{【C级年序桥接】保留1220 年的多政权并行检索入口；本条不主张所列政权在当年发生直接互动，具体事件须另以单项材料入账。}\textbf{南宋、金、西夏与蒙古（年序桥接）。}
从 《宋史》；《金史》；《元史》；蒙古文、波斯文及地方材料（据事核） 的 1220 年（共享年序桥接） 条目看，〔【C级年序桥接】保留1220 年的多政权并行检索入口；本条不主张所列政权在当年发生直接互动，具体事件须另以单项材料入账。〕使 南宋、金、西夏与蒙古（年序桥接） 的一个具体行动进入叙事；题名保留的城、路、水道或机构名称，是目前可以落地的空间线索。

前一层压力可由以下线索辨认：相邻已入账事件之间缺少该年度的共享年序入口，易使跨政权材料被单一政权叙事遮蔽。 这一处置留下的后续条件是：保留该年度的检索与补账位置；后续仅在获得可核单项材料时拆分为具体政治、财政、军事、社会或文化事件。

关于可说范围，须保留：本条只确认该年位于可追踪的共享年序中；正史、地方文书和外部材料的保存密度不一，不能由桥接标签推断行动、统属、疆界或社会影响。 进一步讨论可参照 蒙古征服、区域政权转换与跨语种年序研究

\subsection{1220年（元中统-39年；公元换算据《元史》本纪纪年）}
\noindent\eventanchor{SYD-1220-MONGOL-015}{冬，金邢州節度使武貴降}\textbf{蒙古。}
1220年（元中统-39年；公元换算据《元史》本纪纪年） 的记载将 蒙古 与〔冬，金邢州節度使武貴降〕相连。此处可确认的是这项被记录的行动；题名保留的城、路、水道或机构名称，是目前可以落地的空间线索。

本纪从朝廷视角记录军政行动；出兵与战果需分辨。 记录所见的结果则是：后续路线、控制与损失应与对方记录和遗址材料复核。 日期相邻不等于因果已经证实。

证据的边界是：《元史》为明初仓促官修，本纪保存大量实录条目但有编次与纪年问题；行省执行、数字、族称及对外关系须以条格、多语文书和对方史料互校。 蒙古征服、草原政治与跨区域文献比较研究 可用于继续检验这一结论。

\noindent\eventanchor{SYD-1220-MONGOL-065}{東平嚴實籍彰德、大名、磁、洺、恩、博、滑、濬等州戶三十萬來歸，木華黎承制授實金紫光祿大夫、行尚書省事}\textbf{蒙古。}
〔東平嚴實籍彰德、大名、磁、洺、恩、博、滑、濬等州戶三十萬來歸，木華黎承制授實金紫光祿大夫、行尚書省事〕见于 《元史》卷001〈本纪第1卷本年条〉 的 1220年（元中统-39年；公元换算据《元史》本纪纪年） 记录。蒙古 的选择因此有了可回查的时间位置；题名保留的城、路、水道或机构名称，是目前可以落地的空间线索。

把本项放回事件链，能够看到的前提为：本纪年条记录政令或机构处置；实施背景须参照制度材料。 而它所打开或收紧的下一步是：可在同年及后续政令中追踪；条文不单独证明地方执行。

《元史》为明初仓促官修，本纪保存大量实录条目但有编次与纪年问题；行省执行、数字、族称及对外关系须以条格、多语文书和对方史料互校。 因此，蒙古征服、草原政治与跨区域文献比较研究 只能扩展问题，不能代替未保存的细节。

\noindent\eventanchor{SYD-1220-MONGOL-114}{木華黎攻東平不克，留嚴實守之，撤圍趨洺州，分兵徇河北諸郡}\textbf{蒙古。}
1220年（元中统-39年；公元换算据《元史》本纪纪年），《元史》卷001〈本纪第1卷本年条〉 所见的〔木華黎攻東平不克，留嚴實守之，撤圍趨洺州，分兵徇河北諸郡〕把 蒙古 的处置留在可复核的年序中；题名保留的城、路、水道或机构名称，是目前可以落地的空间线索。

它承接的条件是：本纪保留灾害或工程的报告地点与朝廷处置；灾情范围需另证。 随后可追踪的变化为：可与地方文书、河工和环境材料比较，不将单点记录外推为全域因果。

关于可说范围，须保留：《元史》为明初仓促官修，本纪保存大量实录条目但有编次与纪年问题；行省执行、数字、族称及对外关系须以条格、多语文书和对方史料互校。 进一步讨论可参照 蒙古征服、草原政治与跨区域文献比较研究

\noindent\eventanchor{SYD-1220-MONGOL-159}{木華黎徇地至真定，武仙出降}\textbf{蒙古。}
蒙古 在 1220年（元中统-39年；公元换算据《元史》本纪纪年） 的材料痕迹是〔木華黎徇地至真定，武仙出降〕，出处为 《元史》卷001〈本纪第1卷本年条〉。材料未将地点细化到城址，不能用中枢所在地替它补足。

前一层压力可由以下线索辨认：本纪从朝廷视角记录军政行动；出兵与战果需分辨。 这一处置留下的后续条件是：后续路线、控制与损失应与对方记录和遗址材料复核。

证据的边界是：《元史》为明初仓促官修，本纪保存大量实录条目但有编次与纪年问题；行省执行、数字、族称及对外关系须以条格、多语文书和对方史料互校。 蒙古征服、草原政治与跨区域文献比较研究 可用于继续检验这一结论。

\noindent\eventanchor{SYD-1220-MONGOL-201}{秋，攻斡脫羅兒城，克之}\textbf{蒙古。}
从 《元史》卷001〈本纪第1卷本年条〉 的 1220年（元中统-39年；公元换算据《元史》本纪纪年） 条目看，〔秋，攻斡脫羅兒城，克之〕使 蒙古 的一个具体行动进入叙事；题名保留的城、路、水道或机构名称，是目前可以落地的空间线索。

本纪从朝廷视角记录军政行动；出兵与战果需分辨。 记录所见的结果则是：后续路线、控制与损失应与对方记录和遗址材料复核。 日期相邻不等于因果已经证实。

《元史》为明初仓促官修，本纪保存大量实录条目但有编次与纪年问题；行省执行、数字、族称及对外关系须以条格、多语文书和对方史料互校。 因此，蒙古征服、草原政治与跨区域文献比较研究 只能扩展问题，不能代替未保存的细节。

\subsection{1221 年（共享年序桥接）}
\noindent\eventanchor{SY-1221-CHRONOLOGY-BRIDGE}{【C级年序桥接】保留1221 年的多政权并行检索入口；本条不主张所列政权在当年发生直接互动，具体事件须另以单项材料入账。}\textbf{南宋、金、西夏与蒙古（年序桥接）。}
1221 年（共享年序桥接） 的记载将 南宋、金、西夏与蒙古（年序桥接） 与〔【C级年序桥接】保留1221 年的多政权并行检索入口；本条不主张所列政权在当年发生直接互动，具体事件须另以单项材料入账。〕相连。此处可确认的是这项被记录的行动；题名保留的城、路、水道或机构名称，是目前可以落地的空间线索。

把本项放回事件链，能够看到的前提为：相邻已入账事件之间缺少该年度的共享年序入口，易使跨政权材料被单一政权叙事遮蔽。 而它所打开或收紧的下一步是：保留该年度的检索与补账位置；后续仅在获得可核单项材料时拆分为具体政治、财政、军事、社会或文化事件。

关于可说范围，须保留：本条只确认该年位于可追踪的共享年序中；正史、地方文书和外部材料的保存密度不一，不能由桥接标签推断行动、统属、疆界或社会影响。 进一步讨论可参照 蒙古征服、区域政权转换与跨语种年序研究

\subsection{1221年（元中统-38年；公元换算据《元史》本纪纪年）}
\noindent\eventanchor{SYD-1221-MONGOL-016}{金東平行省事忙古棄城遁，嚴實入守之}\textbf{蒙古。}
〔金東平行省事忙古棄城遁，嚴實入守之〕见于 《元史》卷001〈本纪第1卷本年条〉 的 1221年（元中统-38年；公元换算据《元史》本纪纪年） 记录。蒙古 的选择因此有了可回查的时间位置；题名保留的城、路、水道或机构名称，是目前可以落地的空间线索。

它承接的条件是：本纪年条记录政令或机构处置；实施背景须参照制度材料。 随后可追踪的变化为：可在同年及后续政令中追踪；条文不单独证明地方执行。

证据的边界是：《元史》为明初仓促官修，本纪保存大量实录条目但有编次与纪年问题；行省执行、数字、族称及对外关系须以条格、多语文书和对方史料互校。 蒙古征服、草原政治与跨区域文献比较研究 可用于继续检验这一结论。

\noindent\eventanchor{SYD-1221-MONGOL-066}{木華黎出河西，克葭、綏德、保安、鄜、坊、丹等州，進攻延安，不下}\textbf{蒙古。}
1221年（元中统-38年；公元换算据《元史》本纪纪年），《元史》卷001〈本纪第1卷本年条〉 所见的〔木華黎出河西，克葭、綏德、保安、鄜、坊、丹等州，進攻延安，不下〕把 蒙古 的处置留在可复核的年序中；题名保留的城、路、水道或机构名称，是目前可以落地的空间线索。

前一层压力可由以下线索辨认：本纪保留灾害或工程的报告地点与朝廷处置；灾情范围需另证。 这一处置留下的后续条件是：可与地方文书、河工和环境材料比较，不将单点记录外推为全域因果。

《元史》为明初仓促官修，本纪保存大量实录条目但有编次与纪年问题；行省执行、数字、族称及对外关系须以条格、多语文书和对方史料互校。 因此，蒙古征服、草原政治与跨区域文献比较研究 只能扩展问题，不能代替未保存的细节。

\noindent\eventanchor{SYD-1221-MONGOL-115}{秋，帝攻班勒紇等城，皇子朮赤、察合台、窩闊台分攻玉龍傑赤等城，下之}\textbf{蒙古。}
蒙古 在 1221年（元中统-38年；公元换算据《元史》本纪纪年） 的材料痕迹是〔秋，帝攻班勒紇等城，皇子朮赤、察合台、窩闊台分攻玉龍傑赤等城，下之〕，出处为 《元史》卷001〈本纪第1卷本年条〉。题名保留的城、路、水道或机构名称，是目前可以落地的空间线索。

本纪年条记录政令或机构处置；实施背景须参照制度材料。 记录所见的结果则是：可在同年及后续政令中追踪；条文不单独证明地方执行。 日期相邻不等于因果已经证实。

关于可说范围，须保留：《元史》为明初仓促官修，本纪保存大量实录条目但有编次与纪年问题；行省执行、数字、族称及对外关系须以条格、多语文书和对方史料互校。 进一步讨论可参照 蒙古征服、草原政治与跨区域文献比较研究

\noindent\eventanchor{SYD-1221-MONGOL-160}{十六年辛巳春，帝攻卜哈兒、薛迷思干等城，皇子朮赤攻養吉干、八兒真等城，並下之}\textbf{蒙古。}
从 《元史》卷001〈本纪第1卷本年条〉 的 1221年（元中统-38年；公元换算据《元史》本纪纪年） 条目看，〔十六年辛巳春，帝攻卜哈兒、薛迷思干等城，皇子朮赤攻養吉干、八兒真等城，並下之〕使 蒙古 的一个具体行动进入叙事；题名保留的城、路、水道或机构名称，是目前可以落地的空间线索。

把本项放回事件链，能够看到的前提为：本纪从朝廷视角记录军政行动；出兵与战果需分辨。 而它所打开或收紧的下一步是：后续路线、控制与损失应与对方记录和遗址材料复核。

证据的边界是：《元史》为明初仓促官修，本纪保存大量实录条目但有编次与纪年问题；行省执行、数字、族称及对外关系须以条格、多语文书和对方史料互校。 蒙古征服、草原政治与跨区域文献比较研究 可用于继续检验这一结论。

\noindent\eventanchor{SYD-1221-MONGOL-202}{十一月，宋京東安撫使張琳以京東諸郡來降，以琳為滄、景、濱、棣等州行都元帥}\textbf{蒙古与南宋。}
1221年（元中统-38年；公元换算据《元史》本纪纪年） 的记载将 蒙古与南宋 与〔十一月，宋京東安撫使張琳以京東諸郡來降，以琳為滄、景、濱、棣等州行都元帥〕相连。此处可确认的是这项被记录的行动；题名保留的城、路、水道或机构名称，是目前可以落地的空间线索。

它承接的条件是：本纪从朝廷视角记录军政行动；出兵与战果需分辨。 随后可追踪的变化为：后续路线、控制与损失应与对方记录和遗址材料复核。

《元史》为明初仓促官修，本纪保存大量实录条目但有编次与纪年问题；行省执行、数字、族称及对外关系须以条格、多语文书和对方史料互校。 因此，蒙古征服、草原政治与跨区域文献比较研究 只能扩展问题，不能代替未保存的细节。

\noindent\eventanchor{SYD-1221-MONGOL-238}{夏六月，宋〔漣〕水忠義統轄石珪率眾來降，以珪為濟、兖、單三州總管}\textbf{蒙古与南宋。}
〔夏六月，宋〔漣〕水忠義統轄石珪率眾來降，以珪為濟、兖、單三州總管〕见于 《元史》卷001〈本纪第1卷本年条〉 的 1221年（元中统-38年；公元换算据《元史》本纪纪年） 记录。蒙古与南宋 的选择因此有了可回查的时间位置；题名保留的城、路、水道或机构名称，是目前可以落地的空间线索。

前一层压力可由以下线索辨认：本纪保留灾害或工程的报告地点与朝廷处置；灾情范围需另证。 这一处置留下的后续条件是：可与地方文书、河工和环境材料比较，不将单点记录外推为全域因果。

关于可说范围，须保留：《元史》为明初仓促官修，本纪保存大量实录条目但有编次与纪年问题；行省执行、数字、族称及对外关系须以条格、多语文书和对方史料互校。 进一步讨论可参照 蒙古征服、草原政治与跨区域文献比较研究

\noindent\eventanchor{SYD-1221-MONGOL-270}{夏四月，駐蹕鐵門關，金主遣烏古孫仲端奉國書請和，稱帝為兄}\textbf{蒙古与金。}
1221年（元中统-38年；公元换算据《元史》本纪纪年），《元史》卷001〈本纪第1卷本年条〉 所见的〔夏四月，駐蹕鐵門關，金主遣烏古孫仲端奉國書請和，稱帝為兄〕把 蒙古与金 的处置留在可复核的年序中；题名保留的城、路、水道或机构名称，是目前可以落地的空间线索。

本纪记录使节、贡赐或往来，不等于稳定统属关系。 记录所见的结果则是：可与对方编年和交通、文书材料核对其后续落实。 日期相邻不等于因果已经证实。

证据的边界是：《元史》为明初仓促官修，本纪保存大量实录条目但有编次与纪年问题；行省执行、数字、族称及对外关系须以条格、多语文书和对方史料互校。 蒙古征服、草原政治与跨区域文献比较研究 可用于继续检验这一结论。

\subsection{1222 年（共享年序桥接）}
\noindent\eventanchor{SY-1222-CHRONOLOGY-BRIDGE}{【C级年序桥接】保留1222 年的多政权并行检索入口；本条不主张所列政权在当年发生直接互动，具体事件须另以单项材料入账。}\textbf{南宋、金、西夏与蒙古（年序桥接）。}
南宋、金、西夏与蒙古（年序桥接） 在 1222 年（共享年序桥接） 的材料痕迹是〔【C级年序桥接】保留1222 年的多政权并行检索入口；本条不主张所列政权在当年发生直接互动，具体事件须另以单项材料入账。〕，出处为 《宋史》；《金史》；《元史》；蒙古文、波斯文及地方材料（据事核）。题名保留的城、路、水道或机构名称，是目前可以落地的空间线索。

把本项放回事件链，能够看到的前提为：相邻已入账事件之间缺少该年度的共享年序入口，易使跨政权材料被单一政权叙事遮蔽。 而它所打开或收紧的下一步是：保留该年度的检索与补账位置；后续仅在获得可核单项材料时拆分为具体政治、财政、军事、社会或文化事件。

本条只确认该年位于可追踪的共享年序中；正史、地方文书和外部材料的保存密度不一，不能由桥接标签推断行动、统属、疆界或社会影响。 因此，蒙古征服、区域政权转换与跨语种年序研究 只能扩展问题，不能代替未保存的细节。

\subsection{1222年（元中统-37年；公元换算据《元史》本纪纪年）}
\noindent\eventanchor{SYD-1222-MONGOL-017}{帝謂曰：「我向欲汝主授我河朔地，令汝主為河南王，彼此罷兵，汝主不從}\textbf{蒙古。}
从 《元史》卷001〈本纪第1卷本年条〉 的 1222年（元中统-37年；公元换算据《元史》本纪纪年） 条目看，〔帝謂曰：「我向欲汝主授我河朔地，令汝主為河南王，彼此罷兵，汝主不從〕使 蒙古 的一个具体行动进入叙事；题名保留的城、路、水道或机构名称，是目前可以落地的空间线索。

它承接的条件是：本纪年条记录政令或机构处置；实施背景须参照制度材料。 随后可追踪的变化为：可在同年及后续政令中追踪；条文不单独证明地方执行。

关于可说范围，须保留：《元史》为明初仓促官修，本纪保存大量实录条目但有编次与纪年问题；行省执行、数字、族称及对外关系须以条格、多语文书和对方史料互校。 进一步讨论可参照 蒙古征服、草原政治与跨区域文献比较研究

\noindent\eventanchor{SYD-1222-MONGOL-067}{帝自將擊之，擒滅里可汗}\textbf{蒙古。}
1222年（元中统-37年；公元换算据《元史》本纪纪年） 的记载将 蒙古 与〔帝自將擊之，擒滅里可汗〕相连。此处可确认的是这项被记录的行动；材料未将地点细化到城址，不能用中枢所在地替它补足。

前一层压力可由以下线索辨认：本纪从朝廷视角记录军政行动；出兵与战果需分辨。 这一处置留下的后续条件是：后续路线、控制与损失应与对方记录和遗址材料复核。

证据的边界是：《元史》为明初仓促官修，本纪保存大量实录条目但有编次与纪年问题；行省执行、数字、族称及对外关系须以条格、多语文书和对方史料互校。 蒙古征服、草原政治与跨区域文献比较研究 可用于继续检验这一结论。

\noindent\eventanchor{SYD-1222-MONGOL-116}{冬十月，金河中府來附，以石天應為兵馬都元帥守之}\textbf{蒙古。}
〔冬十月，金河中府來附，以石天應為兵馬都元帥守之〕见于 《元史》卷001〈本纪第1卷本年条〉 的 1222年（元中统-37年；公元换算据《元史》本纪纪年） 记录。蒙古 的选择因此有了可回查的时间位置；题名保留的城、路、水道或机构名称，是目前可以落地的空间线索。

本纪保留灾害或工程的报告地点与朝廷处置；灾情范围需另证。 记录所见的结果则是：可与地方文书、河工和环境材料比较，不将单点记录外推为全域因果。 日期相邻不等于因果已经证实。

《元史》为明初仓促官修，本纪保存大量实录条目但有编次与纪年问题；行省执行、数字、族称及对外关系须以条格、多语文书和对方史料互校。 因此，蒙古征服、草原政治与跨区域文献比较研究 只能扩展问题，不能代替未保存的细节。

\noindent\eventanchor{SYD-1222-MONGOL-161}{今木華黎已盡取之，乃始來請耶？」仲端乞哀，帝曰：「念汝遠來，河朔既為我有，關西數城未下者，其割付我}\textbf{蒙古。}
1222年（元中统-37年；公元换算据《元史》本纪纪年），《元史》卷001〈本纪第1卷本年条〉 所见的〔今木華黎已盡取之，乃始來請耶？」仲端乞哀，帝曰：「念汝遠來，河朔既為我有，關西數城未下者，其割付我〕把 蒙古 的处置留在可复核的年序中；题名保留的城、路、水道或机构名称，是目前可以落地的空间线索。

把本项放回事件链，能够看到的前提为：本纪保留灾害或工程的报告地点与朝廷处置；灾情范围需另证。 而它所打开或收紧的下一步是：可与地方文书、河工和环境材料比较，不将单点记录外推为全域因果。

关于可说范围，须保留：《元史》为明初仓促官修，本纪保存大量实录条目但有编次与纪年问题；行省执行、数字、族称及对外关系须以条格、多语文书和对方史料互校。 进一步讨论可参照 蒙古征服、草原政治与跨区域文献比较研究

\noindent\eventanchor{SYD-1222-MONGOL-203}{金平陽公胡天〔作〕以青龍堡降}\textbf{蒙古。}
蒙古 在 1222年（元中统-37年；公元换算据《元史》本纪纪年） 的材料痕迹是〔金平陽公胡天〔作〕以青龍堡降〕，出处为 《元史》卷001〈本纪第1卷本年条〉。材料未将地点细化到城址，不能用中枢所在地替它补足。

它承接的条件是：本纪从朝廷视角记录军政行动；出兵与战果需分辨。 随后可追踪的变化为：后续路线、控制与损失应与对方记录和遗址材料复核。

证据的边界是：《元史》为明初仓促官修，本纪保存大量实录条目但有编次与纪年问题；行省执行、数字、族称及对外关系须以条格、多语文书和对方史料互校。 蒙古征服、草原政治与跨区域文献比较研究 可用于继续检验这一结论。

\noindent\eventanchor{SYD-1222-MONGOL-239}{令汝主為河南王，勿復違也}\textbf{蒙古。}
从 《元史》卷001〈本纪第1卷本年条〉 的 1222年（元中统-37年；公元换算据《元史》本纪纪年） 条目看，〔令汝主為河南王，勿復違也〕使 蒙古 的一个具体行动进入叙事；题名保留的城、路、水道或机构名称，是目前可以落地的空间线索。

前一层压力可由以下线索辨认：本纪年条记录政令或机构处置；实施背景须参照制度材料。 这一处置留下的后续条件是：可在同年及后续政令中追踪；条文不单独证明地方执行。

《元史》为明初仓促官修，本纪保存大量实录条目但有编次与纪年问题；行省执行、数字、族称及对外关系须以条格、多语文书和对方史料互校。 因此，蒙古征服、草原政治与跨区域文献比较研究 只能扩展问题，不能代替未保存的细节。

\noindent\eventanchor{SYD-1222-MONGOL-271}{木華黎軍克乾、涇、邠、原等州，攻鳳翔不下}\textbf{蒙古。}
1222年（元中统-37年；公元换算据《元史》本纪纪年） 的记载将 蒙古 与〔木華黎軍克乾、涇、邠、原等州，攻鳳翔不下〕相连。此处可确认的是这项被记录的行动；题名保留的城、路、水道或机构名称，是目前可以落地的空间线索。

本纪从朝廷视角记录军政行动；出兵与战果需分辨。 记录所见的结果则是：后续路线、控制与损失应与对方记录和遗址材料复核。 日期相邻不等于因果已经证实。

关于可说范围，须保留：《元史》为明初仓促官修，本纪保存大量实录条目但有编次与纪年问题；行省执行、数字、族称及对外关系须以条格、多语文书和对方史料互校。 进一步讨论可参照 蒙古征服、草原政治与跨区域文献比较研究

\subsection{1223年（元中统-36年；公元换算据《元史》本纪纪年）}
\noindent\eventanchor{SYD-1223-MONGOL-018}{冬十月，金主珣殂，子守緒立}\textbf{蒙古与金。}
〔冬十月，金主珣殂，子守緒立〕见于 《元史》卷001〈本纪第1卷本年条〉 的 1223年（元中统-36年；公元换算据《元史》本纪纪年） 记录。蒙古与金 的选择因此有了可回查的时间位置；材料未将地点细化到城址，不能用中枢所在地替它补足。

把本项放回事件链，能够看到的前提为：本纪年条提供日期和中枢可见行动；前因不据单条补定。 而它所打开或收紧的下一步是：可回查同年及后续条目，地方影响仍须另证。

证据的边界是：《元史》为明初仓促官修，本纪保存大量实录条目但有编次与纪年问题；行省执行、数字、族称及对外关系须以条格、多语文书和对方史料互校。 蒙古征服、草原政治与跨区域文献比较研究 可用于继续检验这一结论。

\noindent\eventanchor{SYD-1223-MONGOL-068}{皇子朮赤、察合台、窩闊台及八剌之兵來會，遂定西域諸城，置達魯花赤監治之}\textbf{蒙古。}
1223年（元中统-36年；公元换算据《元史》本纪纪年），《元史》卷001〈本纪第1卷本年条〉 所见的〔皇子朮赤、察合台、窩闊台及八剌之兵來會，遂定西域諸城，置達魯花赤監治之〕把 蒙古 的处置留在可复核的年序中；题名保留的城、路、水道或机构名称，是目前可以落地的空间线索。

它承接的条件是：本纪年条记录政令或机构处置；实施背景须参照制度材料。 随后可追踪的变化为：可在同年及后续政令中追踪；条文不单独证明地方执行。

《元史》为明初仓促官修，本纪保存大量实录条目但有编次与纪年问题；行省执行、数字、族称及对外关系须以条格、多语文书和对方史料互校。 因此，蒙古征服、草原政治与跨区域文献比较研究 只能扩展问题，不能代替未保存的细节。

\noindent\eventanchor{SYD-1223-MONGOL-117}{是歲，宋復遣苟夢玉來}\textbf{蒙古与南宋。}
蒙古与南宋 在 1223年（元中统-36年；公元换算据《元史》本纪纪年） 的材料痕迹是〔是歲，宋復遣苟夢玉來〕，出处为 《元史》卷001〈本纪第1卷本年条〉。材料未将地点细化到城址，不能用中枢所在地替它补足。

前一层压力可由以下线索辨认：本纪年条记录政令或机构处置；实施背景须参照制度材料。 这一处置留下的后续条件是：可在同年及后续政令中追踪；条文不单独证明地方执行。

关于可说范围，须保留：《元史》为明初仓促官修，本纪保存大量实录条目但有编次与纪年问题；行省执行、数字、族称及对外关系须以条格、多语文书和对方史料互校。 进一步讨论可参照 蒙古征服、草原政治与跨区域文献比较研究

\subsection{1224年（元中统-35年；公元换算据《元史》本纪纪年）}
\noindent\eventanchor{SYD-1224-MONGOL-019}{史天倪與戰於恩州，敗之}\textbf{蒙古。}
从 《元史》卷001〈本纪第1卷本年条〉 的 1224年（元中统-35年；公元换算据《元史》本纪纪年） 条目看，〔史天倪與戰於恩州，敗之〕使 蒙古 的一个具体行动进入叙事；题名保留的城、路、水道或机构名称，是目前可以落地的空间线索。

本纪从朝廷视角记录军政行动；出兵与战果需分辨。 记录所见的结果则是：后续路线、控制与损失应与对方记录和遗址材料复核。 日期相邻不等于因果已经证实。

证据的边界是：《元史》为明初仓促官修，本纪保存大量实录条目但有编次与纪年问题；行省执行、数字、族称及对外关系须以条格、多语文书和对方史料互校。 蒙古征服、草原政治与跨区域文献比较研究 可用于继续检验这一结论。

\noindent\eventanchor{SYD-1224-MONGOL-069}{是歲，帝至東印度國，角端見，班師}\textbf{蒙古。}
1224年（元中统-35年；公元换算据《元史》本纪纪年） 的记载将 蒙古 与〔是歲，帝至東印度國，角端見，班師〕相连。此处可确认的是这项被记录的行动；材料未将地点细化到城址，不能用中枢所在地替它补足。

把本项放回事件链，能够看到的前提为：本纪年条记录政令或机构处置；实施背景须参照制度材料。 而它所打开或收紧的下一步是：可在同年及后续政令中追踪；条文不单独证明地方执行。

《元史》为明初仓促官修，本纪保存大量实录条目但有编次与纪年问题；行省执行、数字、族称及对外关系须以条格、多语文书和对方史料互校。 因此，蒙古征服、草原政治与跨区域文献比较研究 只能扩展问题，不能代替未保存的细节。

\subsection{1225 年（共享年序桥接）}
\noindent\eventanchor{SY-1225-CHRONOLOGY-BRIDGE}{【C级年序桥接】保留1225 年的多政权并行检索入口；本条不主张所列政权在当年发生直接互动，具体事件须另以单项材料入账。}\textbf{南宋、金、西夏与蒙古（年序桥接）。}
〔【C级年序桥接】保留1225 年的多政权并行检索入口；本条不主张所列政权在当年发生直接互动，具体事件须另以单项材料入账。〕见于 《宋史》；《金史》；《元史》；蒙古文、波斯文及地方材料（据事核） 的 1225 年（共享年序桥接） 记录。南宋、金、西夏与蒙古（年序桥接） 的选择因此有了可回查的时间位置；题名保留的城、路、水道或机构名称，是目前可以落地的空间线索。

它承接的条件是：相邻已入账事件之间缺少该年度的共享年序入口，易使跨政权材料被单一政权叙事遮蔽。 随后可追踪的变化为：保留该年度的检索与补账位置；后续仅在获得可核单项材料时拆分为具体政治、财政、军事、社会或文化事件。

关于可说范围，须保留：本条只确认该年位于可追踪的共享年序中；正史、地方文书和外部材料的保存密度不一，不能由桥接标签推断行动、统属、疆界或社会影响。 进一步讨论可参照 蒙古征服、区域政权转换与跨语种年序研究

\subsection{1225年（元中统-34年；公元换算据《元史》本纪纪年）}
\noindent\eventanchor{SYD-1225-MONGOL-020}{董俊判官李全亦以中山叛}\textbf{蒙古。}
1225年（元中统-34年；公元换算据《元史》本纪纪年），《元史》卷001〈本纪第1卷本年条〉 所见的〔董俊判官李全亦以中山叛〕把 蒙古 的处置留在可复核的年序中；题名保留的城、路、水道或机构名称，是目前可以落地的空间线索。

前一层压力可由以下线索辨认：本纪从朝廷视角记录军政行动；出兵与战果需分辨。 这一处置留下的后续条件是：后续路线、控制与损失应与对方记录和遗址材料复核。

证据的边界是：《元史》为明初仓促官修，本纪保存大量实录条目但有编次与纪年问题；行省执行、数字、族称及对外关系须以条格、多语文书和对方史料互校。 蒙古征服、草原政治与跨区域文献比较研究 可用于继续检验这一结论。

\noindent\eventanchor{SYD-1225-MONGOL-070}{二月，武仙以真定叛，殺史天倪}\textbf{蒙古。}
蒙古 在 1225年（元中统-34年；公元换算据《元史》本纪纪年） 的材料痕迹是〔二月，武仙以真定叛，殺史天倪〕，出处为 《元史》卷001〈本纪第1卷本年条〉。材料未将地点细化到城址，不能用中枢所在地替它补足。

本纪从朝廷视角记录军政行动；出兵与战果需分辨。 记录所见的结果则是：后续路线、控制与损失应与对方记录和遗址材料复核。 日期相邻不等于因果已经证实。

《元史》为明初仓促官修，本纪保存大量实录条目但有编次与纪年问题；行省执行、数字、族称及对外关系须以条格、多语文书和对方史料互校。 因此，蒙古征服、草原政治与跨区域文献比较研究 只能扩展问题，不能代替未保存的细节。

\noindent\eventanchor{SYD-1225-MONGOL-118}{三月，史天澤擊仙走之，復真定}\textbf{蒙古。}
从 《元史》卷001〈本纪第1卷本年条〉 的 1225年（元中统-34年；公元换算据《元史》本纪纪年） 条目看，〔三月，史天澤擊仙走之，復真定〕使 蒙古 的一个具体行动进入叙事；材料未将地点细化到城址，不能用中枢所在地替它补足。

把本项放回事件链，能够看到的前提为：本纪年条记录政令或机构处置；实施背景须参照制度材料。 而它所打开或收紧的下一步是：可在同年及后续政令中追踪；条文不单独证明地方执行。

关于可说范围，须保留：《元史》为明初仓促官修，本纪保存大量实录条目但有编次与纪年问题；行省执行、数字、族称及对外关系须以条格、多语文书和对方史料互校。 进一步讨论可参照 蒙古征服、草原政治与跨区域文献比较研究

\subsection{1226年（元中统-33年；公元换算据《元史》本纪纪年）}
\noindent\eventanchor{SYD-1226-MONGOL-021}{丙寅，帝渡河擊夏師，敗之}\textbf{蒙古。}
1226年（元中统-33年；公元换算据《元史》本纪纪年） 的记载将 蒙古 与〔丙寅，帝渡河擊夏師，敗之〕相连。此处可确认的是这项被记录的行动；题名保留的城、路、水道或机构名称，是目前可以落地的空间线索。

它承接的条件是：本纪保留灾害或工程的报告地点与朝廷处置；灾情范围需另证。 随后可追踪的变化为：可与地方文书、河工和环境材料比较，不将单点记录外推为全域因果。

证据的边界是：《元史》为明初仓促官修，本纪保存大量实录条目但有编次与纪年问题；行省执行、数字、族称及对外关系须以条格、多语文书和对方史料互校。 蒙古征服、草原政治与跨区域文献比较研究 可用于继续检验这一结论。

\noindent\eventanchor{SYD-1226-MONGOL-071}{冬十一月庚申，帝攻靈州，夏遣嵬名令公來援}\textbf{蒙古。}
〔冬十一月庚申，帝攻靈州，夏遣嵬名令公來援〕见于 《元史》卷001〈本纪第1卷本年条〉 的 1226年（元中统-33年；公元换算据《元史》本纪纪年） 记录。蒙古 的选择因此有了可回查的时间位置；题名保留的城、路、水道或机构名称，是目前可以落地的空间线索。

前一层压力可由以下线索辨认：本纪从朝廷视角记录军政行动；出兵与战果需分辨。 这一处置留下的后续条件是：后续路线、控制与损失应与对方记录和遗址材料复核。

《元史》为明初仓促官修，本纪保存大量实录条目但有编次与纪年问题；行省执行、数字、族称及对外关系须以条格、多语文书和对方史料互校。 因此，蒙古征服、草原政治与跨区域文献比较研究 只能扩展问题，不能代替未保存的细节。

\noindent\eventanchor{SYD-1226-MONGOL-119}{二十一年〔丙戌〕春正月，帝以西夏納仇人〔亦〕臘喝翔昆及不遣質子，自將伐之}\textbf{蒙古与西夏。}
1226年（元中统-33年；公元换算据《元史》本纪纪年），《元史》卷001〈本纪第1卷本年条〉 所见的〔二十一年〔丙戌〕春正月，帝以西夏納仇人〔亦〕臘喝翔昆及不遣質子，自將伐之〕把 蒙古与西夏 的处置留在可复核的年序中；材料未将地点细化到城址，不能用中枢所在地替它补足。

本纪年条提供日期和中枢可见行动；前因不据单条补定。 记录所见的结果则是：可回查同年及后续条目，地方影响仍须另证。 日期相邻不等于因果已经证实。

关于可说范围，须保留：《元史》为明初仓促官修，本纪保存大量实录条目但有编次与纪年问题；行省执行、数字、族称及对外关系须以条格、多语文书和对方史料互校。 进一步讨论可参照 蒙古征服、草原政治与跨区域文献比较研究

\subsection{1227年（元中统-32年；公元换算据《元史》本纪纪年）}
\noindent\eventanchor{SYD-1227-MONGOL-022}{帝謂群臣曰：「朕自去冬五星聚時，已嘗許不殺掠，遽忘下詔耶}\textbf{蒙古。}
蒙古 在 1227年（元中统-32年；公元换算据《元史》本纪纪年） 的材料痕迹是〔帝謂群臣曰：「朕自去冬五星聚時，已嘗許不殺掠，遽忘下詔耶〕，出处为 《元史》卷001〈本纪第1卷本年条〉。材料未将地点细化到城址，不能用中枢所在地替它补足。

把本项放回事件链，能够看到的前提为：本纪年条记录政令或机构处置；实施背景须参照制度材料。 而它所打开或收紧的下一步是：可在同年及后续政令中追踪；条文不单独证明地方执行。

证据的边界是：《元史》为明初仓促官修，本纪保存大量实录条目但有编次与纪年问题；行省执行、数字、族称及对外关系须以条格、多语文书和对方史料互校。 蒙古征服、草原政治与跨区域文献比较研究 可用于继续检验这一结论。

\noindent\eventanchor{SYD-1227-MONGOL-072}{二十二年丁亥春，帝留兵攻夏王城，自率師渡河攻積石州}\textbf{蒙古。}
从 《元史》卷001〈本纪第1卷本年条〉 的 1227年（元中统-32年；公元换算据《元史》本纪纪年） 条目看，〔二十二年丁亥春，帝留兵攻夏王城，自率師渡河攻積石州〕使 蒙古 的一个具体行动进入叙事；题名保留的城、路、水道或机构名称，是目前可以落地的空间线索。

它承接的条件是：本纪保留灾害或工程的报告地点与朝廷处置；灾情范围需另证。 随后可追踪的变化为：可与地方文书、河工和环境材料比较，不将单点记录外推为全域因果。

《元史》为明初仓促官修，本纪保存大量实录条目但有编次与纪年问题；行省执行、数字、族称及对外关系须以条格、多语文书和对方史料互校。 因此，蒙古征服、草原政治与跨区域文献比较研究 只能扩展问题，不能代替未保存的细节。

\noindent\eventanchor{SYD-1227-MONGOL-120}{臨崩謂左右曰：「金精兵在潼關，南據連山，北限大河，難以遽破}\textbf{蒙古。}
1227年（元中统-32年；公元换算据《元史》本纪纪年） 的记载将 蒙古 与〔臨崩謂左右曰：「金精兵在潼關，南據連山，北限大河，難以遽破〕相连。此处可确认的是这项被记录的行动；题名保留的城、路、水道或机构名称，是目前可以落地的空间线索。

前一层压力可由以下线索辨认：本纪保留灾害或工程的报告地点与朝廷处置；灾情范围需另证。 这一处置留下的后续条件是：可与地方文书、河工和环境材料比较，不将单点记录外推为全域因果。

关于可说范围，须保留：《元史》为明初仓促官修，本纪保存大量实录条目但有编次与纪年问题；行省执行、数字、族称及对外关系须以条格、多语文书和对方史料互校。 进一步讨论可参照 蒙古征服、草原政治与跨区域文献比较研究

\noindent\eventanchor{SYD-1227-MONGOL-162}{六月，金遣完顏合周、奧屯阿虎來請和}\textbf{蒙古。}
〔六月，金遣完顏合周、奧屯阿虎來請和〕见于 《元史》卷001〈本纪第1卷本年条〉 的 1227年（元中统-32年；公元换算据《元史》本纪纪年） 记录。蒙古 的选择因此有了可回查的时间位置；材料未将地点细化到城址，不能用中枢所在地替它补足。

本纪记录使节、贡赐或往来，不等于稳定统属关系。 记录所见的结果则是：可与对方编年和交通、文书材料核对其后续落实。 日期相邻不等于因果已经证实。

证据的边界是：《元史》为明初仓促官修，本纪保存大量实录条目但有编次与纪年问题；行省执行、数字、族称及对外关系须以条格、多语文书和对方史料互校。 蒙古征服、草原政治与跨区域文献比较研究 可用于继续检验这一结论。

\noindent\eventanchor{SYD-1227-MONGOL-204}{遣斡陳那顏攻信都府，拔之}\textbf{蒙古。}
1227年（元中统-32年；公元换算据《元史》本纪纪年），《元史》卷001〈本纪第1卷本年条〉 所见的〔遣斡陳那顏攻信都府，拔之〕把 蒙古 的处置留在可复核的年序中；题名保留的城、路、水道或机构名称，是目前可以落地的空间线索。

把本项放回事件链，能够看到的前提为：本纪从朝廷视角记录军政行动；出兵与战果需分辨。 而它所打开或收紧的下一步是：后续路线、控制与损失应与对方记录和遗址材料复核。

《元史》为明初仓促官修，本纪保存大量实录条目但有编次与纪年问题；行省执行、数字、族称及对外关系须以条格、多语文书和对方史料互校。 因此，蒙古征服、草原政治与跨区域文献比较研究 只能扩展问题，不能代替未保存的细节。

\noindent\eventanchor{SYD-1227-MONGOL-240}{然以數萬之眾，千里赴援，人馬疲弊，雖至弗能戰，破之必矣}\textbf{蒙古。}
蒙古 在 1227年（元中统-32年；公元换算据《元史》本纪纪年） 的材料痕迹是〔然以數萬之眾，千里赴援，人馬疲弊，雖至弗能戰，破之必矣〕，出处为 《元史》卷001〈本纪第1卷本年条〉。材料未将地点细化到城址，不能用中枢所在地替它补足。

它承接的条件是：本纪从朝廷视角记录军政行动；出兵与战果需分辨。 随后可追踪的变化为：后续路线、控制与损失应与对方记录和遗址材料复核。

关于可说范围，须保留：《元史》为明初仓促官修，本纪保存大量实录条目但有编次与纪年问题；行省执行、数字、族称及对外关系须以条格、多语文书和对方史料互校。 进一步讨论可参照 蒙古征服、草原政治与跨区域文献比较研究

\noindent\eventanchor{SYD-1227-MONGOL-272}{三月，破洮、河、西寧二州}\textbf{蒙古。}
从 《元史》卷001〈本纪第1卷本年条〉 的 1227年（元中统-32年；公元换算据《元史》本纪纪年） 条目看，〔三月，破洮、河、西寧二州〕使 蒙古 的一个具体行动进入叙事；题名保留的城、路、水道或机构名称，是目前可以落地的空间线索。

前一层压力可由以下线索辨认：本纪保留灾害或工程的报告地点与朝廷处置；灾情范围需另证。 这一处置留下的后续条件是：可与地方文书、河工和环境材料比较，不将单点记录外推为全域因果。

证据的边界是：《元史》为明初仓促官修，本纪保存大量实录条目但有编次与纪年问题；行省执行、数字、族称及对外关系须以条格、多语文书和对方史料互校。 蒙古征服、草原政治与跨区域文献比较研究 可用于继续检验这一结论。

\subsection{1227年}
\noindent\eventanchor{YUAN-1227-EVENT-157}{西夏覆亡，同年成吉思汗去世}\textbf{蒙古与西夏。}
1227年 的记载将 蒙古与西夏 与〔西夏覆亡，同年成吉思汗去世〕相连。此处可确认的是这项被记录的行动；材料未将地点细化到城址，不能用中枢所在地替它补足。

蒙古诸部整合及对西夏、金的战争中前任统治的结束与宫廷、部族或官僚的权力安排相互交织，促成“西夏覆亡，同年成吉思汗去世”这一继承节点。 记录所见的结果则是：继承并不自动消除既有矛盾；“西夏覆亡，同年成吉思汗去世”后的任官、军权和对外策略需由后续条目追踪。 日期相邻不等于因果已经证实。

《元史》明初速修，纪年与人名有异同；蒙古大汗政权不等同所有汗国，征服、赋役和身份分类须按地区及材料语言分列。 因此，蒙古帝国、元朝行省财政、多语文书与区域社会研究 只能扩展问题，不能代替未保存的细节。

\subsection{1228 年（共享年序桥接）}
\noindent\eventanchor{SY-1228-CHRONOLOGY-BRIDGE}{【C级年序桥接】保留1228 年的多政权并行检索入口；本条不主张所列政权在当年发生直接互动，具体事件须另以单项材料入账。}\textbf{南宋、金、西夏与蒙古（年序桥接）。}
〔【C级年序桥接】保留1228 年的多政权并行检索入口；本条不主张所列政权在当年发生直接互动，具体事件须另以单项材料入账。〕见于 《宋史》；《金史》；《元史》；蒙古文、波斯文及地方材料（据事核） 的 1228 年（共享年序桥接） 记录。南宋、金、西夏与蒙古（年序桥接） 的选择因此有了可回查的时间位置；题名保留的城、路、水道或机构名称，是目前可以落地的空间线索。

把本项放回事件链，能够看到的前提为：相邻已入账事件之间缺少该年度的共享年序入口，易使跨政权材料被单一政权叙事遮蔽。 而它所打开或收紧的下一步是：保留该年度的检索与补账位置；后续仅在获得可核单项材料时拆分为具体政治、财政、军事、社会或文化事件。

关于可说范围，须保留：本条只确认该年位于可追踪的共享年序中；正史、地方文书和外部材料的保存密度不一，不能由桥接标签推断行动、统属、疆界或社会影响。 进一步讨论可参照 蒙古征服、区域政权转换与跨语种年序研究

\subsection{1229年（元中统-30年；公元换算据《元史》本纪纪年）}
\noindent\eventanchor{SYD-1229-MONGOL-023}{敕蒙古民有馬百者輸牝馬一，牛百者輸牸牛一，羊百者輸羒羊一，為永制}\textbf{蒙古。}
1229年（元中统-30年；公元换算据《元史》本纪纪年），《元史》卷002〈本纪第2卷本年条〉 所见的〔敕蒙古民有馬百者輸牝馬一，牛百者輸牸牛一，羊百者輸羒羊一，為永制〕把 蒙古 的处置留在可复核的年序中；材料未将地点细化到城址，不能用中枢所在地替它补足。

它承接的条件是：本纪所记财用、赈济或征发属于特定报告场景。 随后可追踪的变化为：可与食货志、文书和物质材料比较，不外推为全境总量。

证据的边界是：《元史》为明初仓促官修，本纪保存大量实录条目但有编次与纪年问题；行省执行、数字、族称及对外关系须以条格、多语文书和对方史料互校。 蒙古征服、草原政治与跨区域文献比较研究 可用于继续检验这一结论。

\noindent\eventanchor{SYD-1229-MONGOL-073}{帝雖御極，而朝政猶出於六皇后云}\textbf{蒙古。}
蒙古 在 1229年（元中统-30年；公元换算据《元史》本纪纪年） 的材料痕迹是〔帝雖御極，而朝政猶出於六皇后云〕，出处为 《元史》卷002〈本纪第2卷本年条〉。材料未将地点细化到城址，不能用中枢所在地替它补足。

前一层压力可由以下线索辨认：本纪年条提供日期和中枢可见行动；前因不据单条补定。 这一处置留下的后续条件是：可回查同年及后续条目，地方影响仍须另证。

《元史》为明初仓促官修，本纪保存大量实录条目但有编次与纪年问题；行省执行、数字、族称及对外关系须以条格、多语文书和对方史料互校。 因此，蒙古征服、草原政治与跨区域文献比较研究 只能扩展问题，不能代替未保存的细节。

\noindent\eventanchor{SYD-1229-MONGOL-121}{金遣阿虎帶來歸太祖之賵，帝曰：「汝主久不降，使先帝老于兵間，吾豈能忘也，賵何為哉！」却之}\textbf{蒙古。}
从 《元史》卷002〈本纪第2卷本年条〉 的 1229年（元中统-30年；公元换算据《元史》本纪纪年） 条目看，〔金遣阿虎帶來歸太祖之賵，帝曰：「汝主久不降，使先帝老于兵間，吾豈能忘也，賵何為哉！」却之〕使 蒙古 的一个具体行动进入叙事；材料未将地点细化到城址，不能用中枢所在地替它补足。

本纪从朝廷视角记录军政行动；出兵与战果需分辨。 记录所见的结果则是：后续路线、控制与损失应与对方记录和遗址材料复核。 日期相邻不等于因果已经证实。

关于可说范围，须保留：《元史》为明初仓促官修，本纪保存大量实录条目但有编次与纪年问题；行省执行、数字、族称及对外关系须以条格、多语文书和对方史料互校。 进一步讨论可参照 蒙古征服、草原政治与跨区域文献比较研究

\noindent\eventanchor{SYD-1229-MONGOL-163}{命河北漢民以戶計，出賦調，耶律楚材主之}\textbf{蒙古。}
1229年（元中统-30年；公元换算据《元史》本纪纪年） 的记载将 蒙古 与〔命河北漢民以戶計，出賦調，耶律楚材主之〕相连。此处可确认的是这项被记录的行动；题名保留的城、路、水道或机构名称，是目前可以落地的空间线索。

把本项放回事件链，能够看到的前提为：本纪年条记录政令或机构处置；实施背景须参照制度材料。 而它所打开或收紧的下一步是：可在同年及后续政令中追踪；条文不单独证明地方执行。

证据的边界是：《元史》为明初仓促官修，本纪保存大量实录条目但有编次与纪年问题；行省执行、数字、族称及对外关系须以条格、多语文书和对方史料互校。 蒙古征服、草原政治与跨区域文献比较研究 可用于继续检验这一结论。

\noindent\eventanchor{SYD-1229-MONGOL-205}{秋八月己未，諸王百官大會于怯綠連河曲雕阿蘭之地，以太祖遺詔即皇帝位于庫鐵烏阿剌里}\textbf{蒙古。}
〔秋八月己未，諸王百官大會于怯綠連河曲雕阿蘭之地，以太祖遺詔即皇帝位于庫鐵烏阿剌里〕见于 《元史》卷002〈本纪第2卷本年条〉 的 1229年（元中统-30年；公元换算据《元史》本纪纪年） 记录。蒙古 的选择因此有了可回查的时间位置；题名保留的城、路、水道或机构名称，是目前可以落地的空间线索。

它承接的条件是：本纪年条记录政令或机构处置；实施背景须参照制度材料。 随后可追踪的变化为：可在同年及后续政令中追踪；条文不单独证明地方执行。

《元史》为明初仓促官修，本纪保存大量实录条目但有编次与纪年问题；行省执行、数字、族称及对外关系须以条格、多语文书和对方史料互校。 因此，蒙古征服、草原政治与跨区域文献比较研究 只能扩展问题，不能代替未保存的细节。

\noindent\eventanchor{SYD-1229-MONGOL-241}{權萬戶史權等耀兵淮南，攻虎頭關寨，拔之，進圍黃州}\textbf{蒙古。}
1229年（元中统-30年；公元换算据《元史》本纪纪年），《元史》卷002〈本纪第2卷本年条〉 所见的〔權萬戶史權等耀兵淮南，攻虎頭關寨，拔之，進圍黃州〕把 蒙古 的处置留在可复核的年序中；题名保留的城、路、水道或机构名称，是目前可以落地的空间线索。

前一层压力可由以下线索辨认：本纪从朝廷视角记录军政行动；出兵与战果需分辨。 这一处置留下的后续条件是：后续路线、控制与损失应与对方记录和遗址材料复核。

关于可说范围，须保留：《元史》为明初仓促官修，本纪保存大量实录条目但有编次与纪年问题；行省执行、数字、族称及对外关系须以条格、多语文书和对方史料互校。 进一步讨论可参照 蒙古征服、草原政治与跨区域文献比较研究

\noindent\eventanchor{SYD-1229-MONGOL-273}{始立朝儀，皇族尊屬皆拜}\textbf{蒙古。}
蒙古 在 1229年（元中统-30年；公元换算据《元史》本纪纪年） 的材料痕迹是〔始立朝儀，皇族尊屬皆拜〕，出处为 《元史》卷002〈本纪第2卷本年条〉。材料未将地点细化到城址，不能用中枢所在地替它补足。

本纪年条提供日期和中枢可见行动；前因不据单条补定。 记录所见的结果则是：可回查同年及后续条目，地方影响仍须另证。 日期相邻不等于因果已经证实。

证据的边界是：《元史》为明初仓促官修，本纪保存大量实录条目但有编次与纪年问题；行省执行、数字、族称及对外关系须以条格、多语文书和对方史料互校。 蒙古征服、草原政治与跨区域文献比较研究 可用于继续检验这一结论。

\subsection{1229年}
\noindent\eventanchor{YUAN-1229-EVENT-158}{窝阔台即大汗位}\textbf{蒙古。}
从 《元史》相关本纪；《元朝秘史》、波斯文史籍或《高丽史》对应叙事；《元典章》《通制条格》、多语碑刻与文书（据事核） 的 1229年 条目看，〔窝阔台即大汗位〕使 蒙古 的一个具体行动进入叙事；材料未将地点细化到城址，不能用中枢所在地替它补足。

把本项放回事件链，能够看到的前提为：“窝阔台即大汗位”发生在蒙古诸部整合及对西夏、金的战争的具体压力与机会之中，相关行动者的选择不能脱离此前事件链解释。 而它所打开或收紧的下一步是：“窝阔台即大汗位”为后续政治、边防或资源配置留下条件；其社会影响与实际覆盖范围仍须以异类材料限定。

《元史》明初速修，纪年与人名有异同；蒙古大汗政权不等同所有汗国，征服、赋役和身份分类须按地区及材料语言分列。 因此，蒙古帝国、元朝行省财政、多语文书与区域社会研究 只能扩展问题，不能代替未保存的细节。

\subsection{1230年（元中统-29年；公元换算据《元史》本纪纪年）}
\noindent\eventanchor{SYD-1230-MONGOL-024}{八月，命野里知吉帶率搠思蠻部兵征西}\textbf{蒙古。}
1230年（元中统-29年；公元换算据《元史》本纪纪年） 的记载将 蒙古 与〔八月，命野里知吉帶率搠思蠻部兵征西〕相连。此处可确认的是这项被记录的行动；材料未将地点细化到城址，不能用中枢所在地替它补足。

它承接的条件是：本纪年条记录政令或机构处置；实施背景须参照制度材料。 随后可追踪的变化为：可在同年及后续政令中追踪；条文不单独证明地方执行。

关于可说范围，须保留：《元史》为明初仓促官修，本纪保存大量实录条目但有编次与纪年问题；行省执行、数字、族称及对外关系须以条格、多语文书和对方史料互校。 进一步讨论可参照 蒙古征服、草原政治与跨区域文献比较研究

\noindent\eventanchor{SYD-1230-MONGOL-074}{冬十一月，始置十路徵收課稅使，以陳時可、趙昉使燕京，劉中、劉桓使宣德，周立和、王貞使西京，呂振、劉子振使太原，楊簡、高廷英使平陽，王晉、賈從使真定，張瑜、王銳使東平，王德亨、侯顯使北京，夾谷永、程泰使平州，田木西、李天翼使濟南}\textbf{蒙古。}
〔冬十一月，始置十路徵收課稅使，以陳時可、趙昉使燕京，劉中、劉桓使宣德，周立和、王貞使西京，呂振、劉子振使太原，楊簡、高廷英使平陽，王晉、賈從使真定，張瑜、王銳使東平，王德亨、侯顯使北京，夾谷永、程泰使平州，田木西、李天翼使濟南〕见于 《元史》卷002〈本纪第2卷本年条〉 的 1230年（元中统-29年；公元换算据《元史》本纪纪年） 记录。蒙古 的选择因此有了可回查的时间位置；题名保留的城、路、水道或机构名称，是目前可以落地的空间线索。

前一层压力可由以下线索辨认：本纪年条记录政令或机构处置；实施背景须参照制度材料。 这一处置留下的后续条件是：可在同年及后续政令中追踪；条文不单独证明地方执行。

证据的边界是：《元史》为明初仓促官修，本纪保存大量实录条目但有编次与纪年问题；行省执行、数字、族称及对外关系须以条格、多语文书和对方史料互校。 蒙古征服、草原政治与跨区域文献比较研究 可用于继续检验这一结论。

\noindent\eventanchor{SYD-1230-MONGOL-122}{朶忽魯及金兵戰，敗績，命速不台援之}\textbf{蒙古与金。}
1230年（元中统-29年；公元换算据《元史》本纪纪年），《元史》卷002〈本纪第2卷本年条〉 所见的〔朶忽魯及金兵戰，敗績，命速不台援之〕把 蒙古与金 的处置留在可复核的年序中；材料未将地点细化到城址，不能用中枢所在地替它补足。

本纪年条记录政令或机构处置；实施背景须参照制度材料。 记录所见的结果则是：可在同年及后续政令中追踪；条文不单独证明地方执行。 日期相邻不等于因果已经证实。

《元史》为明初仓促官修，本纪保存大量实录条目但有编次与纪年问题；行省执行、数字、族称及对外关系须以条格、多语文书和对方史料互校。 因此，蒙古征服、草原政治与跨区域文献比较研究 只能扩展问题，不能代替未保存的细节。

\noindent\eventanchor{SYD-1230-MONGOL-164}{二年丁未春，張柔攻泗州}\textbf{蒙古。}
蒙古 在 1230年（元中统-29年；公元换算据《元史》本纪纪年） 的材料痕迹是〔二年丁未春，張柔攻泗州〕，出处为 《元史》卷002〈本纪第2卷本年条〉。题名保留的城、路、水道或机构名称，是目前可以落地的空间线索。

把本项放回事件链，能够看到的前提为：本纪从朝廷视角记录军政行动；出兵与战果需分辨。 而它所打开或收紧的下一步是：后续路线、控制与损失应与对方记录和遗址材料复核。

关于可说范围，须保留：《元史》为明初仓促官修，本纪保存大量实录条目但有编次与纪年问题；行省执行、数字、族称及对外关系须以条格、多语文书和对方史料互校。 进一步讨论可参照 蒙古征服、草原政治与跨区域文献比较研究

\noindent\eventanchor{SYD-1230-MONGOL-206}{二年庚寅春正月，詔自今以前事勿問}\textbf{蒙古。}
从 《元史》卷002〈本纪第2卷本年条〉 的 1230年（元中统-29年；公元换算据《元史》本纪纪年） 条目看，〔二年庚寅春正月，詔自今以前事勿問〕使 蒙古 的一个具体行动进入叙事；材料未将地点细化到城址，不能用中枢所在地替它补足。

它承接的条件是：本纪年条记录政令或机构处置；实施背景须参照制度材料。 随后可追踪的变化为：可在同年及后续政令中追踪；条文不单独证明地方执行。

证据的边界是：《元史》为明初仓促官修，本纪保存大量实录条目但有编次与纪年问题；行省执行、数字、族称及对外关系须以条格、多语文书和对方史料互校。 蒙古征服、草原政治与跨区域文献比较研究 可用于继续检验这一结论。

\noindent\eventanchor{SYD-1230-MONGOL-242}{金主率師來援，敗之，尋拔其城}\textbf{蒙古与金。}
1230年（元中统-29年；公元换算据《元史》本纪纪年） 的记载将 蒙古与金 与〔金主率師來援，敗之，尋拔其城〕相连。此处可确认的是这项被记录的行动；题名保留的城、路、水道或机构名称，是目前可以落地的空间线索。

前一层压力可由以下线索辨认：本纪从朝廷视角记录军政行动；出兵与战果需分辨。 这一处置留下的后续条件是：后续路线、控制与损失应与对方记录和遗址材料复核。

《元史》为明初仓促官修，本纪保存大量实录条目但有编次与纪年问题；行省执行、数字、族称及对外关系须以条格、多语文书和对方史料互校。 因此，蒙古征服、草原政治与跨区域文献比较研究 只能扩展问题，不能代替未保存的细节。

\noindent\eventanchor{SYD-1230-MONGOL-274}{秋七月，帝自將南伐，皇弟拖雷、皇姪蒙哥率師從，拔天成等堡，遂渡河攻鳳翔}\textbf{蒙古。}
〔秋七月，帝自將南伐，皇弟拖雷、皇姪蒙哥率師從，拔天成等堡，遂渡河攻鳳翔〕见于 《元史》卷002〈本纪第2卷本年条〉 的 1230年（元中统-29年；公元换算据《元史》本纪纪年） 记录。蒙古 的选择因此有了可回查的时间位置；题名保留的城、路、水道或机构名称，是目前可以落地的空间线索。

本纪保留灾害或工程的报告地点与朝廷处置；灾情范围需另证。 记录所见的结果则是：可与地方文书、河工和环境材料比较，不将单点记录外推为全域因果。 日期相邻不等于因果已经证实。

关于可说范围，须保留：《元史》为明初仓促官修，本纪保存大量实录条目但有编次与纪年问题；行省执行、数字、族称及对外关系须以条格、多语文书和对方史料互校。 进一步讨论可参照 蒙古征服、草原政治与跨区域文献比较研究

\subsection{1231年（元中统-28年；公元换算据《元史》本纪纪年）}
\noindent\eventanchor{SYD-1231-MONGOL-025}{定宗崩後，議所立未決}\textbf{蒙古。}
1231年（元中统-28年；公元换算据《元史》本纪纪年），《元史》卷002〈本纪第2卷本年条〉 所见的〔定宗崩後，議所立未決〕把 蒙古 的处置留在可复核的年序中；材料未将地点细化到城址，不能用中枢所在地替它补足。

把本项放回事件链，能够看到的前提为：本纪年条提供日期和中枢可见行动；前因不据单条补定。 而它所打开或收紧的下一步是：可回查同年及后续条目，地方影响仍须另证。

证据的边界是：《元史》为明初仓促官修，本纪保存大量实录条目但有编次与纪年问题；行省执行、数字、族称及对外关系须以条格、多语文书和对方史料互校。 蒙古征服、草原政治与跨区域文献比较研究 可用于继续检验这一结论。

\noindent\eventanchor{SYD-1231-MONGOL-075}{高麗王㬚遣其弟懷安公請降}\textbf{蒙古。}
蒙古 在 1231年（元中统-28年；公元换算据《元史》本纪纪年） 的材料痕迹是〔高麗王㬚遣其弟懷安公請降〕，出处为 《元史》卷002〈本纪第2卷本年条〉。材料未将地点细化到城址，不能用中枢所在地替它补足。

它承接的条件是：本纪从朝廷视角记录军政行动；出兵与战果需分辨。 随后可追踪的变化为：后续路线、控制与损失应与对方记录和遗址材料复核。

《元史》为明初仓促官修，本纪保存大量实录条目但有编次与纪年问题；行省执行、数字、族称及对外关系须以条格、多语文书和对方史料互校。 因此，蒙古征服、草原政治与跨区域文献比较研究 只能扩展问题，不能代替未保存的细节。

\noindent\eventanchor{SYD-1231-MONGOL-123}{遣搠不罕使宋假道，宋殺之}\textbf{蒙古与南宋。}
从 《元史》卷002〈本纪第2卷本年条〉 的 1231年（元中统-28年；公元换算据《元史》本纪纪年） 条目看，〔遣搠不罕使宋假道，宋殺之〕使 蒙古与南宋 的一个具体行动进入叙事；材料未将地点细化到城址，不能用中枢所在地替它补足。

前一层压力可由以下线索辨认：本纪从朝廷视角记录军政行动；出兵与战果需分辨。 这一处置留下的后续条件是：后续路线、控制与损失应与对方记录和遗址材料复核。

关于可说范围，须保留：《元史》为明初仓促官修，本纪保存大量实录条目但有编次与纪年问题；行省执行、数字、族称及对外关系须以条格、多语文书和对方史料互校。 进一步讨论可参照 蒙古征服、草原政治与跨区域文献比较研究

\noindent\eventanchor{SYD-1231-MONGOL-165}{然自壬寅以來，法度不一，內外離心，而太宗之政衰矣}\textbf{蒙古。}
1231年（元中统-28年；公元换算据《元史》本纪纪年） 的记载将 蒙古 与〔然自壬寅以來，法度不一，內外離心，而太宗之政衰矣〕相连。此处可确认的是这项被记录的行动；材料未将地点细化到城址，不能用中枢所在地替它补足。

本纪年条提供日期和中枢可见行动；前因不据单条补定。 记录所见的结果则是：可回查同年及后续条目，地方影响仍须另证。 日期相邻不等于因果已经证实。

证据的边界是：《元史》为明初仓促官修，本纪保存大量实录条目但有编次与纪年问题；行省执行、数字、族称及对外关系须以条格、多语文书和对方史料互校。 蒙古征服、草原政治与跨区域文献比较研究 可用于继续检验这一结论。

\noindent\eventanchor{SYD-1231-MONGOL-207}{三年辛卯春二月，克鳳翔，攻洛陽、河中諸城，下之}\textbf{蒙古。}
〔三年辛卯春二月，克鳳翔，攻洛陽、河中諸城，下之〕见于 《元史》卷002〈本纪第2卷本年条〉 的 1231年（元中统-28年；公元换算据《元史》本纪纪年） 记录。蒙古 的选择因此有了可回查的时间位置；题名保留的城、路、水道或机构名称，是目前可以落地的空间线索。

把本项放回事件链，能够看到的前提为：本纪保留灾害或工程的报告地点与朝廷处置；灾情范围需另证。 而它所打开或收紧的下一步是：可与地方文书、河工和环境材料比较，不将单点记录外推为全域因果。

《元史》为明初仓促官修，本纪保存大量实录条目但有编次与纪年问题；行省执行、数字、族称及对外关系须以条格、多语文书和对方史料互校。 因此，蒙古征服、草原政治与跨区域文献比较研究 只能扩展问题，不能代替未保存的细节。

\noindent\eventanchor{SYD-1231-MONGOL-243}{始立中書省，改侍從官名}\textbf{蒙古。}
1231年（元中统-28年；公元换算据《元史》本纪纪年），《元史》卷002〈本纪第2卷本年条〉 所见的〔始立中書省，改侍從官名〕把 蒙古 的处置留在可复核的年序中；材料未将地点细化到城址，不能用中枢所在地替它补足。

它承接的条件是：本纪年条记录政令或机构处置；实施背景须参照制度材料。 随后可追踪的变化为：可在同年及后续政令中追踪；条文不单独证明地方执行。

关于可说范围，须保留：《元史》为明初仓促官修，本纪保存大量实录条目但有编次与纪年问题；行省执行、数字、族称及对外关系须以条格、多语文书和对方史料互校。 进一步讨论可参照 蒙古征服、草原政治与跨区域文献比较研究

\noindent\eventanchor{SYD-1231-MONGOL-275}{是歲大旱，河水盡涸，野草自焚，牛馬十死八九，人不聊生}\textbf{蒙古。}
蒙古 在 1231年（元中统-28年；公元换算据《元史》本纪纪年） 的材料痕迹是〔是歲大旱，河水盡涸，野草自焚，牛馬十死八九，人不聊生〕，出处为 《元史》卷002〈本纪第2卷本年条〉。题名保留的城、路、水道或机构名称，是目前可以落地的空间线索。

前一层压力可由以下线索辨认：本纪保留灾害或工程的报告地点与朝廷处置；灾情范围需另证。 这一处置留下的后续条件是：可与地方文书、河工和环境材料比较，不将单点记录外推为全域因果。

证据的边界是：《元史》为明初仓促官修，本纪保存大量实录条目但有编次与纪年问题；行省执行、数字、族称及对外关系须以条格、多语文书和对方史料互校。 蒙古征服、草原政治与跨区域文献比较研究 可用于继续检验这一结论。

\subsection{1231年}
\noindent\eventanchor{YUAN-1231-EVENT-159}{三峰山等战事削弱金军主力}\textbf{蒙古与金。}
从 《元史》相关本纪；《元朝秘史》、波斯文史籍或《高丽史》对应叙事；《元典章》《通制条格》、多语碑刻与文书（据事核） 的 1231年 条目看，〔三峰山等战事削弱金军主力〕使 蒙古与金 的一个具体行动进入叙事；题名保留的城、路、水道或机构名称，是目前可以落地的空间线索。

蒙古诸部整合及对西夏、金的战争中既有的联盟、交通与补给压力，为“三峰山等战事削弱金军主力”提供了直接的军事或动员背景。 记录所见的结果则是：“三峰山等战事削弱金军主力”改变了相关城镇、道路或政权间的军事选择；伤亡、俘获和控制范围仅可按各方材料分别确认。 日期相邻不等于因果已经证实。

《元史》明初速修，纪年与人名有异同；蒙古大汗政权不等同所有汗国，征服、赋役和身份分类须按地区及材料语言分列。 因此，蒙古帝国、元朝行省财政、多语文书与区域社会研究 只能扩展问题，不能代替未保存的细节。

\subsection{1232年（元中统-27年；公元换算据《元史》本纪纪年）}
\noindent\eventanchor{SYD-1232-MONGOL-026}{八月，撒禮塔復征高麗，中矢卒}\textbf{蒙古。}
1232年（元中统-27年；公元换算据《元史》本纪纪年） 的记载将 蒙古 与〔八月，撒禮塔復征高麗，中矢卒〕相连。此处可确认的是这项被记录的行动；材料未将地点细化到城址，不能用中枢所在地替它补足。

把本项放回事件链，能够看到的前提为：本纪年条记录政令或机构处置；实施背景须参照制度材料。 而它所打开或收紧的下一步是：可在同年及后续政令中追踪；条文不单独证明地方执行。

关于可说范围，须保留：《元史》为明初仓促官修，本纪保存大量实录条目但有编次与纪年问题；行省执行、数字、族称及对外关系须以条格、多语文书和对方史料互校。 进一步讨论可参照 蒙古征服、草原政治与跨区域文献比较研究

\noindent\eventanchor{SYD-1232-MONGOL-076}{冬十一月，獵于納蘭赤剌溫之野}\textbf{蒙古。}
〔冬十一月，獵于納蘭赤剌溫之野〕见于 《元史》卷002〈本纪第2卷本年条〉 的 1232年（元中统-27年；公元换算据《元史》本纪纪年） 记录。蒙古 的选择因此有了可回查的时间位置；材料未将地点细化到城址，不能用中枢所在地替它补足。

它承接的条件是：本纪年条提供日期和中枢可见行动；前因不据单条补定。 随后可追踪的变化为：可回查同年及后续条目，地方影响仍须另证。

证据的边界是：《元史》为明初仓促官修，本纪保存大量实录条目但有编次与纪年问题；行省执行、数字、族称及对外关系须以条格、多语文书和对方史料互校。 蒙古征服、草原政治与跨区域文献比较研究 可用于继续检验这一结论。

\noindent\eventanchor{SYD-1232-MONGOL-124}{高麗叛，殺所置官吏，徙居江華島}\textbf{蒙古。}
1232年（元中统-27年；公元换算据《元史》本纪纪年），《元史》卷002〈本纪第2卷本年条〉 所见的〔高麗叛，殺所置官吏，徙居江華島〕把 蒙古 的处置留在可复核的年序中；题名保留的城、路、水道或机构名称，是目前可以落地的空间线索。

前一层压力可由以下线索辨认：本纪年条记录政令或机构处置；实施背景须参照制度材料。 这一处置留下的后续条件是：可在同年及后续政令中追踪；条文不单独证明地方执行。

《元史》为明初仓促官修，本纪保存大量实录条目但有编次与纪年问题；行省执行、数字、族称及对外关系须以条格、多语文书和对方史料互校。 因此，蒙古征服、草原政治与跨区域文献比较研究 只能扩展问题，不能代替未保存的细节。

\noindent\eventanchor{SYD-1232-MONGOL-166}{庚寅，拖雷渡漢江，遣使來報，即詔諸軍進發}\textbf{蒙古。}
蒙古 在 1232年（元中统-27年；公元换算据《元史》本纪纪年） 的材料痕迹是〔庚寅，拖雷渡漢江，遣使來報，即詔諸軍進發〕，出处为 《元史》卷002〈本纪第2卷本年条〉。题名保留的城、路、水道或机构名称，是目前可以落地的空间线索。

本纪年条记录政令或机构处置；实施背景须参照制度材料。 记录所见的结果则是：可在同年及后续政令中追踪；条文不单独证明地方执行。 日期相邻不等于因果已经证实。

关于可说范围，须保留：《元史》为明初仓促官修，本纪保存大量实录条目但有编次与纪年问题；行省执行、数字、族称及对外关系须以条格、多语文书和对方史料互校。 进一步讨论可参照 蒙古征服、草原政治与跨区域文献比较研究

\noindent\eventanchor{SYD-1232-MONGOL-208}{金參政完顏思烈、恒山公武仙救南京，諸軍與戰，敗之}\textbf{蒙古。}
从 《元史》卷002〈本纪第2卷本年条〉 的 1232年（元中统-27年；公元换算据《元史》本纪纪年） 条目看，〔金參政完顏思烈、恒山公武仙救南京，諸軍與戰，敗之〕使 蒙古 的一个具体行动进入叙事；题名保留的城、路、水道或机构名称，是目前可以落地的空间线索。

把本项放回事件链，能够看到的前提为：本纪从朝廷视角记录军政行动；出兵与战果需分辨。 而它所打开或收紧的下一步是：后续路线、控制与损失应与对方记录和遗址材料复核。

证据的边界是：《元史》为明初仓促官修，本纪保存大量实录条目但有编次与纪年问题；行省执行、数字、族称及对外关系须以条格、多语文书和对方史料互校。 蒙古征服、草原政治与跨区域文献比较研究 可用于继续检验这一结论。

\noindent\eventanchor{SYD-1232-MONGOL-244}{金防城提控馬伯堅降，授伯堅金符，使守之}\textbf{蒙古。}
1232年（元中统-27年；公元换算据《元史》本纪纪年） 的记载将 蒙古 与〔金防城提控馬伯堅降，授伯堅金符，使守之〕相连。此处可确认的是这项被记录的行动；题名保留的城、路、水道或机构名称，是目前可以落地的空间线索。

它承接的条件是：本纪从朝廷视角记录军政行动；出兵与战果需分辨。 随后可追踪的变化为：后续路线、控制与损失应与对方记录和遗址材料复核。

《元史》为明初仓促官修，本纪保存大量实录条目但有编次与纪年问题；行省执行、数字、族称及对外关系须以条格、多语文书和对方史料互校。 因此，蒙古征服、草原政治与跨区域文献比较研究 只能扩展问题，不能代替未保存的细节。

\noindent\eventanchor{SYD-1232-MONGOL-276}{秋七月，遣唐慶使金諭降，金殺之}\textbf{蒙古。}
〔秋七月，遣唐慶使金諭降，金殺之〕见于 《元史》卷002〈本纪第2卷本年条〉 的 1232年（元中统-27年；公元换算据《元史》本纪纪年） 记录。蒙古 的选择因此有了可回查的时间位置；材料未将地点细化到城址，不能用中枢所在地替它补足。

前一层压力可由以下线索辨认：本纪从朝廷视角记录军政行动；出兵与战果需分辨。 这一处置留下的后续条件是：后续路线、控制与损失应与对方记录和遗址材料复核。

关于可说范围，须保留：《元史》为明初仓促官修，本纪保存大量实录条目但有编次与纪年问题；行省执行、数字、族称及对外关系须以条格、多语文书和对方史料互校。 进一步讨论可参照 蒙古征服、草原政治与跨区域文献比较研究

\subsection{1233年（元中统-26年；公元换算据《元史》本纪纪年）}
\noindent\eventanchor{SYD-1233-MONGOL-027}{冬十一月，宋遣荊鄂都統孟珙以兵糧來助}\textbf{蒙古与南宋。}
1233年（元中统-26年；公元换算据《元史》本纪纪年），《元史》卷002〈本纪第2卷本年条〉 所见的〔冬十一月，宋遣荊鄂都統孟珙以兵糧來助〕把 蒙古与南宋 的处置留在可复核的年序中；题名保留的城、路、水道或机构名称，是目前可以落地的空间线索。

本纪年条提供日期和中枢可见行动；前因不据单条补定。 记录所见的结果则是：可回查同年及后续条目，地方影响仍须另证。 日期相邻不等于因果已经证实。

证据的边界是：《元史》为明初仓促官修，本纪保存大量实录条目但有编次与纪年问题；行省执行、数字、族称及对外关系须以条格、多语文书和对方史料互校。 蒙古征服、草原政治与跨区域文献比较研究 可用于继续检验这一结论。

\noindent\eventanchor{SYD-1233-MONGOL-077}{金人以海、沂、萊、濰等州降}\textbf{蒙古与金。}
蒙古与金 在 1233年（元中统-26年；公元换算据《元史》本纪纪年） 的材料痕迹是〔金人以海、沂、萊、濰等州降〕，出处为 《元史》卷002〈本纪第2卷本年条〉。题名保留的城、路、水道或机构名称，是目前可以落地的空间线索。

把本项放回事件链，能够看到的前提为：本纪从朝廷视角记录军政行动；出兵与战果需分辨。 而它所打开或收紧的下一步是：后续路线、控制与损失应与对方记录和遗址材料复核。

《元史》为明初仓促官修，本纪保存大量实录条目但有编次与纪年问题；行省执行、数字、族称及对外关系须以条格、多语文书和对方史料互校。 因此，蒙古征服、草原政治与跨区域文献比较研究 只能扩展问题，不能代替未保存的细节。

\noindent\eventanchor{SYD-1233-MONGOL-125}{十二月，諸軍與宋兵合攻蔡，敗武仙于息州}\textbf{蒙古与南宋。}
从 《元史》卷002〈本纪第2卷本年条〉 的 1233年（元中统-26年；公元换算据《元史》本纪纪年） 条目看，〔十二月，諸軍與宋兵合攻蔡，敗武仙于息州〕使 蒙古与南宋 的一个具体行动进入叙事；题名保留的城、路、水道或机构名称，是目前可以落地的空间线索。

它承接的条件是：本纪从朝廷视角记录军政行动；出兵与战果需分辨。 随后可追踪的变化为：后续路线、控制与损失应与对方记录和遗址材料复核。

关于可说范围，须保留：《元史》为明初仓促官修，本纪保存大量实录条目但有编次与纪年问题；行省执行、数字、族称及对外关系须以条格、多语文书和对方史料互校。 进一步讨论可参照 蒙古征服、草原政治与跨区域文献比较研究

\noindent\eventanchor{SYD-1233-MONGOL-167}{戊辰，金西面元帥崔立殺留守完顏奴申、完顏習捏阿不，以南京降}\textbf{蒙古。}
1233年（元中统-26年；公元换算据《元史》本纪纪年） 的记载将 蒙古 与〔戊辰，金西面元帥崔立殺留守完顏奴申、完顏習捏阿不，以南京降〕相连。此处可确认的是这项被记录的行动；题名保留的城、路、水道或机构名称，是目前可以落地的空间线索。

前一层压力可由以下线索辨认：本纪从朝廷视角记录军政行动；出兵与战果需分辨。 这一处置留下的后续条件是：后续路线、控制与损失应与对方记录和遗址材料复核。

证据的边界是：《元史》为明初仓促官修，本纪保存大量实录条目但有编次与纪年问题；行省执行、数字、族称及对外关系须以条格、多语文书和对方史料互校。 蒙古征服、草原政治与跨区域文献比较研究 可用于继续检验这一结论。

\noindent\eventanchor{SYD-1233-MONGOL-209}{夏四月，速不台進至青城，崔立以金太后王氏、后徒單氏及〔梁〕王從恪、〔荊〕王守純等至軍中，速不台遣送行在，遂入南京}\textbf{蒙古。}
〔夏四月，速不台進至青城，崔立以金太后王氏、后徒單氏及〔梁〕王從恪、〔荊〕王守純等至軍中，速不台遣送行在，遂入南京〕见于 《元史》卷002〈本纪第2卷本年条〉 的 1233年（元中统-26年；公元换算据《元史》本纪纪年） 记录。蒙古 的选择因此有了可回查的时间位置；题名保留的城、路、水道或机构名称，是目前可以落地的空间线索。

本纪年条提供日期和中枢可见行动；前因不据单条补定。 记录所见的结果则是：可回查同年及后续条目，地方影响仍须另证。 日期相邻不等于因果已经证实。

《元史》为明初仓促官修，本纪保存大量实录条目但有编次与纪年问题；行省执行、数字、族称及对外关系须以条格、多语文书和对方史料互校。 因此，蒙古征服、草原政治与跨区域文献比较研究 只能扩展问题，不能代替未保存的细节。

\noindent\eventanchor{SYD-1233-MONGOL-245}{詔以孔子五十一世孫元〔措〕襲封衍聖公}\textbf{蒙古。}
1233年（元中统-26年；公元换算据《元史》本纪纪年），《元史》卷002〈本纪第2卷本年条〉 所见的〔詔以孔子五十一世孫元〔措〕襲封衍聖公〕把 蒙古 的处置留在可复核的年序中；材料未将地点细化到城址，不能用中枢所在地替它补足。

把本项放回事件链，能够看到的前提为：本纪年条记录政令或机构处置；实施背景须参照制度材料。 而它所打开或收紧的下一步是：可在同年及后续政令中追踪；条文不单独证明地方执行。

关于可说范围，须保留：《元史》为明初仓促官修，本纪保存大量实录条目但有编次与纪年问题；行省执行、数字、族称及对外关系须以条格、多语文书和对方史料互校。 进一步讨论可参照 蒙古征服、草原政治与跨区域文献比较研究

\noindent\eventanchor{SYD-1233-MONGOL-277}{詔諸王議伐萬奴，遂命皇子貴由及諸王按赤帶將左翼軍討之}\textbf{蒙古。}
蒙古 在 1233年（元中统-26年；公元换算据《元史》本纪纪年） 的材料痕迹是〔詔諸王議伐萬奴，遂命皇子貴由及諸王按赤帶將左翼軍討之〕，出处为 《元史》卷002〈本纪第2卷本年条〉。材料未将地点细化到城址，不能用中枢所在地替它补足。

它承接的条件是：本纪年条记录政令或机构处置；实施背景须参照制度材料。 随后可追踪的变化为：可在同年及后续政令中追踪；条文不单独证明地方执行。

证据的边界是：《元史》为明初仓促官修，本纪保存大量实录条目但有编次与纪年问题；行省执行、数字、族称及对外关系须以条格、多语文书和对方史料互校。 蒙古征服、草原政治与跨区域文献比较研究 可用于继续检验这一结论。

\subsection{1234年（元中统-25年；公元换算据《元史》本纪纪年）}
\noindent\eventanchor{SYD-1234-MONGOL-028}{但盜馬一二者，即論死}\textbf{蒙古。}
从 《元史》卷002〈本纪第2卷本年条〉 的 1234年（元中统-25年；公元换算据《元史》本纪纪年） 条目看，〔但盜馬一二者，即論死〕使 蒙古 的一个具体行动进入叙事；材料未将地点细化到城址，不能用中枢所在地替它补足。

前一层压力可由以下线索辨认：本纪年条提供日期和中枢可见行动；前因不据单条补定。 这一处置留下的后续条件是：可回查同年及后续条目，地方影响仍须另证。

《元史》为明初仓促官修，本纪保存大量实录条目但有编次与纪年问题；行省执行、数字、族称及对外关系须以条格、多语文书和对方史料互校。 因此，蒙古征服、草原政治与跨区域文献比较研究 只能扩展问题，不能代替未保存的细节。

\noindent\eventanchor{SYD-1234-MONGOL-078}{凡來會，用善馬五十匹為一羈，守者五人，飼羸馬三人，守乞烈思三人}\textbf{蒙古。}
1234年（元中统-25年；公元换算据《元史》本纪纪年） 的记载将 蒙古 与〔凡來會，用善馬五十匹為一羈，守者五人，飼羸馬三人，守乞烈思三人〕相连。此处可确认的是这项被记录的行动；材料未将地点细化到城址，不能用中枢所在地替它补足。

本纪年条提供日期和中枢可见行动；前因不据单条补定。 记录所见的结果则是：可回查同年及后续条目，地方影响仍须另证。 日期相邻不等于因果已经证实。

关于可说范围，须保留：《元史》为明初仓促官修，本纪保存大量实录条目但有编次与纪年问题；行省执行、数字、族称及对外关系须以条格、多语文书和对方史料互校。 进一步讨论可参照 蒙古征服、草原政治与跨区域文献比较研究

\noindent\eventanchor{SYD-1234-MONGOL-126}{今後來會諸軍，甲內數不足，於近翼抽〔補〕足之}\textbf{蒙古。}
〔今後來會諸軍，甲內數不足，於近翼抽〔補〕足之〕见于 《元史》卷002〈本纪第2卷本年条〉 的 1234年（元中统-25年；公元换算据《元史》本纪纪年） 记录。蒙古 的选择因此有了可回查的时间位置；材料未将地点细化到城址，不能用中枢所在地替它补足。

把本项放回事件链，能够看到的前提为：本纪年条提供日期和中枢可见行动；前因不据单条补定。 而它所打开或收紧的下一步是：可回查同年及后续条目，地方影响仍须另证。

证据的边界是：《元史》为明初仓促官修，本纪保存大量实录条目但有编次与纪年问题；行省执行、数字、族称及对外关系须以条格、多语文书和对方史料互校。 蒙古征服、草原政治与跨区域文献比较研究 可用于继续检验这一结论。

\noindent\eventanchor{SYD-1234-MONGOL-168}{軍中凡十人置甲長，聽其指揮，專擅者論罪}\textbf{蒙古。}
1234年（元中统-25年；公元换算据《元史》本纪纪年），《元史》卷002〈本纪第2卷本年条〉 所见的〔軍中凡十人置甲長，聽其指揮，專擅者論罪〕把 蒙古 的处置留在可复核的年序中；材料未将地点细化到城址，不能用中枢所在地替它补足。

它承接的条件是：本纪年条记录政令或机构处置；实施背景须参照制度材料。 随后可追踪的变化为：可在同年及后续政令中追踪；条文不单独证明地方执行。

《元史》为明初仓促官修，本纪保存大量实录条目但有编次与纪年问题；行省执行、数字、族称及对外关系须以条格、多语文书和对方史料互校。 因此，蒙古征服、草原政治与跨区域文献比较研究 只能扩展问题，不能代替未保存的细节。

\noindent\eventanchor{SYD-1234-MONGOL-210}{其甲長以事來宮中，即置權攝一人、甲外一人，二人不得擅自往來，違者罪之}\textbf{蒙古。}
蒙古 在 1234年（元中统-25年；公元换算据《元史》本纪纪年） 的材料痕迹是〔其甲長以事來宮中，即置權攝一人、甲外一人，二人不得擅自往來，違者罪之〕，出处为 《元史》卷002〈本纪第2卷本年条〉。题名保留的城、路、水道或机构名称，是目前可以落地的空间线索。

前一层压力可由以下线索辨认：本纪年条记录政令或机构处置；实施背景须参照制度材料。 这一处置留下的后续条件是：可在同年及后续政令中追踪；条文不单独证明地方执行。

关于可说范围，须保留：《元史》为明初仓促官修，本纪保存大量实录条目但有编次与纪年问题；行省执行、数字、族称及对外关系须以条格、多语文书和对方史料互校。 进一步讨论可参照 蒙古征服、草原政治与跨区域文献比较研究

\noindent\eventanchor{SYD-1234-MONGOL-246}{是秋，帝在八里里答闌答八思之地，議自將伐宋，國王查老溫請行，遂遣之}\textbf{蒙古与南宋。}
从 《元史》卷002〈本纪第2卷本年条〉 的 1234年（元中统-25年；公元换算据《元史》本纪纪年） 条目看，〔是秋，帝在八里里答闌答八思之地，議自將伐宋，國王查老溫請行，遂遣之〕使 蒙古与南宋 的一个具体行动进入叙事；材料未将地点细化到城址，不能用中枢所在地替它补足。

本纪年条提供日期和中枢可见行动；前因不据单条补定。 记录所见的结果则是：可回查同年及后续条目，地方影响仍须另证。 日期相邻不等于因果已经证实。

证据的边界是：《元史》为明初仓促官修，本纪保存大量实录条目但有编次与纪年问题；行省执行、数字、族称及对外关系须以条格、多语文书和对方史料互校。 蒙古征服、草原政治与跨区域文献比较研究 可用于继续检验这一结论。

\noindent\eventanchor{SYD-1234-MONGOL-278}{諸出入宮禁，各有從者，男女止以十人為朋，出入毋得相雜}\textbf{蒙古。}
1234年（元中统-25年；公元换算据《元史》本纪纪年） 的记载将 蒙古 与〔諸出入宮禁，各有從者，男女止以十人為朋，出入毋得相雜〕相连。此处可确认的是这项被记录的行动；题名保留的城、路、水道或机构名称，是目前可以落地的空间线索。

把本项放回事件链，能够看到的前提为：本纪年条记录政令或机构处置；实施背景须参照制度材料。 而它所打开或收紧的下一步是：可在同年及后续政令中追踪；条文不单独证明地方执行。

《元史》为明初仓促官修，本纪保存大量实录条目但有编次与纪年问题；行省执行、数字、族称及对外关系须以条格、多语文书和对方史料互校。 因此，蒙古征服、草原政治与跨区域文献比较研究 只能扩展问题，不能代替未保存的细节。

\subsection{1234年（蒙古窝阔台六年、金天兴三年、宋端平元年；干支、月日待核；公元换算据各方本纪年号对照）}
\noindent\eventanchor{YUAN-1234-EVENT-160}{蒙古与南宋攻灭金朝}\textbf{蒙古南宋金。}
〔蒙古与南宋攻灭金朝〕见于 《元史》相关本纪；《元朝秘史》、波斯文史籍或《高丽史》对应叙事；《元典章》《通制条格》、多语碑刻与文书（据事核） 的 1234年（蒙古窝阔台六年、金天兴三年、宋端平元年；干支、月日待核；公元换算据各方本纪年号对照） 记录。蒙古南宋金 的选择因此有了可回查的时间位置；材料未将地点细化到城址，不能用中枢所在地替它补足。

它承接的条件是：蒙古诸部整合及对西夏、金的战争中既有的联盟、交通与补给压力，为“蒙古与南宋攻灭金朝”提供了直接的军事或动员背景。 随后可追踪的变化为：“蒙古与南宋攻灭金朝”改变了相关城镇、道路或政权间的军事选择；伤亡、俘获和控制范围仅可按各方材料分别确认。

关于可说范围，须保留：《元史》明初速修，纪年与人名有异同；蒙古大汗政权不等同所有汗国，征服、赋役和身份分类须按地区及材料语言分列。 进一步讨论可参照 蒙古帝国、元朝行省财政、多语文书与区域社会研究

\subsection{1235年（元中统-24年；公元换算据《元史》本纪纪年）}
\noindent\eventanchor{SYD-1235-MONGOL-029}{冬十月，曲出圍棗陽，拔之，遂徇襄、鄧，入郢，虜人民牛馬數萬而還}\textbf{蒙古。}
1235年（元中统-24年；公元换算据《元史》本纪纪年），《元史》卷002〈本纪第2卷本年条〉 所见的〔冬十月，曲出圍棗陽，拔之，遂徇襄、鄧，入郢，虜人民牛馬數萬而還〕把 蒙古 的处置留在可复核的年序中；材料未将地点细化到城址，不能用中枢所在地替它补足。

前一层压力可由以下线索辨认：本纪从朝廷视角记录军政行动；出兵与战果需分辨。 这一处置留下的后续条件是：后续路线、控制与损失应与对方记录和遗址材料复核。

证据的边界是：《元史》为明初仓促官修，本纪保存大量实录条目但有编次与纪年问题；行省执行、数字、族称及对外关系须以条格、多语文书和对方史料互校。 蒙古征服、草原政治与跨区域文献比较研究 可用于继续检验这一结论。

\noindent\eventanchor{SYD-1235-MONGOL-079}{遣諸王拔都及皇子貴由、皇姪蒙哥征西域，皇子闊端征秦、鞏，皇子曲出及胡土虎伐宋，唐古征高麗}\textbf{蒙古与南宋。}
蒙古与南宋 在 1235年（元中统-24年；公元换算据《元史》本纪纪年） 的材料痕迹是〔遣諸王拔都及皇子貴由、皇姪蒙哥征西域，皇子闊端征秦、鞏，皇子曲出及胡土虎伐宋，唐古征高麗〕，出处为 《元史》卷002〈本纪第2卷本年条〉。题名保留的城、路、水道或机构名称，是目前可以落地的空间线索。

本纪年条提供日期和中枢可见行动；前因不据单条补定。 记录所见的结果则是：可回查同年及后续条目，地方影响仍须另证。 日期相邻不等于因果已经证实。

《元史》为明初仓促官修，本纪保存大量实录条目但有编次与纪年问题；行省执行、数字、族称及对外关系须以条格、多语文书和对方史料互校。 因此，蒙古征服、草原政治与跨区域文献比较研究 只能扩展问题，不能代替未保存的细节。

\noindent\eventanchor{SYD-1235-MONGOL-127}{十一月，闊端攻石門，金便宜都總帥汪世顯降}\textbf{蒙古。}
从 《元史》卷002〈本纪第2卷本年条〉 的 1235年（元中统-24年；公元换算据《元史》本纪纪年） 条目看，〔十一月，闊端攻石門，金便宜都總帥汪世顯降〕使 蒙古 的一个具体行动进入叙事；题名保留的城、路、水道或机构名称，是目前可以落地的空间线索。

把本项放回事件链，能够看到的前提为：本纪从朝廷视角记录军政行动；出兵与战果需分辨。 而它所打开或收紧的下一步是：后续路线、控制与损失应与对方记录和遗址材料复核。

关于可说范围，须保留：《元史》为明初仓促官修，本纪保存大量实录条目但有编次与纪年问题；行省执行、数字、族称及对外关系须以条格、多语文书和对方史料互校。 进一步讨论可参照 蒙古征服、草原政治与跨区域文献比较研究

\noindent\eventanchor{SYD-1235-MONGOL-169}{中書省臣請契勘大明曆，從之}\textbf{蒙古。}
1235年（元中统-24年；公元换算据《元史》本纪纪年） 的记载将 蒙古 与〔中書省臣請契勘大明曆，從之〕相连。此处可确认的是这项被记录的行动；材料未将地点细化到城址，不能用中枢所在地替它补足。

它承接的条件是：本纪年条记录政令或机构处置；实施背景须参照制度材料。 随后可追踪的变化为：可在同年及后续政令中追踪；条文不单独证明地方执行。

证据的边界是：《元史》为明初仓促官修，本纪保存大量实录条目但有编次与纪年问题；行省执行、数字、族称及对外关系须以条格、多语文书和对方史料互校。 蒙古征服、草原政治与跨区域文献比较研究 可用于继续检验这一结论。

\subsection{1235年}
\noindent\eventanchor{YUAN-1235-EVENT-161}{蒙古贵族会议决定多方向扩张}\textbf{蒙古。}
〔蒙古贵族会议决定多方向扩张〕见于 《元史》相关本纪；《元朝秘史》、波斯文史籍或《高丽史》对应叙事；《元典章》《通制条格》、多语碑刻与文书（据事核） 的 1235年 记录。蒙古 的选择因此有了可回查的时间位置；材料未将地点细化到城址，不能用中枢所在地替它补足。

前一层压力可由以下线索辨认：“蒙古贵族会议决定多方向扩张”发生在蒙古诸部整合及对西夏、金的战争的具体压力与机会之中，相关行动者的选择不能脱离此前事件链解释。 这一处置留下的后续条件是：“蒙古贵族会议决定多方向扩张”为后续政治、边防或资源配置留下条件；其社会影响与实际覆盖范围仍须以异类材料限定。

《元史》明初速修，纪年与人名有异同；蒙古大汗政权不等同所有汗国，征服、赋役和身份分类须按地区及材料语言分列。 因此，蒙古帝国、元朝行省财政、多语文书与区域社会研究 只能扩展问题，不能代替未保存的细节。

\subsection{1236 年（共享年序桥接）}
\noindent\eventanchor{SY-1236-CHRONOLOGY-BRIDGE}{【C级年序桥接】保留1236 年的多政权并行检索入口；本条不主张所列政权在当年发生直接互动，具体事件须另以单项材料入账。}\textbf{南宋、金、西夏与蒙古（年序桥接）。}
1236 年（共享年序桥接），《宋史》；《金史》；《元史》；蒙古文、波斯文及地方材料（据事核） 所见的〔【C级年序桥接】保留1236 年的多政权并行检索入口；本条不主张所列政权在当年发生直接互动，具体事件须另以单项材料入账。〕把 南宋、金、西夏与蒙古（年序桥接） 的处置留在可复核的年序中；题名保留的城、路、水道或机构名称，是目前可以落地的空间线索。

相邻已入账事件之间缺少该年度的共享年序入口，易使跨政权材料被单一政权叙事遮蔽。 记录所见的结果则是：保留该年度的检索与补账位置；后续仅在获得可核单项材料时拆分为具体政治、财政、军事、社会或文化事件。 日期相邻不等于因果已经证实。

关于可说范围，须保留：本条只确认该年位于可追踪的共享年序中；正史、地方文书和外部材料的保存密度不一，不能由桥接标签推断行动、统属、疆界或社会影响。 进一步讨论可参照 蒙古征服、区域政权转换与跨语种年序研究

\subsection{1236年（元中统-23年；公元换算据《元史》本纪纪年）}
\noindent\eventanchor{SYD-1236-MONGOL-030}{八年丙申春正月，諸王各治具來會宴}\textbf{蒙古。}
蒙古 在 1236年（元中统-23年；公元换算据《元史》本纪纪年） 的材料痕迹是〔八年丙申春正月，諸王各治具來會宴〕，出处为 《元史》卷002〈本纪第2卷本年条〉。材料未将地点细化到城址，不能用中枢所在地替它补足。

把本项放回事件链，能够看到的前提为：本纪年条提供日期和中枢可见行动；前因不据单条补定。 而它所打开或收紧的下一步是：可回查同年及后续条目，地方影响仍须另证。

证据的边界是：《元史》为明初仓促官修，本纪保存大量实录条目但有编次与纪年问题；行省执行、数字、族称及对外关系须以条格、多语文书和对方史料互校。 蒙古征服、草原政治与跨区域文献比较研究 可用于继续检验这一结论。

\noindent\eventanchor{SYD-1236-MONGOL-080}{二月，命應州郭勝、鈞州孛朮魯九住、鄧州趙祥從曲出充先鋒伐宋}\textbf{蒙古与南宋。}
从 《元史》卷002〈本纪第2卷本年条〉 的 1236年（元中统-23年；公元换算据《元史》本纪纪年） 条目看，〔二月，命應州郭勝、鈞州孛朮魯九住、鄧州趙祥從曲出充先鋒伐宋〕使 蒙古与南宋 的一个具体行动进入叙事；题名保留的城、路、水道或机构名称，是目前可以落地的空间线索。

它承接的条件是：本纪年条记录政令或机构处置；实施背景须参照制度材料。 随后可追踪的变化为：可在同年及后续政令中追踪；条文不单独证明地方执行。

《元史》为明初仓促官修，本纪保存大量实录条目但有编次与纪年问题；行省执行、数字、族称及对外关系须以条格、多语文书和对方史料互校。 因此，蒙古征服、草原政治与跨区域文献比较研究 只能扩展问题，不能代替未保存的细节。

\noindent\eventanchor{SYD-1236-MONGOL-128}{皇子闊端、駙馬赤苦、公主阿剌海、公主果真、國王查剌溫、茶合帶、鍛真、蒙古寒札、按赤那顏、坼那顏、火斜、朮思，並於東平府戶內撥賜有差}\textbf{蒙古。}
1236年（元中统-23年；公元换算据《元史》本纪纪年） 的记载将 蒙古 与〔皇子闊端、駙馬赤苦、公主阿剌海、公主果真、國王查剌溫、茶合帶、鍛真、蒙古寒札、按赤那顏、坼那顏、火斜、朮思，並於東平府戶內撥賜有差〕相连。此处可确认的是这项被记录的行动；题名保留的城、路、水道或机构名称，是目前可以落地的空间线索。

前一层压力可由以下线索辨认：本纪年条提供日期和中枢可见行动；前因不据单条补定。 这一处置留下的后续条件是：可回查同年及后续条目，地方影响仍须另证。

关于可说范围，须保留：《元史》为明初仓促官修，本纪保存大量实录条目但有编次与纪年问题；行省执行、数字、族称及对外关系须以条格、多语文书和对方史料互校。 进一步讨论可参照 蒙古征服、草原政治与跨区域文献比较研究

\noindent\eventanchor{SYD-1236-MONGOL-170}{闊端率汪世顯等入蜀，取宋關外數州，斬蜀將曹友聞}\textbf{蒙古与南宋。}
〔闊端率汪世顯等入蜀，取宋關外數州，斬蜀將曹友聞〕见于 《元史》卷002〈本纪第2卷本年条〉 的 1236年（元中统-23年；公元换算据《元史》本纪纪年） 记录。蒙古与南宋 的选择因此有了可回查的时间位置；题名保留的城、路、水道或机构名称，是目前可以落地的空间线索。

本纪年条提供日期和中枢可见行动；前因不据单条补定。 记录所见的结果则是：可回查同年及后续条目，地方影响仍须另证。 日期相邻不等于因果已经证实。

证据的边界是：《元史》为明初仓促官修，本纪保存大量实录条目但有编次与纪年问题；行省执行、数字、族称及对外关系须以条格、多语文书和对方史料互校。 蒙古征服、草原政治与跨区域文献比较研究 可用于继续检验这一结论。

\noindent\eventanchor{SYD-1236-MONGOL-211}{秋七月，命陳時可閱刑名、科差、課稅等案，赴闕磨照}\textbf{蒙古。}
1236年（元中统-23年；公元换算据《元史》本纪纪年），《元史》卷002〈本纪第2卷本年条〉 所见的〔秋七月，命陳時可閱刑名、科差、課稅等案，赴闕磨照〕把 蒙古 的处置留在可复核的年序中；材料未将地点细化到城址，不能用中枢所在地替它补足。

把本项放回事件链，能够看到的前提为：本纪年条记录政令或机构处置；实施背景须参照制度材料。 而它所打开或收紧的下一步是：可在同年及后续政令中追踪；条文不单独证明地方执行。

《元史》为明初仓促官修，本纪保存大量实录条目但有编次与纪年问题；行省执行、数字、族称及对外关系须以条格、多语文书和对方史料互校。 因此，蒙古征服、草原政治与跨区域文献比较研究 只能扩展问题，不能代替未保存的细节。

\noindent\eventanchor{SYD-1236-MONGOL-247}{三月，復修孔子廟及司天臺}\textbf{蒙古。}
蒙古 在 1236年（元中统-23年；公元换算据《元史》本纪纪年） 的材料痕迹是〔三月，復修孔子廟及司天臺〕，出处为 《元史》卷002〈本纪第2卷本年条〉。材料未将地点细化到城址，不能用中枢所在地替它补足。

它承接的条件是：本纪年条记录政令或机构处置；实施背景须参照制度材料。 随后可追踪的变化为：可在同年及后续政令中追踪；条文不单独证明地方执行。

关于可说范围，须保留：《元史》为明初仓促官修，本纪保存大量实录条目但有编次与纪年问题；行省执行、数字、族称及对外关系须以条格、多语文书和对方史料互校。 进一步讨论可参照 蒙古征服、草原政治与跨区域文献比较研究

\noindent\eventanchor{SYD-1236-MONGOL-279}{夏六月，復括中州戶口，得續戶一百一十餘萬}\textbf{蒙古。}
从 《元史》卷002〈本纪第2卷本年条〉 的 1236年（元中统-23年；公元换算据《元史》本纪纪年） 条目看，〔夏六月，復括中州戶口，得續戶一百一十餘萬〕使 蒙古 的一个具体行动进入叙事；题名保留的城、路、水道或机构名称，是目前可以落地的空间线索。

前一层压力可由以下线索辨认：本纪年条记录政令或机构处置；实施背景须参照制度材料。 这一处置留下的后续条件是：可在同年及后续政令中追踪；条文不单独证明地方执行。

证据的边界是：《元史》为明初仓促官修，本纪保存大量实录条目但有编次与纪年问题；行省执行、数字、族称及对外关系须以条格、多语文书和对方史料互校。 蒙古征服、草原政治与跨区域文献比较研究 可用于继续检验这一结论。

\subsection{1237 年（共享年序桥接）}
\noindent\eventanchor{SY-1237-CHRONOLOGY-BRIDGE}{【C级年序桥接】保留1237 年的多政权并行检索入口；本条不主张所列政权在当年发生直接互动，具体事件须另以单项材料入账。}\textbf{南宋、金、西夏与蒙古（年序桥接）。}
1237 年（共享年序桥接） 的记载将 南宋、金、西夏与蒙古（年序桥接） 与〔【C级年序桥接】保留1237 年的多政权并行检索入口；本条不主张所列政权在当年发生直接互动，具体事件须另以单项材料入账。〕相连。此处可确认的是这项被记录的行动；题名保留的城、路、水道或机构名称，是目前可以落地的空间线索。

相邻已入账事件之间缺少该年度的共享年序入口，易使跨政权材料被单一政权叙事遮蔽。 记录所见的结果则是：保留该年度的检索与补账位置；后续仅在获得可核单项材料时拆分为具体政治、财政、军事、社会或文化事件。 日期相邻不等于因果已经证实。

本条只确认该年位于可追踪的共享年序中；正史、地方文书和外部材料的保存密度不一，不能由桥接标签推断行动、统属、疆界或社会影响。 因此，蒙古征服、区域政权转换与跨语种年序研究 只能扩展问题，不能代替未保存的细节。

\subsection{1237年（元中统-22年；公元换算据《元史》本纪纪年）}
\noindent\eventanchor{SYD-1237-MONGOL-031}{蒙哥征欽察部，破之，擒其酋八赤蠻}\textbf{蒙古。}
〔蒙哥征欽察部，破之，擒其酋八赤蠻〕见于 《元史》卷002〈本纪第2卷本年条〉 的 1237年（元中统-22年；公元换算据《元史》本纪纪年） 记录。蒙古 的选择因此有了可回查的时间位置；材料未将地点细化到城址，不能用中枢所在地替它补足。

把本项放回事件链，能够看到的前提为：本纪从朝廷视角记录军政行动；出兵与战果需分辨。 而它所打开或收紧的下一步是：后续路线、控制与损失应与对方记录和遗址材料复核。

关于可说范围，须保留：《元史》为明初仓促官修，本纪保存大量实录条目但有编次与纪年问题；行省执行、数字、族称及对外关系须以条格、多语文书和对方史料互校。 进一步讨论可参照 蒙古征服、草原政治与跨区域文献比较研究

\noindent\eventanchor{SYD-1237-MONGOL-081}{秋八月，命朮虎乃、劉中試諸路儒士，中選者除本貫議事官，得四千三十人}\textbf{蒙古。}
1237年（元中统-22年；公元换算据《元史》本纪纪年），《元史》卷002〈本纪第2卷本年条〉 所见的〔秋八月，命朮虎乃、劉中試諸路儒士，中選者除本貫議事官，得四千三十人〕把 蒙古 的处置留在可复核的年序中；题名保留的城、路、水道或机构名称，是目前可以落地的空间线索。

它承接的条件是：本纪年条记录政令或机构处置；实施背景须参照制度材料。 随后可追踪的变化为：可在同年及后续政令中追踪；条文不单独证明地方执行。

证据的边界是：《元史》为明初仓促官修，本纪保存大量实录条目但有编次与纪年问题；行省执行、数字、族称及对外关系须以条格、多语文书和对方史料互校。 蒙古征服、草原政治与跨区域文献比较研究 可用于继续检验这一结论。

\noindent\eventanchor{SYD-1237-MONGOL-129}{是冬，口溫不花等圍光州，命張柔、鞏彥暉、史天澤攻下之}\textbf{蒙古。}
蒙古 在 1237年（元中统-22年；公元换算据《元史》本纪纪年） 的材料痕迹是〔是冬，口溫不花等圍光州，命張柔、鞏彥暉、史天澤攻下之〕，出处为 《元史》卷002〈本纪第2卷本年条〉。题名保留的城、路、水道或机构名称，是目前可以落地的空间线索。

前一层压力可由以下线索辨认：本纪年条记录政令或机构处置；实施背景须参照制度材料。 这一处置留下的后续条件是：可在同年及后续政令中追踪；条文不单独证明地方执行。

《元史》为明初仓促官修，本纪保存大量实录条目但有编次与纪年问题；行省执行、数字、族称及对外关系须以条格、多语文书和对方史料互校。 因此，蒙古征服、草原政治与跨区域文献比较研究 只能扩展问题，不能代替未保存的细节。

\noindent\eventanchor{SYD-1237-MONGOL-171}{遂別攻蘄州，降隨州，略地至黃州}\textbf{蒙古。}
从 《元史》卷002〈本纪第2卷本年条〉 的 1237年（元中统-22年；公元换算据《元史》本纪纪年） 条目看，〔遂別攻蘄州，降隨州，略地至黃州〕使 蒙古 的一个具体行动进入叙事；题名保留的城、路、水道或机构名称，是目前可以落地的空间线索。

本纪从朝廷视角记录军政行动；出兵与战果需分辨。 记录所见的结果则是：后续路线、控制与损失应与对方记录和遗址材料复核。 日期相邻不等于因果已经证实。

关于可说范围，须保留：《元史》为明初仓促官修，本纪保存大量实录条目但有编次与纪年问题；行省执行、数字、族称及对外关系须以条格、多语文书和对方史料互校。 进一步讨论可参照 蒙古征服、草原政治与跨区域文献比较研究

\noindent\eventanchor{SYD-1237-MONGOL-212}{夏四月，築掃隣城，作迦堅茶寒殿}\textbf{蒙古。}
1237年（元中统-22年；公元换算据《元史》本纪纪年） 的记载将 蒙古 与〔夏四月，築掃隣城，作迦堅茶寒殿〕相连。此处可确认的是这项被记录的行动；题名保留的城、路、水道或机构名称，是目前可以落地的空间线索。

把本项放回事件链，能够看到的前提为：本纪年条提供日期和中枢可见行动；前因不据单条补定。 而它所打开或收紧的下一步是：可回查同年及后续条目，地方影响仍须另证。

证据的边界是：《元史》为明初仓促官修，本纪保存大量实录条目但有编次与纪年问题；行省执行、数字、族称及对外关系须以条格、多语文书和对方史料互校。 蒙古征服、草原政治与跨区域文献比较研究 可用于继续检验这一结论。

\subsection{1238 年（共享年序桥接）}
\noindent\eventanchor{SY-1238-CHRONOLOGY-BRIDGE}{【C级年序桥接】保留1238 年的多政权并行检索入口；本条不主张所列政权在当年发生直接互动，具体事件须另以单项材料入账。}\textbf{南宋、金、西夏与蒙古（年序桥接）。}
〔【C级年序桥接】保留1238 年的多政权并行检索入口；本条不主张所列政权在当年发生直接互动，具体事件须另以单项材料入账。〕见于 《宋史》；《金史》；《元史》；蒙古文、波斯文及地方材料（据事核） 的 1238 年（共享年序桥接） 记录。南宋、金、西夏与蒙古（年序桥接） 的选择因此有了可回查的时间位置；题名保留的城、路、水道或机构名称，是目前可以落地的空间线索。

它承接的条件是：相邻已入账事件之间缺少该年度的共享年序入口，易使跨政权材料被单一政权叙事遮蔽。 随后可追踪的变化为：保留该年度的检索与补账位置；后续仅在获得可核单项材料时拆分为具体政治、财政、军事、社会或文化事件。

本条只确认该年位于可追踪的共享年序中；正史、地方文书和外部材料的保存密度不一，不能由桥接标签推断行动、统属、疆界或社会影响。 因此，蒙古征服、区域政权转换与跨语种年序研究 只能扩展问题，不能代替未保存的细节。

\subsection{1238年（元中统-21年；公元换算据《元史》本纪纪年）}
\noindent\eventanchor{SYD-1238-MONGOL-032}{秋八月，陳時可、高慶民等言諸路旱蝗，詔免今年田租，仍停舊未輸納者，俟豐歲議之}\textbf{蒙古。}
1238年（元中统-21年；公元换算据《元史》本纪纪年），《元史》卷002〈本纪第2卷本年条〉 所见的〔秋八月，陳時可、高慶民等言諸路旱蝗，詔免今年田租，仍停舊未輸納者，俟豐歲議之〕把 蒙古 的处置留在可复核的年序中；题名保留的城、路、水道或机构名称，是目前可以落地的空间线索。

前一层压力可由以下线索辨认：本纪年条记录政令或机构处置；实施背景须参照制度材料。 这一处置留下的后续条件是：可在同年及后续政令中追踪；条文不单独证明地方执行。

关于可说范围，须保留：《元史》为明初仓促官修，本纪保存大量实录条目但有编次与纪年问题；行省执行、数字、族称及对外关系须以条格、多语文书和对方史料互校。 进一步讨论可参照 蒙古征服、草原政治与跨区域文献比较研究

\noindent\eventanchor{SYD-1238-MONGOL-082}{十年戊戌春，塔思軍至北峽關，宋將汪統制降}\textbf{蒙古与南宋。}
蒙古与南宋 在 1238年（元中统-21年；公元换算据《元史》本纪纪年） 的材料痕迹是〔十年戊戌春，塔思軍至北峽關，宋將汪統制降〕，出处为 《元史》卷002〈本纪第2卷本年条〉。题名保留的城、路、水道或机构名称，是目前可以落地的空间线索。

本纪从朝廷视角记录军政行动；出兵与战果需分辨。 记录所见的结果则是：后续路线、控制与损失应与对方记录和遗址材料复核。 日期相邻不等于因果已经证实。

证据的边界是：《元史》为明初仓促官修，本纪保存大量实录条目但有编次与纪年问题；行省执行、数字、族称及对外关系须以条格、多语文书和对方史料互校。 蒙古征服、草原政治与跨区域文献比较研究 可用于继续检验这一结论。

\noindent\eventanchor{SYD-1238-MONGOL-130}{夏，襄陽別將劉義叛，執游顯等降宋}\textbf{蒙古与南宋。}
从 《元史》卷002〈本纪第2卷本年条〉 的 1238年（元中统-21年；公元换算据《元史》本纪纪年） 条目看，〔夏，襄陽別將劉義叛，執游顯等降宋〕使 蒙古与南宋 的一个具体行动进入叙事；材料未将地点细化到城址，不能用中枢所在地替它补足。

把本项放回事件链，能够看到的前提为：本纪从朝廷视角记录军政行动；出兵与战果需分辨。 而它所打开或收紧的下一步是：后续路线、控制与损失应与对方记录和遗址材料复核。

《元史》为明初仓促官修，本纪保存大量实录条目但有编次与纪年问题；行省执行、数字、族称及对外关系须以条格、多语文书和对方史料互校。 因此，蒙古征服、草原政治与跨区域文献比较研究 只能扩展问题，不能代替未保存的细节。

\noindent\eventanchor{SYD-1238-MONGOL-172}{築圖蘇湖城，作迎駕殿}\textbf{蒙古。}
1238年（元中统-21年；公元换算据《元史》本纪纪年） 的记载将 蒙古 与〔築圖蘇湖城，作迎駕殿〕相连。此处可确认的是这项被记录的行动；题名保留的城、路、水道或机构名称，是目前可以落地的空间线索。

它承接的条件是：本纪年条提供日期和中枢可见行动；前因不据单条补定。 随后可追踪的变化为：可回查同年及后续条目，地方影响仍须另证。

关于可说范围，须保留：《元史》为明初仓促官修，本纪保存大量实录条目但有编次与纪年问题；行省执行、数字、族称及对外关系须以条格、多语文书和对方史料互校。 进一步讨论可参照 蒙古征服、草原政治与跨区域文献比较研究

\subsection{1239 年（共享年序桥接）}
\noindent\eventanchor{SY-1239-CHRONOLOGY-BRIDGE}{【C级年序桥接】保留1239 年的多政权并行检索入口；本条不主张所列政权在当年发生直接互动，具体事件须另以单项材料入账。}\textbf{南宋、金、西夏与蒙古（年序桥接）。}
〔【C级年序桥接】保留1239 年的多政权并行检索入口；本条不主张所列政权在当年发生直接互动，具体事件须另以单项材料入账。〕见于 《宋史》；《金史》；《元史》；蒙古文、波斯文及地方材料（据事核） 的 1239 年（共享年序桥接） 记录。南宋、金、西夏与蒙古（年序桥接） 的选择因此有了可回查的时间位置；题名保留的城、路、水道或机构名称，是目前可以落地的空间线索。

前一层压力可由以下线索辨认：相邻已入账事件之间缺少该年度的共享年序入口，易使跨政权材料被单一政权叙事遮蔽。 这一处置留下的后续条件是：保留该年度的检索与补账位置；后续仅在获得可核单项材料时拆分为具体政治、财政、军事、社会或文化事件。

证据的边界是：本条只确认该年位于可追踪的共享年序中；正史、地方文书和外部材料的保存密度不一，不能由桥接标签推断行动、统属、疆界或社会影响。 蒙古征服、区域政权转换与跨语种年序研究 可用于继续检验这一结论。

\subsection{1239年（元中统-20年；公元换算据《元史》本纪纪年）}
\noindent\eventanchor{SYD-1239-MONGOL-033}{十一年己亥春，復獵于揭揭察哈之澤}\textbf{蒙古。}
1239年（元中统-20年；公元换算据《元史》本纪纪年），《元史》卷002〈本纪第2卷本年条〉 所见的〔十一年己亥春，復獵于揭揭察哈之澤〕把 蒙古 的处置留在可复核的年序中；材料未将地点细化到城址，不能用中枢所在地替它补足。

本纪年条记录政令或机构处置；实施背景须参照制度材料。 记录所见的结果则是：可在同年及后续政令中追踪；条文不单独证明地方执行。 日期相邻不等于因果已经证实。

《元史》为明初仓促官修，本纪保存大量实录条目但有编次与纪年问题；行省执行、数字、族称及对外关系须以条格、多语文书和对方史料互校。 因此，蒙古征服、草原政治与跨区域文献比较研究 只能扩展问题，不能代替未保存的细节。

\noindent\eventanchor{SYD-1239-MONGOL-083}{以山東諸路災，免其稅糧}\textbf{蒙古。}
蒙古 在 1239年（元中统-20年；公元换算据《元史》本纪纪年） 的材料痕迹是〔以山東諸路災，免其稅糧〕，出处为 《元史》卷002〈本纪第2卷本年条〉。题名保留的城、路、水道或机构名称，是目前可以落地的空间线索。

把本项放回事件链，能够看到的前提为：本纪所记财用、赈济或征发属于特定报告场景。 而它所打开或收紧的下一步是：可与食货志、文书和物质材料比较，不外推为全境总量。

关于可说范围，须保留：《元史》为明初仓促官修，本纪保存大量实录条目但有编次与纪年问题；行省执行、数字、族称及对外关系须以条格、多语文书和对方史料互校。 进一步讨论可参照 蒙古征服、草原政治与跨区域文献比较研究

\subsection{1240 年（共享年序桥接）}
\noindent\eventanchor{SY-1240-CHRONOLOGY-BRIDGE}{【C级年序桥接】保留1240 年的多政权并行检索入口；本条不主张所列政权在当年发生直接互动，具体事件须另以单项材料入账。}\textbf{南宋、金、西夏与蒙古（年序桥接）。}
从 《宋史》；《金史》；《元史》；蒙古文、波斯文及地方材料（据事核） 的 1240 年（共享年序桥接） 条目看，〔【C级年序桥接】保留1240 年的多政权并行检索入口；本条不主张所列政权在当年发生直接互动，具体事件须另以单项材料入账。〕使 南宋、金、西夏与蒙古（年序桥接） 的一个具体行动进入叙事；题名保留的城、路、水道或机构名称，是目前可以落地的空间线索。

它承接的条件是：相邻已入账事件之间缺少该年度的共享年序入口，易使跨政权材料被单一政权叙事遮蔽。 随后可追踪的变化为：保留该年度的检索与补账位置；后续仅在获得可核单项材料时拆分为具体政治、财政、军事、社会或文化事件。

证据的边界是：本条只确认该年位于可追踪的共享年序中；正史、地方文书和外部材料的保存密度不一，不能由桥接标签推断行动、统属、疆界或社会影响。 蒙古征服、区域政权转换与跨语种年序研究 可用于继续检验这一结论。

\subsection{1240年（元中统-19年；公元换算据《元史》本纪纪年）}
\noindent\eventanchor{SYD-1240-MONGOL-034}{冬十二月，詔貴由班師}\textbf{蒙古。}
1240年（元中统-19年；公元换算据《元史》本纪纪年） 的记载将 蒙古 与〔冬十二月，詔貴由班師〕相连。此处可确认的是这项被记录的行动；材料未将地点细化到城址，不能用中枢所在地替它补足。

前一层压力可由以下线索辨认：本纪年条记录政令或机构处置；实施背景须参照制度材料。 这一处置留下的后续条件是：可在同年及后续政令中追踪；条文不单独证明地方执行。

《元史》为明初仓促官修，本纪保存大量实录条目但有编次与纪年问题；行省执行、数字、族称及对外关系须以条格、多语文书和对方史料互校。 因此，蒙古征服、草原政治与跨区域文献比较研究 只能扩展问题，不能代替未保存的细节。

\noindent\eventanchor{SYD-1240-MONGOL-084}{國初，令民代償，民多亡命，至是罷之}\textbf{蒙古。}
〔國初，令民代償，民多亡命，至是罷之〕见于 《元史》卷002〈本纪第2卷本年条〉 的 1240年（元中统-19年；公元换算据《元史》本纪纪年） 记录。蒙古 的选择因此有了可回查的时间位置；材料未将地点细化到城址，不能用中枢所在地替它补足。

本纪年条记录政令或机构处置；实施背景须参照制度材料。 记录所见的结果则是：可在同年及后续政令中追踪；条文不单独证明地方执行。 日期相邻不等于因果已经证实。

关于可说范围，须保留：《元史》为明初仓促官修，本纪保存大量实录条目但有编次与纪年问题；行省执行、数字、族称及对外关系须以条格、多语文书和对方史料互校。 进一步讨论可参照 蒙古征服、草原政治与跨区域文献比较研究

\noindent\eventanchor{SYD-1240-MONGOL-131}{皇子貴由克西域未下諸部，遣使奏捷}\textbf{蒙古。}
1240年（元中统-19年；公元换算据《元史》本纪纪年），《元史》卷002〈本纪第2卷本年条〉 所见的〔皇子貴由克西域未下諸部，遣使奏捷〕把 蒙古 的处置留在可复核的年序中；材料未将地点细化到城址，不能用中枢所在地替它补足。

把本项放回事件链，能够看到的前提为：本纪记录使节、贡赐或往来，不等于稳定统属关系。 而它所打开或收紧的下一步是：可与对方编年和交通、文书材料核对其后续落实。

证据的边界是：《元史》为明初仓促官修，本纪保存大量实录条目但有编次与纪年问题；行省执行、数字、族称及对外关系须以条格、多语文书和对方史料互校。 蒙古征服、草原政治与跨区域文献比较研究 可用于继续检验这一结论。

\noindent\eventanchor{SYD-1240-MONGOL-173}{仍命凡假貸歲久，惟子本相侔而止，著為令}\textbf{蒙古。}
蒙古 在 1240年（元中统-19年；公元换算据《元史》本纪纪年） 的材料痕迹是〔仍命凡假貸歲久，惟子本相侔而止，著為令〕，出处为 《元史》卷002〈本纪第2卷本年条〉。材料未将地点细化到城址，不能用中枢所在地替它补足。

它承接的条件是：本纪年条记录政令或机构处置；实施背景须参照制度材料。 随后可追踪的变化为：可在同年及后续政令中追踪；条文不单独证明地方执行。

《元史》为明初仓促官修，本纪保存大量实录条目但有编次与纪年问题；行省执行、数字、族称及对外关系须以条格、多语文书和对方史料互校。 因此，蒙古征服、草原政治与跨区域文献比较研究 只能扩展问题，不能代替未保存的细节。

\noindent\eventanchor{SYD-1240-MONGOL-213}{是歲，以官民貸回鶻金償官者歲加倍，名羊羔息，其害為甚，詔以官物代還，凡七萬六千錠}\textbf{蒙古。}
从 《元史》卷002〈本纪第2卷本年条〉 的 1240年（元中统-19年；公元换算据《元史》本纪纪年） 条目看，〔是歲，以官民貸回鶻金償官者歲加倍，名羊羔息，其害為甚，詔以官物代還，凡七萬六千錠〕使 蒙古 的一个具体行动进入叙事；材料未将地点细化到城址，不能用中枢所在地替它补足。

前一层压力可由以下线索辨认：本纪年条记录政令或机构处置；实施背景须参照制度材料。 这一处置留下的后续条件是：可在同年及后续政令中追踪；条文不单独证明地方执行。

关于可说范围，须保留：《元史》为明初仓促官修，本纪保存大量实录条目但有编次与纪年问题；行省执行、数字、族称及对外关系须以条格、多语文书和对方史料互校。 进一步讨论可参照 蒙古征服、草原政治与跨区域文献比较研究

\subsection{1241年（元中统-18年；公元换算据《元史》本纪纪年）}
\noindent\eventanchor{SYD-1241-MONGOL-035}{奧都剌合蠻進酒，帝歡飲，極夜乃罷}\textbf{蒙古。}
1241年（元中统-18年；公元换算据《元史》本纪纪年） 的记载将 蒙古 与〔奧都剌合蠻進酒，帝歡飲，極夜乃罷〕相连。此处可确认的是这项被记录的行动；题名保留的城、路、水道或机构名称，是目前可以落地的空间线索。

本纪年条记录政令或机构处置；实施背景须参照制度材料。 记录所见的结果则是：可在同年及后续政令中追踪；条文不单独证明地方执行。 日期相邻不等于因果已经证实。

证据的边界是：《元史》为明初仓促官修，本纪保存大量实录条目但有编次与纪年问题；行省执行、数字、族称及对外关系须以条格、多语文书和对方史料互校。 蒙古征服、草原政治与跨区域文献比较研究 可用于继续检验这一结论。

\noindent\eventanchor{SYD-1241-MONGOL-085}{帝有疾，詔赦天下囚徒}\textbf{蒙古。}
〔帝有疾，詔赦天下囚徒〕见于 《元史》卷002〈本纪第2卷本年条〉 的 1241年（元中统-18年；公元换算据《元史》本纪纪年） 记录。蒙古 的选择因此有了可回查的时间位置；材料未将地点细化到城址，不能用中枢所在地替它补足。

把本项放回事件链，能够看到的前提为：本纪年条记录政令或机构处置；实施背景须参照制度材料。 而它所打开或收紧的下一步是：可在同年及后续政令中追踪；条文不单独证明地方执行。

《元史》为明初仓促官修，本纪保存大量实录条目但有编次与纪年问题；行省执行、数字、族称及对外关系须以条格、多语文书和对方史料互校。 因此，蒙古征服、草原政治与跨区域文献比较研究 只能扩展问题，不能代替未保存的细节。

\noindent\eventanchor{SYD-1241-MONGOL-132}{冬十月，命牙老瓦赤主管漢民公事}\textbf{蒙古。}
1241年（元中统-18年；公元换算据《元史》本纪纪年），《元史》卷002〈本纪第2卷本年条〉 所见的〔冬十月，命牙老瓦赤主管漢民公事〕把 蒙古 的处置留在可复核的年序中；材料未将地点细化到城址，不能用中枢所在地替它补足。

它承接的条件是：本纪年条记录政令或机构处置；实施背景须参照制度材料。 随后可追踪的变化为：可在同年及后续政令中追踪；条文不单独证明地方执行。

关于可说范围，须保留：《元史》为明初仓促官修，本纪保存大量实录条目但有编次与纪年问题；行省执行、数字、族称及对外关系须以条格、多语文书和对方史料互校。 进一步讨论可参照 蒙古征服、草原政治与跨区域文献比较研究

\noindent\eventanchor{SYD-1241-MONGOL-174}{秋，高麗國王㬚以族子綧入質}\textbf{蒙古。}
蒙古 在 1241年（元中统-18年；公元换算据《元史》本纪纪年） 的材料痕迹是〔秋，高麗國王㬚以族子綧入質〕，出处为 《元史》卷002〈本纪第2卷本年条〉。材料未将地点细化到城址，不能用中枢所在地替它补足。

前一层压力可由以下线索辨认：本纪年条提供日期和中枢可见行动；前因不据单条补定。 这一处置留下的后续条件是：可回查同年及后续条目，地方影响仍须另证。

证据的边界是：《元史》为明初仓促官修，本纪保存大量实录条目但有编次与纪年问题；行省执行、数字、族称及对外关系须以条格、多语文书和对方史料互校。 蒙古征服、草原政治与跨区域文献比较研究 可用于继续检验这一结论。

\noindent\eventanchor{SYD-1241-MONGOL-214}{秋，后命張柔總兵戍杞}\textbf{蒙古。}
从 《元史》卷002〈本纪第2卷本年条〉 的 1241年（元中统-18年；公元换算据《元史》本纪纪年） 条目看，〔秋，后命張柔總兵戍杞〕使 蒙古 的一个具体行动进入叙事；材料未将地点细化到城址，不能用中枢所在地替它补足。

本纪年条记录政令或机构处置；实施背景须参照制度材料。 记录所见的结果则是：可在同年及后续政令中追踪；条文不单独证明地方执行。 日期相邻不等于因果已经证实。

《元史》为明初仓促官修，本纪保存大量实录条目但有编次与纪年问题；行省执行、数字、族称及对外关系须以条格、多语文书和对方史料互校。 因此，蒙古征服、草原政治与跨区域文献比较研究 只能扩展问题，不能代替未保存的细节。

\noindent\eventanchor{SYD-1241-MONGOL-248}{秋七月，張柔自五河口渡淮，攻宋揚、滁、和等州}\textbf{蒙古与南宋。}
1241年（元中统-18年；公元换算据《元史》本纪纪年） 的记载将 蒙古与南宋 与〔秋七月，張柔自五河口渡淮，攻宋揚、滁、和等州〕相连。此处可确认的是这项被记录的行动；题名保留的城、路、水道或机构名称，是目前可以落地的空间线索。

把本项放回事件链，能够看到的前提为：本纪保留灾害或工程的报告地点与朝廷处置；灾情范围需另证。 而它所打开或收紧的下一步是：可与地方文书、河工和环境材料比较，不将单点记录外推为全域因果。

关于可说范围，须保留：《元史》为明初仓促官修，本纪保存大量实录条目但有编次与纪年问题；行省执行、数字、族称及对外关系须以条格、多语文书和对方史料互校。 进一步讨论可参照 蒙古征服、草原政治与跨区域文献比较研究

\noindent\eventanchor{SYD-1241-MONGOL-280}{宋制置趙蔡請和，乃還}\textbf{蒙古与南宋。}
〔宋制置趙蔡請和，乃還〕见于 《元史》卷002〈本纪第2卷本年条〉 的 1241年（元中统-18年；公元换算据《元史》本纪纪年） 记录。蒙古与南宋 的选择因此有了可回查的时间位置；材料未将地点细化到城址，不能用中枢所在地替它补足。

它承接的条件是：本纪年条记录政令或机构处置；实施背景须参照制度材料。 随后可追踪的变化为：可在同年及后续政令中追踪；条文不单独证明地方执行。

证据的边界是：《元史》为明初仓促官修，本纪保存大量实录条目但有编次与纪年问题；行省执行、数字、族称及对外关系须以条格、多语文书和对方史料互校。 蒙古征服、草原政治与跨区域文献比较研究 可用于继续检验这一结论。

\subsection{1241年}
\noindent\eventanchor{YUAN-1241-EVENT-162}{窝阔台去世，帝国进入摄政与继承协商阶段}\textbf{蒙古。}
1241年，《元史》相关本纪；《元朝秘史》、波斯文史籍或《高丽史》对应叙事；《元典章》《通制条格》、多语碑刻与文书（据事核） 所见的〔窝阔台去世，帝国进入摄政与继承协商阶段〕把 蒙古 的处置留在可复核的年序中；材料未将地点细化到城址，不能用中枢所在地替它补足。

前一层压力可由以下线索辨认：窝阔台至蒙哥时期的扩张与继承中前任统治的结束与宫廷、部族或官僚的权力安排相互交织，促成“窝阔台去世，帝国进入摄政与继承协商阶段”这一继承节点。 这一处置留下的后续条件是：继承并不自动消除既有矛盾；“窝阔台去世，帝国进入摄政与继承协商阶段”后的任官、军权和对外策略需由后续条目追踪。

《元史》明初速修，纪年与人名有异同；蒙古大汗政权不等同所有汗国，征服、赋役和身份分类须按地区及材料语言分列。 因此，蒙古帝国、元朝行省财政、多语文书与区域社会研究 只能扩展问题，不能代替未保存的细节。

\subsection{1242 年（共享年序桥接）}
\noindent\eventanchor{SY-1242-CHRONOLOGY-BRIDGE}{【C级年序桥接】保留1242 年的多政权并行检索入口；本条不主张所列政权在当年发生直接互动，具体事件须另以单项材料入账。}\textbf{南宋、金、西夏与蒙古（年序桥接）。}
南宋、金、西夏与蒙古（年序桥接） 在 1242 年（共享年序桥接） 的材料痕迹是〔【C级年序桥接】保留1242 年的多政权并行检索入口；本条不主张所列政权在当年发生直接互动，具体事件须另以单项材料入账。〕，出处为 《宋史》；《金史》；《元史》；蒙古文、波斯文及地方材料（据事核）。题名保留的城、路、水道或机构名称，是目前可以落地的空间线索。

相邻已入账事件之间缺少该年度的共享年序入口，易使跨政权材料被单一政权叙事遮蔽。 记录所见的结果则是：保留该年度的检索与补账位置；后续仅在获得可核单项材料时拆分为具体政治、财政、军事、社会或文化事件。 日期相邻不等于因果已经证实。

关于可说范围，须保留：本条只确认该年位于可追踪的共享年序中；正史、地方文书和外部材料的保存密度不一，不能由桥接标签推断行动、统属、疆界或社会影响。 进一步讨论可参照 蒙古征服、区域政权转换与跨语种年序研究

\subsection{1243 年（共享年序桥接）}
\noindent\eventanchor{SY-1243-CHRONOLOGY-BRIDGE}{【C级年序桥接】保留1243 年的多政权并行检索入口；本条不主张所列政权在当年发生直接互动，具体事件须另以单项材料入账。}\textbf{南宋、金、西夏与蒙古（年序桥接）。}
从 《宋史》；《金史》；《元史》；蒙古文、波斯文及地方材料（据事核） 的 1243 年（共享年序桥接） 条目看，〔【C级年序桥接】保留1243 年的多政权并行检索入口；本条不主张所列政权在当年发生直接互动，具体事件须另以单项材料入账。〕使 南宋、金、西夏与蒙古（年序桥接） 的一个具体行动进入叙事；题名保留的城、路、水道或机构名称，是目前可以落地的空间线索。

把本项放回事件链，能够看到的前提为：相邻已入账事件之间缺少该年度的共享年序入口，易使跨政权材料被单一政权叙事遮蔽。 而它所打开或收紧的下一步是：保留该年度的检索与补账位置；后续仅在获得可核单项材料时拆分为具体政治、财政、军事、社会或文化事件。

证据的边界是：本条只确认该年位于可追踪的共享年序中；正史、地方文书和外部材料的保存密度不一，不能由桥接标签推断行动、统属、疆界或社会影响。 蒙古征服、区域政权转换与跨语种年序研究 可用于继续检验这一结论。

\subsection{1244 年（共享年序桥接）}
\noindent\eventanchor{SY-1244-CHRONOLOGY-BRIDGE}{【C级年序桥接】保留1244 年的多政权并行检索入口；本条不主张所列政权在当年发生直接互动，具体事件须另以单项材料入账。}\textbf{南宋、金、西夏与蒙古（年序桥接）。}
1244 年（共享年序桥接） 的记载将 南宋、金、西夏与蒙古（年序桥接） 与〔【C级年序桥接】保留1244 年的多政权并行检索入口；本条不主张所列政权在当年发生直接互动，具体事件须另以单项材料入账。〕相连。此处可确认的是这项被记录的行动；题名保留的城、路、水道或机构名称，是目前可以落地的空间线索。

它承接的条件是：相邻已入账事件之间缺少该年度的共享年序入口，易使跨政权材料被单一政权叙事遮蔽。 随后可追踪的变化为：保留该年度的检索与补账位置；后续仅在获得可核单项材料时拆分为具体政治、财政、军事、社会或文化事件。

本条只确认该年位于可追踪的共享年序中；正史、地方文书和外部材料的保存密度不一，不能由桥接标签推断行动、统属、疆界或社会影响。 因此，蒙古征服、区域政权转换与跨语种年序研究 只能扩展问题，不能代替未保存的细节。

\subsection{1245 年（共享年序桥接）}
\noindent\eventanchor{SY-1245-CHRONOLOGY-BRIDGE}{【C级年序桥接】保留1245 年的多政权并行检索入口；本条不主张所列政权在当年发生直接互动，具体事件须另以单项材料入账。}\textbf{南宋、金、西夏与蒙古（年序桥接）。}
〔【C级年序桥接】保留1245 年的多政权并行检索入口；本条不主张所列政权在当年发生直接互动，具体事件须另以单项材料入账。〕见于 《宋史》；《金史》；《元史》；蒙古文、波斯文及地方材料（据事核） 的 1245 年（共享年序桥接） 记录。南宋、金、西夏与蒙古（年序桥接） 的选择因此有了可回查的时间位置；题名保留的城、路、水道或机构名称，是目前可以落地的空间线索。

前一层压力可由以下线索辨认：相邻已入账事件之间缺少该年度的共享年序入口，易使跨政权材料被单一政权叙事遮蔽。 这一处置留下的后续条件是：保留该年度的检索与补账位置；后续仅在获得可核单项材料时拆分为具体政治、财政、军事、社会或文化事件。

关于可说范围，须保留：本条只确认该年位于可追踪的共享年序中；正史、地方文书和外部材料的保存密度不一，不能由桥接标签推断行动、统属、疆界或社会影响。 进一步讨论可参照 蒙古征服、区域政权转换与跨语种年序研究

\subsection{1246年（元中统-13年；公元换算据《元史》本纪纪年）}
\noindent\eventanchor{SYD-1246-MONGOL-036}{帝聞之，即班師，而水已至，後軍有浮渡者}\textbf{蒙古。}
1246年（元中统-13年；公元换算据《元史》本纪纪年），《元史》卷003〈本纪第3卷本年条〉 所见的〔帝聞之，即班師，而水已至，後軍有浮渡者〕把 蒙古 的处置留在可复核的年序中；题名保留的城、路、水道或机构名称，是目前可以落地的空间线索。

本纪年条记录政令或机构处置；实施背景须参照制度材料。 记录所见的结果则是：可在同年及后续政令中追踪；条文不单独证明地方执行。 日期相邻不等于因果已经证实。

证据的边界是：《元史》为明初仓促官修，本纪保存大量实录条目但有编次与纪年问题；行省执行、数字、族称及对外关系须以条格、多语文书和对方史料互校。 蒙古征服、草原政治与跨区域文献比较研究 可用于继续检验这一结论。

\noindent\eventanchor{SYD-1246-MONGOL-086}{八赤蠻謂守者曰：「我之竄入于海，與魚何異？然終見擒，天也}\textbf{蒙古。}
蒙古 在 1246年（元中统-13年；公元换算据《元史》本纪纪年） 的材料痕迹是〔八赤蠻謂守者曰：「我之竄入于海，與魚何異？然終見擒，天也〕，出处为 《元史》卷003〈本纪第3卷本年条〉。题名保留的城、路、水道或机构名称，是目前可以落地的空间线索。

把本项放回事件链，能够看到的前提为：本纪年条提供日期和中枢可见行动；前因不据单条补定。 而它所打开或收紧的下一步是：可回查同年及后续条目，地方影响仍须另证。

《元史》为明初仓促官修，本纪保存大量实录条目但有编次与纪年问题；行省执行、数字、族称及对外关系须以条格、多语文书和对方史料互校。 因此，蒙古征服、草原政治与跨区域文献比较研究 只能扩展问题，不能代替未保存的细节。

\noindent\eventanchor{SYD-1246-MONGOL-133}{八赤蠻曰：「我為一國主，豈苟求生？且身非駝，何以跪人為？」乃命囚之}\textbf{蒙古。}
从 《元史》卷003〈本纪第3卷本年条〉 的 1246年（元中统-13年；公元换算据《元史》本纪纪年） 条目看，〔八赤蠻曰：「我為一國主，豈苟求生？且身非駝，何以跪人為？」乃命囚之〕使 蒙古 的一个具体行动进入叙事；材料未将地点细化到城址，不能用中枢所在地替它补足。

它承接的条件是：本纪年条记录政令或机构处置；实施背景须参照制度材料。 随后可追踪的变化为：可在同年及后续政令中追踪；条文不单独证明地方执行。

关于可说范围，须保留：《元史》为明初仓促官修，本纪保存大量实录条目但有编次与纪年问题；行省执行、数字、族称及对外关系须以条格、多语文书和对方史料互校。 进一步讨论可参照 蒙古征服、草原政治与跨区域文献比较研究

\noindent\eventanchor{SYD-1246-MONGOL-175}{罷築和林城役千五百人}\textbf{蒙古。}
1246年（元中统-13年；公元换算据《元史》本纪纪年） 的记载将 蒙古 与〔罷築和林城役千五百人〕相连。此处可确认的是这项被记录的行动；题名保留的城、路、水道或机构名称，是目前可以落地的空间线索。

前一层压力可由以下线索辨认：本纪年条记录政令或机构处置；实施背景须参照制度材料。 这一处置留下的后续条件是：可在同年及后续政令中追踪；条文不单独证明地方执行。

证据的边界是：《元史》为明初仓促官修，本纪保存大量实录条目但有编次与纪年问题；行省执行、数字、族称及对外关系须以条格、多语文书和对方史料互校。 蒙古征服、草原政治与跨区域文献比较研究 可用于继续检验这一结论。

\noindent\eventanchor{SYD-1246-MONGOL-215}{嘗攻欽察部，其酋八赤蠻逃于海島}\textbf{蒙古。}
〔嘗攻欽察部，其酋八赤蠻逃于海島〕见于 《元史》卷003〈本纪第3卷本年条〉 的 1246年（元中统-13年；公元换算据《元史》本纪纪年） 记录。蒙古 的选择因此有了可回查的时间位置；题名保留的城、路、水道或机构名称，是目前可以落地的空间线索。

本纪从朝廷视角记录军政行动；出兵与战果需分辨。 记录所见的结果则是：后续路线、控制与损失应与对方记录和遗址材料复核。 日期相邻不等于因果已经证实。

《元史》为明初仓促官修，本纪保存大量实录条目但有编次与纪年问题；行省执行、数字、族称及对外关系须以条格、多语文书和对方史料互校。 因此，蒙古征服、草原政治与跨区域文献比较研究 只能扩展问题，不能代替未保存的细节。

\noindent\eventanchor{SYD-1246-MONGOL-249}{帝遣諸王旭烈與忙可撒兒帥兵覘之}\textbf{蒙古。}
1246年（元中统-13年；公元换算据《元史》本纪纪年），《元史》卷003〈本纪第3卷本年条〉 所见的〔帝遣諸王旭烈與忙可撒兒帥兵覘之〕把 蒙古 的处置留在可复核的年序中；材料未将地点细化到城址，不能用中枢所在地替它补足。

把本项放回事件链，能够看到的前提为：本纪年条提供日期和中枢可见行动；前因不据单条补定。 而它所打开或收紧的下一步是：可回查同年及后续条目，地方影响仍须另证。

关于可说范围，须保留：《元史》为明初仓促官修，本纪保存大量实录条目但有编次与纪年问题；行省执行、数字、族称及对外关系须以条格、多语文书和对方史料互校。 进一步讨论可参照 蒙古征服、草原政治与跨区域文献比较研究

\noindent\eventanchor{SYD-1246-MONGOL-281}{帝聞，亟進師，至其地，適大風刮海水去，其淺可渡}\textbf{蒙古。}
蒙古 在 1246年（元中统-13年；公元换算据《元史》本纪纪年） 的材料痕迹是〔帝聞，亟進師，至其地，適大風刮海水去，其淺可渡〕，出处为 《元史》卷003〈本纪第3卷本年条〉。题名保留的城、路、水道或机构名称，是目前可以落地的空间线索。

它承接的条件是：本纪保留灾害或工程的报告地点与朝廷处置；灾情范围需另证。 随后可追踪的变化为：可与地方文书、河工和环境材料比较，不将单点记录外推为全域因果。

证据的边界是：《元史》为明初仓促官修，本纪保存大量实录条目但有编次与纪年问题；行省执行、数字、族称及对外关系须以条格、多语文书和对方史料互校。 蒙古征服、草原政治与跨区域文献比较研究 可用于继续检验这一结论。

\subsection{1246年}
\noindent\eventanchor{YUAN-1246-EVENT-163}{贵由即位为大汗}\textbf{蒙古。}
从 《元史》相关本纪；《元朝秘史》、波斯文史籍或《高丽史》对应叙事；《元典章》《通制条格》、多语碑刻与文书（据事核） 的 1246年 条目看，〔贵由即位为大汗〕使 蒙古 的一个具体行动进入叙事；材料未将地点细化到城址，不能用中枢所在地替它补足。

前一层压力可由以下线索辨认：窝阔台至蒙哥时期的扩张与继承中前任统治的结束与宫廷、部族或官僚的权力安排相互交织，促成“贵由即位为大汗”这一继承节点。 这一处置留下的后续条件是：继承并不自动消除既有矛盾；“贵由即位为大汗”后的任官、军权和对外策略需由后续条目追踪。

《元史》明初速修，纪年与人名有异同；蒙古大汗政权不等同所有汗国，征服、赋役和身份分类须按地区及材料语言分列。 因此，蒙古帝国、元朝行省财政、多语文书与区域社会研究 只能扩展问题，不能代替未保存的细节。

\subsection{1247 年（共享年序桥接）}
\noindent\eventanchor{SY-1247-CHRONOLOGY-BRIDGE}{【C级年序桥接】保留1247 年的多政权并行检索入口；本条不主张所列政权在当年发生直接互动，具体事件须另以单项材料入账。}\textbf{南宋、金、西夏与蒙古（年序桥接）。}
1247 年（共享年序桥接） 的记载将 南宋、金、西夏与蒙古（年序桥接） 与〔【C级年序桥接】保留1247 年的多政权并行检索入口；本条不主张所列政权在当年发生直接互动，具体事件须另以单项材料入账。〕相连。此处可确认的是这项被记录的行动；题名保留的城、路、水道或机构名称，是目前可以落地的空间线索。

相邻已入账事件之间缺少该年度的共享年序入口，易使跨政权材料被单一政权叙事遮蔽。 记录所见的结果则是：保留该年度的检索与补账位置；后续仅在获得可核单项材料时拆分为具体政治、财政、军事、社会或文化事件。 日期相邻不等于因果已经证实。

关于可说范围，须保留：本条只确认该年位于可追踪的共享年序中；正史、地方文书和外部材料的保存密度不一，不能由桥接标签推断行动、统属、疆界或社会影响。 进一步讨论可参照 蒙古征服、区域政权转换与跨语种年序研究

\subsection{1247年（元中统-12年；公元换算据《元史》本纪纪年）}
\noindent\eventanchor{SYD-1247-MONGOL-037}{八月，忽必烈次臨洮，命總帥汪田哥以城利州聞，欲為取蜀之計}\textbf{蒙古。}
〔八月，忽必烈次臨洮，命總帥汪田哥以城利州聞，欲為取蜀之計〕见于 《元史》卷003〈本纪第3卷本年条〉 的 1247年（元中统-12年；公元换算据《元史》本纪纪年） 记录。蒙古 的选择因此有了可回查的时间位置；题名保留的城、路、水道或机构名称，是目前可以落地的空间线索。

把本项放回事件链，能够看到的前提为：本纪年条记录政令或机构处置；实施背景须参照制度材料。 而它所打开或收紧的下一步是：可在同年及后续政令中追踪；条文不单独证明地方执行。

证据的边界是：《元史》为明初仓促官修，本纪保存大量实录条目但有编次与纪年问题；行省执行、数字、族称及对外关系须以条格、多语文书和对方史料互校。 蒙古征服、草原政治与跨区域文献比较研究 可用于继续检验这一结论。

\noindent\eventanchor{SYD-1247-MONGOL-087}{定宗后及失烈門母以厭禳事覺，並賜死}\textbf{蒙古。}
1247年（元中统-12年；公元换算据《元史》本纪纪年），《元史》卷003〈本纪第3卷本年条〉 所见的〔定宗后及失烈門母以厭禳事覺，並賜死〕把 蒙古 的处置留在可复核的年序中；材料未将地点细化到城址，不能用中枢所在地替它补足。

它承接的条件是：本纪年条提供日期和中枢可见行动；前因不据单条补定。 随后可追踪的变化为：可回查同年及后续条目，地方影响仍须另证。

《元史》为明初仓促官修，本纪保存大量实录条目但有编次与纪年问题；行省执行、数字、族称及对外关系须以条格、多语文书和对方史料互校。 因此，蒙古征服、草原政治与跨区域文献比较研究 只能扩展问题，不能代替未保存的细节。

\noindent\eventanchor{SYD-1247-MONGOL-134}{冬十月，命諸王也古征高麗}\textbf{蒙古。}
蒙古 在 1247年（元中统-12年；公元换算据《元史》本纪纪年） 的材料痕迹是〔冬十月，命諸王也古征高麗〕，出处为 《元史》卷003〈本纪第3卷本年条〉。材料未将地点细化到城址，不能用中枢所在地替它补足。

前一层压力可由以下线索辨认：本纪年条记录政令或机构处置；实施背景须参照制度材料。 这一处置留下的后续条件是：可在同年及后续政令中追踪；条文不单独证明地方执行。

关于可说范围，须保留：《元史》为明初仓促官修，本纪保存大量实录条目但有编次与纪年问题；行省执行、数字、族称及对外关系须以条格、多语文书和对方史料互校。 进一步讨论可参照 蒙古征服、草原政治与跨区域文献比较研究

\noindent\eventanchor{SYD-1247-MONGOL-176}{分遷諸王於各所：合丹於別石八里地，蔑里於葉兒的石河，海都於海押立地，別兒哥於曲兒只地，脫脫於葉密立地，蒙哥都及太宗皇后乞里吉忽帖尼於擴端所居地之西}\textbf{蒙古。}
从 《元史》卷003〈本纪第3卷本年条〉 的 1247年（元中统-12年；公元换算据《元史》本纪纪年） 条目看，〔分遷諸王於各所：合丹於別石八里地，蔑里於葉兒的石河，海都於海押立地，別兒哥於曲兒只地，脫脫於葉密立地，蒙哥都及太宗皇后乞里吉忽帖尼於擴端所居地之西〕使 蒙古 的一个具体行动进入叙事；题名保留的城、路、水道或机构名称，是目前可以落地的空间线索。

本纪保留灾害或工程的报告地点与朝廷处置；灾情范围需另证。 记录所见的结果则是：可与地方文书、河工和环境材料比较，不将单点记录外推为全域因果。 日期相邻不等于因果已经证实。

证据的边界是：《元史》为明初仓促官修，本纪保存大量实录条目但有编次与纪年问题；行省执行、数字、族称及对外关系须以条格、多语文书和对方史料互校。 蒙古征服、草原政治与跨区域文献比较研究 可用于继续检验这一结论。

\noindent\eventanchor{SYD-1247-MONGOL-216}{禁錮和只、納忽、〔也〕孫脫等於軍營}\textbf{蒙古。}
1247年（元中统-12年；公元换算据《元史》本纪纪年） 的记载将 蒙古 与〔禁錮和只、納忽、〔也〕孫脫等於軍營〕相连。此处可确认的是这项被记录的行动；材料未将地点细化到城址，不能用中枢所在地替它补足。

把本项放回事件链，能够看到的前提为：本纪年条记录政令或机构处置；实施背景须参照制度材料。 而它所打开或收紧的下一步是：可在同年及后续政令中追踪；条文不单独证明地方执行。

《元史》为明初仓促官修，本纪保存大量实录条目但有编次与纪年问题；行省执行、数字、族称及对外关系须以条格、多语文书和对方史料互校。 因此，蒙古征服、草原政治与跨区域文献比较研究 只能扩展问题，不能代替未保存的细节。

\noindent\eventanchor{SYD-1247-MONGOL-250}{遣乞都不花攻末來吉兒都怯寨}\textbf{蒙古。}
〔遣乞都不花攻末來吉兒都怯寨〕见于 《元史》卷003〈本纪第3卷本年条〉 的 1247年（元中统-12年；公元换算据《元史》本纪纪年） 记录。蒙古 的选择因此有了可回查的时间位置；题名保留的城、路、水道或机构名称，是目前可以落地的空间线索。

它承接的条件是：本纪从朝廷视角记录军政行动；出兵与战果需分辨。 随后可追踪的变化为：后续路线、控制与损失应与对方记录和遗址材料复核。

关于可说范围，须保留：《元史》为明初仓促官修，本纪保存大量实录条目但有编次与纪年问题；行省执行、数字、族称及对外关系须以条格、多语文书和对方史料互校。 进一步讨论可参照 蒙古征服、草原政治与跨区域文献比较研究

\noindent\eventanchor{SYD-1247-MONGOL-282}{秋七月，命忽必烈征大理，諸王禿兒花、撒〔立〕征身毒，怯的不花征沒里奚，旭烈征西域素丹諸國}\textbf{蒙古。}
1247年（元中统-12年；公元换算据《元史》本纪纪年），《元史》卷003〈本纪第3卷本年条〉 所见的〔秋七月，命忽必烈征大理，諸王禿兒花、撒〔立〕征身毒，怯的不花征沒里奚，旭烈征西域素丹諸國〕把 蒙古 的处置留在可复核的年序中；材料未将地点细化到城址，不能用中枢所在地替它补足。

前一层压力可由以下线索辨认：本纪年条记录政令或机构处置；实施背景须参照制度材料。 这一处置留下的后续条件是：可在同年及后续政令中追踪；条文不单独证明地方执行。

证据的边界是：《元史》为明初仓促官修，本纪保存大量实录条目但有编次与纪年问题；行省执行、数字、族称及对外关系须以条格、多语文书和对方史料互校。 蒙古征服、草原政治与跨区域文献比较研究 可用于继续检验这一结论。

\subsection{1248 年（共享年序桥接）}
\noindent\eventanchor{SY-1248-CHRONOLOGY-BRIDGE}{【C级年序桥接】保留1248 年的多政权并行检索入口；本条不主张所列政权在当年发生直接互动，具体事件须另以单项材料入账。}\textbf{南宋、金、西夏与蒙古（年序桥接）。}
南宋、金、西夏与蒙古（年序桥接） 在 1248 年（共享年序桥接） 的材料痕迹是〔【C级年序桥接】保留1248 年的多政权并行检索入口；本条不主张所列政权在当年发生直接互动，具体事件须另以单项材料入账。〕，出处为 《宋史》；《金史》；《元史》；蒙古文、波斯文及地方材料（据事核）。题名保留的城、路、水道或机构名称，是目前可以落地的空间线索。

相邻已入账事件之间缺少该年度的共享年序入口，易使跨政权材料被单一政权叙事遮蔽。 记录所见的结果则是：保留该年度的检索与补账位置；后续仅在获得可核单项材料时拆分为具体政治、财政、军事、社会或文化事件。 日期相邻不等于因果已经证实。

本条只确认该年位于可追踪的共享年序中；正史、地方文书和外部材料的保存密度不一，不能由桥接标签推断行动、统属、疆界或社会影响。 因此，蒙古征服、区域政权转换与跨语种年序研究 只能扩展问题，不能代替未保存的细节。

\subsection{1248年（元中统-11年；公元换算据《元史》本纪纪年）}
\noindent\eventanchor{SYD-1248-MONGOL-038}{罷也古征高麗兵，以札剌兒帶為征東元帥}\textbf{蒙古。}
从 《元史》卷003〈本纪第3卷本年条〉 的 1248年（元中统-11年；公元换算据《元史》本纪纪年） 条目看，〔罷也古征高麗兵，以札剌兒帶為征東元帥〕使 蒙古 的一个具体行动进入叙事；材料未将地点细化到城址，不能用中枢所在地替它补足。

把本项放回事件链，能够看到的前提为：本纪年条记录政令或机构处置；实施背景须参照制度材料。 而它所打开或收紧的下一步是：可在同年及后续政令中追踪；条文不单独证明地方执行。

关于可说范围，须保留：《元史》为明初仓促官修，本纪保存大量实录条目但有编次与纪年问题；行省执行、数字、族称及对外关系须以条格、多语文书和对方史料互校。 进一步讨论可参照 蒙古征服、草原政治与跨区域文献比较研究

\noindent\eventanchor{SYD-1248-MONGOL-088}{此銀就充今後歲賜之數}\textbf{蒙古。}
1248年（元中统-11年；公元换算据《元史》本纪纪年） 的记载将 蒙古 与〔此銀就充今後歲賜之數〕相连。此处可确认的是这项被记录的行动；材料未将地点细化到城址，不能用中枢所在地替它补足。

它承接的条件是：本纪年条提供日期和中枢可见行动；前因不据单条补定。 随后可追踪的变化为：可回查同年及后续条目，地方影响仍须另证。

证据的边界是：《元史》为明初仓促官修，本纪保存大量实录条目但有编次与纪年问题；行省执行、数字、族称及对外关系须以条格、多语文书和对方史料互校。 蒙古征服、草原政治与跨区域文献比较研究 可用于继续检验这一结论。

\noindent\eventanchor{SYD-1248-MONGOL-135}{帝遂會諸王于斡難河北，賜予甚厚}\textbf{蒙古。}
〔帝遂會諸王于斡難河北，賜予甚厚〕见于 《元史》卷003〈本纪第3卷本年条〉 的 1248年（元中统-11年；公元换算据《元史》本纪纪年） 记录。蒙古 的选择因此有了可回查的时间位置；题名保留的城、路、水道或机构名称，是目前可以落地的空间线索。

前一层压力可由以下线索辨认：本纪保留灾害或工程的报告地点与朝廷处置；灾情范围需另证。 这一处置留下的后续条件是：可与地方文书、河工和环境材料比较，不将单点记录外推为全域因果。

《元史》为明初仓促官修，本纪保存大量实录条目但有编次与纪年问题；行省执行、数字、族称及对外关系须以条格、多语文书和对方史料互校。 因此，蒙古征服、草原政治与跨区域文献比较研究 只能扩展问题，不能代替未保存的细节。

\noindent\eventanchor{SYD-1248-MONGOL-177}{帝幸火兒忽納要不〔兒〕之地}\textbf{蒙古。}
1248年（元中统-11年；公元换算据《元史》本纪纪年），《元史》卷003〈本纪第3卷本年条〉 所见的〔帝幸火兒忽納要不〔兒〕之地〕把 蒙古 的处置留在可复核的年序中；材料未将地点细化到城址，不能用中枢所在地替它补足。

本纪年条提供日期和中枢可见行动；前因不据单条补定。 记录所见的结果则是：可回查同年及后续条目，地方影响仍须另证。 日期相邻不等于因果已经证实。

关于可说范围，须保留：《元史》为明初仓促官修，本纪保存大量实录条目但有编次与纪年问题；行省执行、数字、族称及对外关系须以条格、多语文书和对方史料互校。 进一步讨论可参照 蒙古征服、草原政治与跨区域文献比较研究

\noindent\eventanchor{SYD-1248-MONGOL-217}{命宗王耶虎與洪福源同領軍征高麗，攻拔禾山、東州、春州、三角山、楊根、天龍等城}\textbf{蒙古。}
蒙古 在 1248年（元中统-11年；公元换算据《元史》本纪纪年） 的材料痕迹是〔命宗王耶虎與洪福源同領軍征高麗，攻拔禾山、東州、春州、三角山、楊根、天龍等城〕，出处为 《元史》卷003〈本纪第3卷本年条〉。题名保留的城、路、水道或机构名称，是目前可以落地的空间线索。

把本项放回事件链，能够看到的前提为：本纪年条记录政令或机构处置；实施背景须参照制度材料。 而它所打开或收紧的下一步是：可在同年及后续政令中追踪；条文不单独证明地方执行。

证据的边界是：《元史》为明初仓促官修，本纪保存大量实录条目但有编次与纪年问题；行省执行、数字、族称及对外关系须以条格、多语文书和对方史料互校。 蒙古征服、草原政治与跨区域文献比较研究 可用于继续检验这一结论。

\noindent\eventanchor{SYD-1248-MONGOL-251}{遣必闍別兒哥括斡羅思戶口}\textbf{蒙古。}
从 《元史》卷003〈本纪第3卷本年条〉 的 1248年（元中统-11年；公元换算据《元史》本纪纪年） 条目看，〔遣必闍別兒哥括斡羅思戶口〕使 蒙古 的一个具体行动进入叙事；材料未将地点细化到城址，不能用中枢所在地替它补足。

它承接的条件是：本纪年条提供日期和中枢可见行动；前因不据单条补定。 随后可追踪的变化为：可回查同年及后续条目，地方影响仍须另证。

《元史》为明初仓促官修，本纪保存大量实录条目但有编次与纪年问题；行省执行、数字、族称及对外关系须以条格、多语文书和对方史料互校。 因此，蒙古征服、草原政治与跨区域文献比较研究 只能扩展问题，不能代替未保存的细节。

\noindent\eventanchor{SYD-1248-MONGOL-283}{三年癸丑春正月，汪田哥修治利州，且屯田，蜀人莫敢侵軼}\textbf{蒙古。}
1248年（元中统-11年；公元换算据《元史》本纪纪年） 的记载将 蒙古 与〔三年癸丑春正月，汪田哥修治利州，且屯田，蜀人莫敢侵軼〕相连。此处可确认的是这项被记录的行动；题名保留的城、路、水道或机构名称，是目前可以落地的空间线索。

前一层压力可由以下线索辨认：本纪年条提供日期和中枢可见行动；前因不据单条补定。 这一处置留下的后续条件是：可回查同年及后续条目，地方影响仍须另证。

关于可说范围，须保留：《元史》为明初仓促官修，本纪保存大量实录条目但有编次与纪年问题；行省执行、数字、族称及对外关系须以条格、多语文书和对方史料互校。 进一步讨论可参照 蒙古征服、草原政治与跨区域文献比较研究

\subsection{1249 年（共享年序桥接）}
\noindent\eventanchor{SY-1249-CHRONOLOGY-BRIDGE}{【C级年序桥接】保留1249 年的多政权并行检索入口；本条不主张所列政权在当年发生直接互动，具体事件须另以单项材料入账。}\textbf{南宋、金、西夏与蒙古（年序桥接）。}
〔【C级年序桥接】保留1249 年的多政权并行检索入口；本条不主张所列政权在当年发生直接互动，具体事件须另以单项材料入账。〕见于 《宋史》；《金史》；《元史》；蒙古文、波斯文及地方材料（据事核） 的 1249 年（共享年序桥接） 记录。南宋、金、西夏与蒙古（年序桥接） 的选择因此有了可回查的时间位置；题名保留的城、路、水道或机构名称，是目前可以落地的空间线索。

相邻已入账事件之间缺少该年度的共享年序入口，易使跨政权材料被单一政权叙事遮蔽。 记录所见的结果则是：保留该年度的检索与补账位置；后续仅在获得可核单项材料时拆分为具体政治、财政、军事、社会或文化事件。 日期相邻不等于因果已经证实。

证据的边界是：本条只确认该年位于可追踪的共享年序中；正史、地方文书和外部材料的保存密度不一，不能由桥接标签推断行动、统属、疆界或社会影响。 蒙古征服、区域政权转换与跨语种年序研究 可用于继续检验这一结论。

\subsection{1249年（元中统-10年；公元换算据《元史》本纪纪年）}
\noindent\eventanchor{SYD-1249-MONGOL-039}{帝謂大臣，求可以慎固封守、閑於將略者}\textbf{蒙古。}
1249年（元中统-10年；公元换算据《元史》本纪纪年），《元史》卷003〈本纪第3卷本年条〉 所见的〔帝謂大臣，求可以慎固封守、閑於將略者〕把 蒙古 的处置留在可复核的年序中；材料未将地点细化到城址，不能用中枢所在地替它补足。

把本项放回事件链，能够看到的前提为：本纪年条提供日期和中枢可见行动；前因不据单条补定。 而它所打开或收紧的下一步是：可回查同年及后续条目，地方影响仍须另证。

《元史》为明初仓促官修，本纪保存大量实录条目但有编次与纪年问题；行省执行、数字、族称及对外关系须以条格、多语文书和对方史料互校。 因此，蒙古征服、草原政治与跨区域文献比较研究 只能扩展问题，不能代替未保存的细节。

\noindent\eventanchor{SYD-1249-MONGOL-089}{忽必烈還自大理，留兀良合台攻諸夷之未附者，入覲於獵所}\textbf{蒙古。}
蒙古 在 1249年（元中统-10年；公元换算据《元史》本纪纪年） 的材料痕迹是〔忽必烈還自大理，留兀良合台攻諸夷之未附者，入覲於獵所〕，出处为 《元史》卷003〈本纪第3卷本年条〉。材料未将地点细化到城址，不能用中枢所在地替它补足。

它承接的条件是：本纪从朝廷视角记录军政行动；出兵与战果需分辨。 随后可追踪的变化为：后续路线、控制与损失应与对方记录和遗址材料复核。

关于可说范围，须保留：《元史》为明初仓促官修，本纪保存大量实录条目但有编次与纪年问题；行省执行、数字、族称及对外关系须以条格、多语文书和对方史料互校。 进一步讨论可参照 蒙古征服、草原政治与跨区域文献比较研究

\noindent\eventanchor{SYD-1249-MONGOL-136}{均州總管孫嗣遣人齎蠟書降，且乞援，史權以精甲備宋人之要，遂援嗣而來}\textbf{蒙古与南宋。}
从 《元史》卷003〈本纪第3卷本年条〉 的 1249年（元中统-10年；公元换算据《元史》本纪纪年） 条目看，〔均州總管孫嗣遣人齎蠟書降，且乞援，史權以精甲備宋人之要，遂援嗣而來〕使 蒙古与南宋 的一个具体行动进入叙事；题名保留的城、路、水道或机构名称，是目前可以落地的空间线索。

前一层压力可由以下线索辨认：本纪从朝廷视角记录军政行动；出兵与战果需分辨。 这一处置留下的后续条件是：后续路线、控制与损失应与对方记录和遗址材料复核。

证据的边界是：《元史》为明初仓促官修，本纪保存大量实录条目但有编次与纪年问题；行省执行、数字、族称及对外关系须以条格、多语文书和对方史料互校。 蒙古征服、草原政治与跨区域文献比较研究 可用于继续检验这一结论。

\noindent\eventanchor{SYD-1249-MONGOL-178}{其後驍將〔鍾〕顯、王梅、杜柔、袁師信各帥所部來降}\textbf{蒙古。}
1249年（元中统-10年；公元换算据《元史》本纪纪年） 的记载将 蒙古 与〔其後驍將〔鍾〕顯、王梅、杜柔、袁師信各帥所部來降〕相连。此处可确认的是这项被记录的行动；材料未将地点细化到城址，不能用中枢所在地替它补足。

本纪从朝廷视角记录军政行动；出兵与战果需分辨。 记录所见的结果则是：后续路线、控制与损失应与对方记录和遗址材料复核。 日期相邻不等于因果已经证实。

《元史》为明初仓促官修，本纪保存大量实录条目但有编次与纪年问题；行省执行、数字、族称及对外关系须以条格、多语文书和对方史料互校。 因此，蒙古征服、草原政治与跨区域文献比较研究 只能扩展问题，不能代替未保存的细节。

\noindent\eventanchor{SYD-1249-MONGOL-218}{遣札剌亦兒部人火兒赤征高麗}\textbf{蒙古。}
〔遣札剌亦兒部人火兒赤征高麗〕见于 《元史》卷003〈本纪第3卷本年条〉 的 1249年（元中统-10年；公元换算据《元史》本纪纪年） 记录。蒙古 的选择因此有了可回查的时间位置；材料未将地点细化到城址，不能用中枢所在地替它补足。

把本项放回事件链，能够看到的前提为：本纪年条提供日期和中枢可见行动；前因不据单条补定。 而它所打开或收紧的下一步是：可回查同年及后续条目，地方影响仍须另证。

关于可说范围，须保留：《元史》为明初仓促官修，本纪保存大量实录条目但有编次与纪年问题；行省执行、数字、族称及对外关系须以条格、多语文书和对方史料互校。 进一步讨论可参照 蒙古征服、草原政治与跨区域文献比较研究

\noindent\eventanchor{SYD-1249-MONGOL-252}{秋七月，詔官吏之赴朝理算錢糧者，許自首不公，仍禁以後浮費}\textbf{蒙古。}
1249年（元中统-10年；公元换算据《元史》本纪纪年），《元史》卷003〈本纪第3卷本年条〉 所见的〔秋七月，詔官吏之赴朝理算錢糧者，許自首不公，仍禁以後浮費〕把 蒙古 的处置留在可复核的年序中；材料未将地点细化到城址，不能用中枢所在地替它补足。

它承接的条件是：本纪年条记录政令或机构处置；实施背景须参照制度材料。 随后可追踪的变化为：可在同年及后续政令中追踪；条文不单独证明地方执行。

证据的边界是：《元史》为明初仓促官修，本纪保存大量实录条目但有编次与纪年问题；行省执行、数字、族称及对外关系须以条格、多语文书和对方史料互校。 蒙古征服、草原政治与跨区域文献比较研究 可用于继续检验这一结论。

\noindent\eventanchor{SYD-1249-MONGOL-284}{柔又以渦水北隘淺不可舟，軍既病涉，曹、濮、魏、博粟皆不至，乃築甬路自亳抵汴，堤百二十里，流深而不能築，復為橋十五，或廣八十尺，橫以二堡戍之}\textbf{蒙古。}
蒙古 在 1249年（元中统-10年；公元换算据《元史》本纪纪年） 的材料痕迹是〔柔又以渦水北隘淺不可舟，軍既病涉，曹、濮、魏、博粟皆不至，乃築甬路自亳抵汴，堤百二十里，流深而不能築，復為橋十五，或廣八十尺，橫以二堡戍之〕，出处为 《元史》卷003〈本纪第3卷本年条〉。题名保留的城、路、水道或机构名称，是目前可以落地的空间线索。

前一层压力可由以下线索辨认：本纪年条记录政令或机构处置；实施背景须参照制度材料。 这一处置留下的后续条件是：可在同年及后续政令中追踪；条文不单独证明地方执行。

《元史》为明初仓促官修，本纪保存大量实录条目但有编次与纪年问题；行省执行、数字、族称及对外关系须以条格、多语文书和对方史料互校。 因此，蒙古征服、草原政治与跨区域文献比较研究 只能扩展问题，不能代替未保存的细节。

\subsection{1250 年（共享年序桥接）}
\noindent\eventanchor{SY-1250-CHRONOLOGY-BRIDGE}{【C级年序桥接】保留1250 年的多政权并行检索入口；本条不主张所列政权在当年发生直接互动，具体事件须另以单项材料入账。}\textbf{南宋、金、西夏与蒙古（年序桥接）。}
从 《宋史》；《金史》；《元史》；蒙古文、波斯文及地方材料（据事核） 的 1250 年（共享年序桥接） 条目看，〔【C级年序桥接】保留1250 年的多政权并行检索入口；本条不主张所列政权在当年发生直接互动，具体事件须另以单项材料入账。〕使 南宋、金、西夏与蒙古（年序桥接） 的一个具体行动进入叙事；题名保留的城、路、水道或机构名称，是目前可以落地的空间线索。

相邻已入账事件之间缺少该年度的共享年序入口，易使跨政权材料被单一政权叙事遮蔽。 记录所见的结果则是：保留该年度的检索与补账位置；后续仅在获得可核单项材料时拆分为具体政治、财政、军事、社会或文化事件。 日期相邻不等于因果已经证实。

关于可说范围，须保留：本条只确认该年位于可追踪的共享年序中；正史、地方文书和外部材料的保存密度不一，不能由桥接标签推断行动、统属、疆界或社会影响。 进一步讨论可参照 蒙古征服、区域政权转换与跨语种年序研究

\subsection{1250年（元中统-9年；公元换算据《元史》本纪纪年）}
\noindent\eventanchor{SYD-1250-MONGOL-040}{後此又連三歲，攻拔其光州、安城、〔忠〕州、玄〔風〕、珍原、甲向、玉果等城}\textbf{蒙古。}
1250年（元中统-9年；公元换算据《元史》本纪纪年） 的记载将 蒙古 与〔後此又連三歲，攻拔其光州、安城、〔忠〕州、玄〔風〕、珍原、甲向、玉果等城〕相连。此处可确认的是这项被记录的行动；题名保留的城、路、水道或机构名称，是目前可以落地的空间线索。

把本项放回事件链，能够看到的前提为：本纪从朝廷视角记录军政行动；出兵与战果需分辨。 而它所打开或收紧的下一步是：后续路线、控制与损失应与对方记录和遗址材料复核。

证据的边界是：《元史》为明初仓促官修，本纪保存大量实录条目但有编次与纪年问题；行省执行、数字、族称及对外关系须以条格、多语文书和对方史料互校。 蒙古征服、草原政治与跨区域文献比较研究 可用于继续检验这一结论。

\noindent\eventanchor{SYD-1250-MONGOL-090}{是歲，改命劄剌䚟與洪福源同征高麗}\textbf{蒙古。}
〔是歲，改命劄剌䚟與洪福源同征高麗〕见于 《元史》卷003〈本纪第3卷本年条〉 的 1250年（元中统-9年；公元换算据《元史》本纪纪年） 记录。蒙古 的选择因此有了可回查的时间位置；材料未将地点细化到城址，不能用中枢所在地替它补足。

它承接的条件是：本纪年条记录政令或机构处置；实施背景须参照制度材料。 随后可追踪的变化为：可在同年及后续政令中追踪；条文不单独证明地方执行。

《元史》为明初仓促官修，本纪保存大量实录条目但有编次与纪年问题；行省执行、数字、族称及对外关系须以条格、多语文书和对方史料互校。 因此，蒙古征服、草原政治与跨区域文献比较研究 只能扩展问题，不能代替未保存的细节。

\noindent\eventanchor{SYD-1250-MONGOL-137}{五年乙卯春，詔徵逋欠錢穀}\textbf{蒙古。}
1250年（元中统-9年；公元换算据《元史》本纪纪年），《元史》卷003〈本纪第3卷本年条〉 所见的〔五年乙卯春，詔徵逋欠錢穀〕把 蒙古 的处置留在可复核的年序中；材料未将地点细化到城址，不能用中枢所在地替它补足。

前一层压力可由以下线索辨认：本纪年条记录政令或机构处置；实施背景须参照制度材料。 这一处置留下的后续条件是：可在同年及后续政令中追踪；条文不单独证明地方执行。

关于可说范围，须保留：《元史》为明初仓促官修，本纪保存大量实录条目但有编次与纪年问题；行省执行、数字、族称及对外关系须以条格、多语文书和对方史料互校。 进一步讨论可参照 蒙古征服、草原政治与跨区域文献比较研究

\noindent\eventanchor{SYD-1250-MONGOL-179}{以百丈口為宋往來之道，可容萬艘，遂築甬路，自亳而南六十餘里，中為橫江堡}\textbf{蒙古与南宋。}
蒙古与南宋 在 1250年（元中统-9年；公元换算据《元史》本纪纪年） 的材料痕迹是〔以百丈口為宋往來之道，可容萬艘，遂築甬路，自亳而南六十餘里，中為橫江堡〕，出处为 《元史》卷003〈本纪第3卷本年条〉。题名保留的城、路、水道或机构名称，是目前可以落地的空间线索。

本纪年条提供日期和中枢可见行动；前因不据单条补定。 记录所见的结果则是：可回查同年及后续条目，地方影响仍须另证。 日期相邻不等于因果已经证实。

证据的边界是：《元史》为明初仓促官修，本纪保存大量实录条目但有编次与纪年问题；行省执行、数字、族称及对外关系须以条格、多语文书和对方史料互校。 蒙古征服、草原政治与跨区域文献比较研究 可用于继续检验这一结论。

\noindent\eventanchor{SYD-1250-MONGOL-219}{又以路東六十里皆水，可致宋舟，乃立柵水中，惟密置偵邏於所達之路}\textbf{蒙古与南宋。}
从 《元史》卷003〈本纪第3卷本年条〉 的 1250年（元中统-9年；公元换算据《元史》本纪纪年） 条目看，〔又以路東六十里皆水，可致宋舟，乃立柵水中，惟密置偵邏於所達之路〕使 蒙古与南宋 的一个具体行动进入叙事；题名保留的城、路、水道或机构名称，是目前可以落地的空间线索。

把本项放回事件链，能够看到的前提为：本纪年条记录政令或机构处置；实施背景须参照制度材料。 而它所打开或收紧的下一步是：可在同年及后续政令中追踪；条文不单独证明地方执行。

《元史》为明初仓促官修，本纪保存大量实录条目但有编次与纪年问题；行省执行、数字、族称及对外关系须以条格、多语文书和对方史料互校。 因此，蒙古征服、草原政治与跨区域文献比较研究 只能扩展问题，不能代替未保存的细节。

\subsection{1251年（元中统-8年；公元换算据《元史》本纪纪年）}
\noindent\eventanchor{SYD-1251-MONGOL-041}{賜金縷織文衣一襲、銀五十兩、綵帛萬二百匹，以賚軍士}\textbf{蒙古。}
1251年（元中统-8年；公元换算据《元史》本纪纪年） 的记载将 蒙古 与〔賜金縷織文衣一襲、銀五十兩、綵帛萬二百匹，以賚軍士〕相连。此处可确认的是这项被记录的行动；材料未将地点细化到城址，不能用中枢所在地替它补足。

它承接的条件是：本纪年条提供日期和中枢可见行动；前因不据单条补定。 随后可追踪的变化为：可回查同年及后续条目，地方影响仍须另证。

关于可说范围，须保留：《元史》为明初仓促官修，本纪保存大量实录条目但有编次与纪年问题；行省执行、数字、族称及对外关系须以条格、多语文书和对方史料互校。 进一步讨论可参照 蒙古征服、草原政治与跨区域文献比较研究

\noindent\eventanchor{SYD-1251-MONGOL-091}{帝會諸王、百官于欲兒陌哥都之地，設宴六十餘日，賜金帛有差，仍定擬諸王歲賜錢穀}\textbf{蒙古。}
〔帝會諸王、百官于欲兒陌哥都之地，設宴六十餘日，賜金帛有差，仍定擬諸王歲賜錢穀〕见于 《元史》卷003〈本纪第3卷本年条〉 的 1251年（元中统-8年；公元换算据《元史》本纪纪年） 记录。蒙古 的选择因此有了可回查的时间位置；题名保留的城、路、水道或机构名称，是目前可以落地的空间线索。

前一层压力可由以下线索辨认：本纪所记财用、赈济或征发属于特定报告场景。 这一处置留下的后续条件是：可与食货志、文书和物质材料比较，不外推为全境总量。

证据的边界是：《元史》为明初仓促官修，本纪保存大量实录条目但有编次与纪年问题；行省执行、数字、族称及对外关系须以条格、多语文书和对方史料互校。 蒙古征服、草原政治与跨区域文献比较研究 可用于继续检验这一结论。

\noindent\eventanchor{SYD-1251-MONGOL-138}{帝亦以宋人違命囚使，會議伐之}\textbf{蒙古与南宋。}
1251年（元中统-8年；公元换算据《元史》本纪纪年），《元史》卷003〈本纪第3卷本年条〉 所见的〔帝亦以宋人違命囚使，會議伐之〕把 蒙古与南宋 的处置留在可复核的年序中；材料未将地点细化到城址，不能用中枢所在地替它补足。

本纪年条记录政令或机构处置；实施背景须参照制度材料。 记录所见的结果则是：可在同年及后续政令中追踪；条文不单独证明地方执行。 日期相邻不等于因果已经证实。

《元史》为明初仓促官修，本纪保存大量实录条目但有编次与纪年问题；行省执行、数字、族称及对外关系须以条格、多语文书和对方史料互校。 因此，蒙古征服、草原政治与跨区域文献比较研究 只能扩展问题，不能代替未保存的细节。

\noindent\eventanchor{SYD-1251-MONGOL-180}{忽必烈遣沒兒合石詣行在所，奏請續簽內郡漢軍，從之}\textbf{蒙古。}
蒙古 在 1251年（元中统-8年；公元换算据《元史》本纪纪年） 的材料痕迹是〔忽必烈遣沒兒合石詣行在所，奏請續簽內郡漢軍，從之〕，出处为 《元史》卷003〈本纪第3卷本年条〉。材料未将地点细化到城址，不能用中枢所在地替它补足。

把本项放回事件链，能够看到的前提为：本纪年条提供日期和中枢可见行动；前因不据单条补定。 而它所打开或收紧的下一步是：可回查同年及后续条目，地方影响仍须另证。

关于可说范围，须保留：《元史》为明初仓促官修，本纪保存大量实录条目但有编次与纪年问题；行省执行、数字、族称及对外关系须以条格、多语文书和对方史料互校。 进一步讨论可参照 蒙古征服、草原政治与跨区域文献比较研究

\noindent\eventanchor{SYD-1251-MONGOL-220}{秋七月，命諸王各還所部以居}\textbf{蒙古。}
从 《元史》卷003〈本纪第3卷本年条〉 的 1251年（元中统-8年；公元换算据《元史》本纪纪年） 条目看，〔秋七月，命諸王各還所部以居〕使 蒙古 的一个具体行动进入叙事；材料未将地点细化到城址，不能用中枢所在地替它补足。

它承接的条件是：本纪年条记录政令或机构处置；实施背景须参照制度材料。 随后可追踪的变化为：可在同年及后续政令中追踪；条文不单独证明地方执行。

证据的边界是：《元史》为明初仓促官修，本纪保存大量实录条目但有编次与纪年问题；行省执行、数字、族称及对外关系须以条格、多语文书和对方史料互校。 蒙古征服、草原政治与跨区域文献比较研究 可用于继续检验这一结论。

\noindent\eventanchor{SYD-1251-MONGOL-253}{是歲，〔波〕麗國王細嵯甫、雲南酋長摩合羅嵯及素丹諸國來覲}\textbf{蒙古。}
1251年（元中统-8年；公元换算据《元史》本纪纪年） 的记载将 蒙古 与〔是歲，〔波〕麗國王細嵯甫、雲南酋長摩合羅嵯及素丹諸國來覲〕相连。此处可确认的是这项被记录的行动；材料未将地点细化到城址，不能用中枢所在地替它补足。

前一层压力可由以下线索辨认：本纪年条提供日期和中枢可见行动；前因不据单条补定。 这一处置留下的后续条件是：可回查同年及后续条目，地方影响仍须另证。

《元史》为明初仓促官修，本纪保存大量实录条目但有编次与纪年问题；行省执行、数字、族称及对外关系须以条格、多语文书和对方史料互校。 因此，蒙古征服、草原政治与跨区域文献比较研究 只能扩展问题，不能代替未保存的细节。

\noindent\eventanchor{SYD-1251-MONGOL-285}{遂自昔八兒地還至重慶府，敗宋將張都統}\textbf{蒙古与南宋。}
〔遂自昔八兒地還至重慶府，敗宋將張都統〕见于 《元史》卷003〈本纪第3卷本年条〉 的 1251年（元中统-8年；公元换算据《元史》本纪纪年） 记录。蒙古与南宋 的选择因此有了可回查的时间位置；题名保留的城、路、水道或机构名称，是目前可以落地的空间线索。

本纪从朝廷视角记录军政行动；出兵与战果需分辨。 记录所见的结果则是：后续路线、控制与损失应与对方记录和遗址材料复核。 日期相邻不等于因果已经证实。

关于可说范围，须保留：《元史》为明初仓促官修，本纪保存大量实录条目但有编次与纪年问题；行省执行、数字、族称及对外关系须以条格、多语文书和对方史料互校。 进一步讨论可参照 蒙古征服、草原政治与跨区域文献比较研究

\subsection{1251年}
\noindent\eventanchor{YUAN-1251-EVENT-164}{蒙哥即大汗位并整肃政治对手}\textbf{蒙古。}
1251年，《元史》相关本纪；《元朝秘史》、波斯文史籍或《高丽史》对应叙事；《元典章》《通制条格》、多语碑刻与文书（据事核） 所见的〔蒙哥即大汗位并整肃政治对手〕把 蒙古 的处置留在可复核的年序中；材料未将地点细化到城址，不能用中枢所在地替它补足。

把本项放回事件链，能够看到的前提为：“蒙哥即大汗位并整肃政治对手”发生在窝阔台至蒙哥时期的扩张与继承的具体压力与机会之中，相关行动者的选择不能脱离此前事件链解释。 而它所打开或收紧的下一步是：“蒙哥即大汗位并整肃政治对手”为后续政治、边防或资源配置留下条件；其社会影响与实际覆盖范围仍须以异类材料限定。

证据的边界是：《元史》明初速修，纪年与人名有异同；蒙古大汗政权不等同所有汗国，征服、赋役和身份分类须按地区及材料语言分列。 蒙古帝国、元朝行省财政、多语文书与区域社会研究 可用于继续检验这一结论。

\subsection{1252 年（共享年序桥接）}
\noindent\eventanchor{SY-1252-CHRONOLOGY-BRIDGE}{【C级年序桥接】保留1252 年的多政权并行检索入口；本条不主张所列政权在当年发生直接互动，具体事件须另以单项材料入账。}\textbf{南宋、金、西夏与蒙古（年序桥接）。}
南宋、金、西夏与蒙古（年序桥接） 在 1252 年（共享年序桥接） 的材料痕迹是〔【C级年序桥接】保留1252 年的多政权并行检索入口；本条不主张所列政权在当年发生直接互动，具体事件须另以单项材料入账。〕，出处为 《宋史》；《金史》；《元史》；蒙古文、波斯文及地方材料（据事核）。题名保留的城、路、水道或机构名称，是目前可以落地的空间线索。

它承接的条件是：相邻已入账事件之间缺少该年度的共享年序入口，易使跨政权材料被单一政权叙事遮蔽。 随后可追踪的变化为：保留该年度的检索与补账位置；后续仅在获得可核单项材料时拆分为具体政治、财政、军事、社会或文化事件。

本条只确认该年位于可追踪的共享年序中；正史、地方文书和外部材料的保存密度不一，不能由桥接标签推断行动、统属、疆界或社会影响。 因此，蒙古征服、区域政权转换与跨语种年序研究 只能扩展问题，不能代替未保存的细节。

\subsection{1252年（元中统-7年；公元换算据《元史》本纪纪年）}
\noindent\eventanchor{SYD-1252-MONGOL-042}{安南主陳日煚竄海島，遂班師}\textbf{蒙古。}
从 《元史》卷003〈本纪第3卷本年条〉 的 1252年（元中统-7年；公元换算据《元史》本纪纪年） 条目看，〔安南主陳日煚竄海島，遂班師〕使 蒙古 的一个具体行动进入叙事；题名保留的城、路、水道或机构名称，是目前可以落地的空间线索。

前一层压力可由以下线索辨认：本纪年条记录政令或机构处置；实施背景须参照制度材料。 这一处置留下的后续条件是：可在同年及后续政令中追踪；条文不单独证明地方执行。

关于可说范围，须保留：《元史》为明初仓促官修，本纪保存大量实录条目但有编次与纪年问题；行省执行、数字、族称及对外关系须以条格、多语文书和对方史料互校。 进一步讨论可参照 蒙古征服、草原政治与跨区域文献比较研究

\noindent\eventanchor{SYD-1252-MONGOL-092}{冬十一月，兀良合台伐交趾，敗之，入其國}\textbf{蒙古。}
1252年（元中统-7年；公元换算据《元史》本纪纪年） 的记载将 蒙古 与〔冬十一月，兀良合台伐交趾，敗之，入其國〕相连。此处可确认的是这项被记录的行动；材料未将地点细化到城址，不能用中枢所在地替它补足。

本纪从朝廷视角记录军政行动；出兵与战果需分辨。 记录所见的结果则是：后续路线、控制与损失应与对方记录和遗址材料复核。 日期相邻不等于因果已经证实。

证据的边界是：《元史》为明初仓促官修，本纪保存大量实录条目但有编次与纪年问题；行省执行、数字、族称及对外关系须以条格、多语文书和对方史料互校。 蒙古征服、草原政治与跨区域文献比较研究 可用于继续检验这一结论。

\noindent\eventanchor{SYD-1252-MONGOL-139}{忽必烈及諸王阿里不哥、八里土、出木哈兒、玉龍塔失、昔烈吉、公主脫滅干等來迎，大燕}\textbf{蒙古。}
〔忽必烈及諸王阿里不哥、八里土、出木哈兒、玉龍塔失、昔烈吉、公主脫滅干等來迎，大燕〕见于 《元史》卷003〈本纪第3卷本年条〉 的 1252年（元中统-7年；公元换算据《元史》本纪纪年） 记录。蒙古 的选择因此有了可回查的时间位置；题名保留的城、路、水道或机构名称，是目前可以落地的空间线索。

把本项放回事件链，能够看到的前提为：本纪年条提供日期和中枢可见行动；前因不据单条补定。 而它所打开或收紧的下一步是：可回查同年及后续条目，地方影响仍须另证。

《元史》为明初仓促官修，本纪保存大量实录条目但有编次与纪年问题；行省执行、数字、族称及对外关系须以条格、多语文书和对方史料互校。 因此，蒙古征服、草原政治与跨区域文献比较研究 只能扩展问题，不能代替未保存的细节。

\noindent\eventanchor{SYD-1252-MONGOL-181}{回鶻獻水精盆、珍珠傘等物，可直銀三萬餘錠}\textbf{蒙古。}
1252年（元中统-7年；公元换算据《元史》本纪纪年），《元史》卷003〈本纪第3卷本年条〉 所见的〔回鶻獻水精盆、珍珠傘等物，可直銀三萬餘錠〕把 蒙古 的处置留在可复核的年序中；材料未将地点细化到城址，不能用中枢所在地替它补足。

它承接的条件是：本纪保留灾害或工程的报告地点与朝廷处置；灾情范围需另证。 随后可追踪的变化为：可与地方文书、河工和环境材料比较，不将单点记录外推为全域因果。

关于可说范围，须保留：《元史》为明初仓促官修，本纪保存大量实录条目但有编次与纪年问题；行省执行、数字、族称及对外关系须以条格、多语文书和对方史料互校。 进一步讨论可参照 蒙古征服、草原政治与跨区域文献比较研究

\noindent\eventanchor{SYD-1252-MONGOL-221}{乞都不花等討末來吉兒都怯寨，平之}\textbf{蒙古。}
蒙古 在 1252年（元中统-7年；公元换算据《元史》本纪纪年） 的材料痕迹是〔乞都不花等討末來吉兒都怯寨，平之〕，出处为 《元史》卷003〈本纪第3卷本年条〉。题名保留的城、路、水道或机构名称，是目前可以落地的空间线索。

前一层压力可由以下线索辨认：本纪年条提供日期和中枢可见行动；前因不据单条补定。 这一处置留下的后续条件是：可回查同年及后续条目，地方影响仍须另证。

证据的边界是：《元史》为明初仓促官修，本纪保存大量实录条目但有编次与纪年问题；行省执行、数字、族称及对外关系须以条格、多语文书和对方史料互校。 蒙古征服、草原政治与跨区域文献比较研究 可用于继续检验这一结论。

\noindent\eventanchor{SYD-1252-MONGOL-254}{遣阿藍答兒、脫因、囊加台等詣陝西等處理算錢穀}\textbf{蒙古。}
从 《元史》卷003〈本纪第3卷本年条〉 的 1252年（元中统-7年；公元换算据《元史》本纪纪年） 条目看，〔遣阿藍答兒、脫因、囊加台等詣陝西等處理算錢穀〕使 蒙古 的一个具体行动进入叙事；材料未将地点细化到城址，不能用中枢所在地替它补足。

本纪所记财用、赈济或征发属于特定报告场景。 记录所见的结果则是：可与食货志、文书和物质材料比较，不外推为全境总量。 日期相邻不等于因果已经证实。

《元史》为明初仓促官修，本纪保存大量实录条目但有编次与纪年问题；行省执行、数字、族称及对外关系须以条格、多语文书和对方史料互校。 因此，蒙古征服、草原政治与跨区域文献比较研究 只能扩展问题，不能代替未保存的细节。

\noindent\eventanchor{SYD-1252-MONGOL-286}{賽典赤以為言，帝稍償其直，且禁其勿復有所獻}\textbf{蒙古。}
1252年（元中统-7年；公元换算据《元史》本纪纪年） 的记载将 蒙古 与〔賽典赤以為言，帝稍償其直，且禁其勿復有所獻〕相连。此处可确认的是这项被记录的行动；材料未将地点细化到城址，不能用中枢所在地替它补足。

把本项放回事件链，能够看到的前提为：本纪年条记录政令或机构处置；实施背景须参照制度材料。 而它所打开或收紧的下一步是：可在同年及后续政令中追踪；条文不单独证明地方执行。

关于可说范围，须保留：《元史》为明初仓促官修，本纪保存大量实录条目但有编次与纪年问题；行省执行、数字、族称及对外关系须以条格、多语文书和对方史料互校。 进一步讨论可参照 蒙古征服、草原政治与跨区域文献比较研究

\subsection{1253年（元中统-6年；公元换算据《元史》本纪纪年）}
\noindent\eventanchor{SYD-1253-MONGOL-043}{八年戊午春正月朔，幸也里本朵哈之地，受朝賀}\textbf{蒙古。}
〔八年戊午春正月朔，幸也里本朵哈之地，受朝賀〕见于 《元史》卷003〈本纪第3卷本年条〉 的 1253年（元中统-6年；公元换算据《元史》本纪纪年） 记录。蒙古 的选择因此有了可回查的时间位置；材料未将地点细化到城址，不能用中枢所在地替它补足。

它承接的条件是：本纪年条提供日期和中枢可见行动；前因不据单条补定。 随后可追踪的变化为：可回查同年及后续条目，地方影响仍须另证。

证据的边界是：《元史》为明初仓促官修，本纪保存大量实录条目但有编次与纪年问题；行省执行、数字、族称及对外关系须以条格、多语文书和对方史料互校。 蒙古征服、草原政治与跨区域文献比较研究 可用于继续检验这一结论。

\noindent\eventanchor{SYD-1253-MONGOL-093}{八月辛丑，璮與宋人戰，殺宋師殆盡}\textbf{蒙古与南宋。}
1253年（元中统-6年；公元换算据《元史》本纪纪年），《元史》卷003〈本纪第3卷本年条〉 所见的〔八月辛丑，璮與宋人戰，殺宋師殆盡〕把 蒙古与南宋 的处置留在可复核的年序中；材料未将地点细化到城址，不能用中枢所在地替它补足。

前一层压力可由以下线索辨认：本纪从朝廷视角记录军政行动；出兵与战果需分辨。 这一处置留下的后续条件是：后续路线、控制与损失应与对方记录和遗址材料复核。

《元史》为明初仓促官修，本纪保存大量实录条目但有编次与纪年问题；行省执行、数字、族称及对外关系须以条格、多语文书和对方史料互校。 因此，蒙古征服、草原政治与跨区域文献比较研究 只能扩展问题，不能代替未保存的细节。

\noindent\eventanchor{SYD-1253-MONGOL-140}{薄暮，宋知縣王仲由鵝頂堡出降}\textbf{蒙古与南宋。}
蒙古与南宋 在 1253年（元中统-6年；公元换算据《元史》本纪纪年） 的材料痕迹是〔薄暮，宋知縣王仲由鵝頂堡出降〕，出处为 《元史》卷003〈本纪第3卷本年条〉。材料未将地点细化到城址，不能用中枢所在地替它补足。

本纪从朝廷视角记录军政行动；出兵与战果需分辨。 记录所见的结果则是：后续路线、控制与损失应与对方记录和遗址材料复核。 日期相邻不等于因果已经证实。

关于可说范围，须保留：《元史》为明初仓促官修，本纪保存大量实录条目但有编次与纪年问题；行省执行、数字、族称及对外关系须以条格、多语文书和对方史料互校。 进一步讨论可参照 蒙古征服、草原政治与跨区域文献比较研究

\noindent\eventanchor{SYD-1253-MONGOL-182}{丙辰，進攻〔大〕獲山，守將〔楊〕大淵降，命大淵為四川侍郎，仍以其兵從}\textbf{蒙古。}
从 《元史》卷003〈本纪第3卷本年条〉 的 1253年（元中统-6年；公元换算据《元史》本纪纪年） 条目看，〔丙辰，進攻〔大〕獲山，守將〔楊〕大淵降，命大淵為四川侍郎，仍以其兵從〕使 蒙古 的一个具体行动进入叙事；题名保留的城、路、水道或机构名称，是目前可以落地的空间线索。

把本项放回事件链，能够看到的前提为：本纪年条记录政令或机构处置；实施背景须参照制度材料。 而它所打开或收紧的下一步是：可在同年及后续政令中追踪；条文不单独证明地方执行。

证据的边界是：《元史》为明初仓促官修，本纪保存大量实录条目但有编次与纪年问题；行省执行、数字、族称及对外关系须以条格、多语文书和对方史料互校。 蒙古征服、草原政治与跨区域文献比较研究 可用于继续检验这一结论。

\noindent\eventanchor{SYD-1253-MONGOL-222}{賜田哥玉帶及犒賞士卒，留精兵五百守之}\textbf{蒙古。}
1253年（元中统-6年；公元换算据《元史》本纪纪年） 的记载将 蒙古 与〔賜田哥玉帶及犒賞士卒，留精兵五百守之〕相连。此处可确认的是这项被记录的行动；材料未将地点细化到城址，不能用中枢所在地替它补足。

它承接的条件是：本纪年条提供日期和中枢可见行动；前因不据单条补定。 随后可追踪的变化为：可回查同年及后续条目，地方影响仍须另证。

《元史》为明初仓促官修，本纪保存大量实录条目但有编次与纪年问题；行省执行、数字、族称及对外关系须以条格、多语文书和对方史料互校。 因此，蒙古征服、草原政治与跨区域文献比较研究 只能扩展问题，不能代替未保存的细节。

\noindent\eventanchor{SYD-1253-MONGOL-255}{大淵遣人招降其守將張大悅，仍以大悅為元帥}\textbf{蒙古。}
〔大淵遣人招降其守將張大悅，仍以大悅為元帥〕见于 《元史》卷003〈本纪第3卷本年条〉 的 1253年（元中统-6年；公元换算据《元史》本纪纪年） 记录。蒙古 的选择因此有了可回查的时间位置；材料未将地点细化到城址，不能用中枢所在地替它补足。

前一层压力可由以下线索辨认：本纪从朝廷视角记录军政行动；出兵与战果需分辨。 这一处置留下的后续条件是：后续路线、控制与损失应与对方记录和遗址材料复核。

关于可说范围，须保留：《元史》为明初仓促官修，本纪保存大量实录条目但有编次与纪年问题；行省执行、数字、族称及对外关系须以条格、多语文书和对方史料互校。 进一步讨论可参照 蒙古征服、草原政治与跨区域文献比较研究

\noindent\eventanchor{SYD-1253-MONGOL-287}{帝渡嘉陵江，至白水江，命田哥造浮梁以濟}\textbf{蒙古。}
1253年（元中统-6年；公元换算据《元史》本纪纪年），《元史》卷003〈本纪第3卷本年条〉 所见的〔帝渡嘉陵江，至白水江，命田哥造浮梁以濟〕把 蒙古 的处置留在可复核的年序中；题名保留的城、路、水道或机构名称，是目前可以落地的空间线索。

本纪年条记录政令或机构处置；实施背景须参照制度材料。 记录所见的结果则是：可在同年及后续政令中追踪；条文不单独证明地方执行。 日期相邻不等于因果已经证实。

证据的边界是：《元史》为明初仓促官修，本纪保存大量实录条目但有编次与纪年问题；行省执行、数字、族称及对外关系须以条格、多语文书和对方史料互校。 蒙古征服、草原政治与跨区域文献比较研究 可用于继续检验这一结论。

\subsection{1253年}
\noindent\eventanchor{YUAN-1253-EVENT-165}{蒙古军征服大理国}\textbf{蒙古与大理。}
蒙古与大理 在 1253年 的材料痕迹是〔蒙古军征服大理国〕，出处为 《元史》相关本纪；《元朝秘史》、波斯文史籍或《高丽史》对应叙事；《元典章》《通制条格》、多语碑刻与文书（据事核）。材料未将地点细化到城址，不能用中枢所在地替它补足。

把本项放回事件链，能够看到的前提为：“蒙古军征服大理国”发生在窝阔台至蒙哥时期的扩张与继承的具体压力与机会之中，相关行动者的选择不能脱离此前事件链解释。 而它所打开或收紧的下一步是：“蒙古军征服大理国”为后续政治、边防或资源配置留下条件；其社会影响与实际覆盖范围仍须以异类材料限定。

《元史》明初速修，纪年与人名有异同；蒙古大汗政权不等同所有汗国，征服、赋役和身份分类须按地区及材料语言分列。 因此，蒙古帝国、元朝行省财政、多语文书与区域社会研究 只能扩展问题，不能代替未保存的细节。

\subsection{1254 年（共享年序桥接）}
\noindent\eventanchor{SY-1254-CHRONOLOGY-BRIDGE}{【C级年序桥接】保留1254 年的多政权并行检索入口；本条不主张所列政权在当年发生直接互动，具体事件须另以单项材料入账。}\textbf{南宋、金、西夏与蒙古（年序桥接）。}
从 《宋史》；《金史》；《元史》；蒙古文、波斯文及地方材料（据事核） 的 1254 年（共享年序桥接） 条目看，〔【C级年序桥接】保留1254 年的多政权并行检索入口；本条不主张所列政权在当年发生直接互动，具体事件须另以单项材料入账。〕使 南宋、金、西夏与蒙古（年序桥接） 的一个具体行动进入叙事；题名保留的城、路、水道或机构名称，是目前可以落地的空间线索。

它承接的条件是：相邻已入账事件之间缺少该年度的共享年序入口，易使跨政权材料被单一政权叙事遮蔽。 随后可追踪的变化为：保留该年度的检索与补账位置；后续仅在获得可核单项材料时拆分为具体政治、财政、军事、社会或文化事件。

关于可说范围，须保留：本条只确认该年位于可追踪的共享年序中；正史、地方文书和外部材料的保存密度不一，不能由桥接标签推断行动、统属、疆界或社会影响。 进一步讨论可参照 蒙古征服、区域政权转换与跨语种年序研究

\subsection{1254年（元中统-5年；公元换算据《元史》本纪纪年）}
\noindent\eventanchor{SYD-1254-MONGOL-044}{遲明，遇雨，梯折，後軍不克進而止}\textbf{蒙古。}
1254年（元中统-5年；公元换算据《元史》本纪纪年） 的记载将 蒙古 与〔遲明，遇雨，梯折，後軍不克進而止〕相连。此处可确认的是这项被记录的行动；材料未将地点细化到城址，不能用中枢所在地替它补足。

前一层压力可由以下线索辨认：本纪保留灾害或工程的报告地点与朝廷处置；灾情范围需另证。 这一处置留下的后续条件是：可与地方文书、河工和环境材料比较，不将单点记录外推为全域因果。

证据的边界是：《元史》为明初仓促官修，本纪保存大量实录条目但有编次与纪年问题；行省执行、数字、族称及对外关系须以条格、多语文书和对方史料互校。 蒙古征服、草原政治与跨区域文献比较研究 可用于继续检验这一结论。

\noindent\eventanchor{SYD-1254-MONGOL-094}{初，太宗朝，羣臣擅權，政出多門}\textbf{蒙古。}
〔初，太宗朝，羣臣擅權，政出多門〕见于 《元史》卷003〈本纪第3卷本年条〉 的 1254年（元中统-5年；公元换算据《元史》本纪纪年） 记录。蒙古 的选择因此有了可回查的时间位置；材料未将地点细化到城址，不能用中枢所在地替它补足。

本纪年条提供日期和中枢可见行动；前因不据单条补定。 记录所见的结果则是：可回查同年及后续条目，地方影响仍须另证。 日期相邻不等于因果已经证实。

《元史》为明初仓促官修，本纪保存大量实录条目但有编次与纪年问题；行省执行、数字、族称及对外关系须以条格、多语文书和对方史料互校。 因此，蒙古征服、草原政治与跨区域文献比较研究 只能扩展问题，不能代替未保存的细节。

\noindent\eventanchor{SYD-1254-MONGOL-141}{丁卯，大淵請攻合州，俘男女八萬餘}\textbf{蒙古。}
1254年（元中统-5年；公元换算据《元史》本纪纪年），《元史》卷003〈本纪第3卷本年条〉 所见的〔丁卯，大淵請攻合州，俘男女八萬餘〕把 蒙古 的处置留在可复核的年序中；题名保留的城、路、水道或机构名称，是目前可以落地的空间线索。

把本项放回事件链，能够看到的前提为：本纪从朝廷视角记录军政行动；出兵与战果需分辨。 而它所打开或收紧的下一步是：后续路线、控制与损失应与对方记录和遗址材料复核。

关于可说范围，须保留：《元史》为明初仓促官修，本纪保存大量实录条目但有编次与纪年问题；行省执行、数字、族称及对外关系须以条格、多语文书和对方史料互校。 进一步讨论可参照 蒙古征服、草原政治与跨区域文献比较研究

\noindent\eventanchor{SYD-1254-MONGOL-183}{丁酉夜，登外城，殺宋兵甚眾}\textbf{蒙古与南宋。}
蒙古与南宋 在 1254年（元中统-5年；公元换算据《元史》本纪纪年） 的材料痕迹是〔丁酉夜，登外城，殺宋兵甚眾〕，出处为 《元史》卷003〈本纪第3卷本年条〉。题名保留的城、路、水道或机构名称，是目前可以落地的空间线索。

它承接的条件是：本纪从朝廷视角记录军政行动；出兵与战果需分辨。 随后可追踪的变化为：后续路线、控制与损失应与对方记录和遗址材料复核。

证据的边界是：《元史》为明初仓促官修，本纪保存大量实录条目但有编次与纪年问题；行省执行、数字、族称及对外关系须以条格、多语文书和对方史料互校。 蒙古征服、草原政治与跨区域文献比较研究 可用于继续检验这一结论。

\noindent\eventanchor{SYD-1254-MONGOL-223}{九年己未春正月乙巳朔，駐蹕重貴山北，置酒大會，因問諸王、駙馬、百官曰：「今在宋境，夏暑且至，汝等其謂可居否乎？」札剌亦兒部人脫歡曰：「南土瘴癘，上宜北還}\textbf{蒙古与南宋。}
从 《元史》卷003〈本纪第3卷本年条〉 的 1254年（元中统-5年；公元换算据《元史》本纪纪年） 条目看，〔九年己未春正月乙巳朔，駐蹕重貴山北，置酒大會，因問諸王、駙馬、百官曰：「今在宋境，夏暑且至，汝等其謂可居否乎？」札剌亦兒部人脫歡曰：「南土瘴癘，上宜北還〕使 蒙古与南宋 的一个具体行动进入叙事；题名保留的城、路、水道或机构名称，是目前可以落地的空间线索。

前一层压力可由以下线索辨认：本纪年条记录政令或机构处置；实施背景须参照制度材料。 这一处置留下的后续条件是：可在同年及后续政令中追踪；条文不单独证明地方执行。

《元史》为明初仓促官修，本纪保存大量实录条目但有编次与纪年问题；行省执行、数字、族称及对外关系须以条格、多语文书和对方史料互校。 因此，蒙古征服、草原政治与跨区域文献比较研究 只能扩展问题，不能代替未保存的细节。

\noindent\eventanchor{SYD-1254-MONGOL-256}{六月丁巳，汪田哥復選兵夜登外城馬軍寨，殺寨主及守城者}\textbf{蒙古。}
1254年（元中统-5年；公元换算据《元史》本纪纪年） 的记载将 蒙古 与〔六月丁巳，汪田哥復選兵夜登外城馬軍寨，殺寨主及守城者〕相连。此处可确认的是这项被记录的行动；题名保留的城、路、水道或机构名称，是目前可以落地的空间线索。

本纪年条记录政令或机构处置；实施背景须参照制度材料。 记录所见的结果则是：可在同年及后续政令中追踪；条文不单独证明地方执行。 日期相邻不等于因果已经证实。

关于可说范围，须保留：《元史》为明初仓促官修，本纪保存大量实录条目但有编次与纪年问题；行省执行、数字、族称及对外关系须以条格、多语文书和对方史料互校。 进一步讨论可参照 蒙古征服、草原政治与跨区域文献比较研究

\noindent\eventanchor{SYD-1254-MONGOL-288}{秋七月辛亥，留精兵三千守之，餘悉攻重慶}\textbf{蒙古。}
〔秋七月辛亥，留精兵三千守之，餘悉攻重慶〕见于 《元史》卷003〈本纪第3卷本年条〉 的 1254年（元中统-5年；公元换算据《元史》本纪纪年） 记录。蒙古 的选择因此有了可回查的时间位置；材料未将地点细化到城址，不能用中枢所在地替它补足。

把本项放回事件链，能够看到的前提为：本纪从朝廷视角记录军政行动；出兵与战果需分辨。 而它所打开或收紧的下一步是：后续路线、控制与损失应与对方记录和遗址材料复核。

证据的边界是：《元史》为明初仓促官修，本纪保存大量实录条目但有编次与纪年问题；行省执行、数字、族称及对外关系须以条格、多语文书和对方史料互校。 蒙古征服、草原政治与跨区域文献比较研究 可用于继续检验这一结论。

\subsection{1255 年（共享年序桥接）}
\noindent\eventanchor{SY-1255-CHRONOLOGY-BRIDGE}{【C级年序桥接】保留1255 年的多政权并行检索入口；本条不主张所列政权在当年发生直接互动，具体事件须另以单项材料入账。}\textbf{南宋、金、西夏与蒙古（年序桥接）。}
1255 年（共享年序桥接），《宋史》；《金史》；《元史》；蒙古文、波斯文及地方材料（据事核） 所见的〔【C级年序桥接】保留1255 年的多政权并行检索入口；本条不主张所列政权在当年发生直接互动，具体事件须另以单项材料入账。〕把 南宋、金、西夏与蒙古（年序桥接） 的处置留在可复核的年序中；题名保留的城、路、水道或机构名称，是目前可以落地的空间线索。

它承接的条件是：相邻已入账事件之间缺少该年度的共享年序入口，易使跨政权材料被单一政权叙事遮蔽。 随后可追踪的变化为：保留该年度的检索与补账位置；后续仅在获得可核单项材料时拆分为具体政治、财政、军事、社会或文化事件。

本条只确认该年位于可追踪的共享年序中；正史、地方文书和外部材料的保存密度不一，不能由桥接标签推断行动、统属、疆界或社会影响。 因此，蒙古征服、区域政权转换与跨语种年序研究 只能扩展问题，不能代替未保存的细节。

\subsection{1256年}
\noindent\eventanchor{YUAN-1256-EVENT-166}{蒙哥朝加强汉地户籍、赋税和征发整顿}\textbf{蒙古。}
蒙古 在 1256年 的材料痕迹是〔蒙哥朝加强汉地户籍、赋税和征发整顿〕，出处为 《元史》相关本纪；《元朝秘史》、波斯文史籍或《高丽史》对应叙事；《元典章》《通制条格》、多语碑刻与文书（据事核）。材料未将地点细化到城址，不能用中枢所在地替它补足。

前一层压力可由以下线索辨认：窝阔台至蒙哥时期的扩张与继承对中枢资源、交通或行政协同的要求，推动了“蒙哥朝加强汉地户籍、赋税和征发整顿”这一制度或空间安排。 这一处置留下的后续条件是：“蒙哥朝加强汉地户籍、赋税和征发整顿”提供新的治理框架或资源通道，但颁行、复申与不同地区的实际执行须分开核验。

关于可说范围，须保留：《元史》明初速修，纪年与人名有异同；蒙古大汗政权不等同所有汗国，征服、赋役和身份分类须按地区及材料语言分列。 进一步讨论可参照 蒙古帝国、元朝行省财政、多语文书与区域社会研究

\subsection{1257 年（共享年序桥接）}
\noindent\eventanchor{SY-1257-CHRONOLOGY-BRIDGE}{【C级年序桥接】保留1257 年的多政权并行检索入口；本条不主张所列政权在当年发生直接互动，具体事件须另以单项材料入账。}\textbf{南宋、金、西夏与蒙古（年序桥接）。}
从 《宋史》；《金史》；《元史》；蒙古文、波斯文及地方材料（据事核） 的 1257 年（共享年序桥接） 条目看，〔【C级年序桥接】保留1257 年的多政权并行检索入口；本条不主张所列政权在当年发生直接互动，具体事件须另以单项材料入账。〕使 南宋、金、西夏与蒙古（年序桥接） 的一个具体行动进入叙事；题名保留的城、路、水道或机构名称，是目前可以落地的空间线索。

相邻已入账事件之间缺少该年度的共享年序入口，易使跨政权材料被单一政权叙事遮蔽。 记录所见的结果则是：保留该年度的检索与补账位置；后续仅在获得可核单项材料时拆分为具体政治、财政、军事、社会或文化事件。 日期相邻不等于因果已经证实。

证据的边界是：本条只确认该年位于可追踪的共享年序中；正史、地方文书和外部材料的保存密度不一，不能由桥接标签推断行动、统属、疆界或社会影响。 蒙古征服、区域政权转换与跨语种年序研究 可用于继续检验这一结论。

\subsection{1259年}
\noindent\eventanchor{YUAN-1259-EVENT-167}{蒙哥在攻宋期间去世}\textbf{蒙古与南宋。}
1259年 的记载将 蒙古与南宋 与〔蒙哥在攻宋期间去世〕相连。此处可确认的是这项被记录的行动；材料未将地点细化到城址，不能用中枢所在地替它补足。

把本项放回事件链，能够看到的前提为：窝阔台至蒙哥时期的扩张与继承中前任统治的结束与宫廷、部族或官僚的权力安排相互交织，促成“蒙哥在攻宋期间去世”这一继承节点。 而它所打开或收紧的下一步是：继承并不自动消除既有矛盾；“蒙哥在攻宋期间去世”后的任官、军权和对外策略需由后续条目追踪。

《元史》明初速修，纪年与人名有异同；蒙古大汗政权不等同所有汗国，征服、赋役和身份分类须按地区及材料语言分列。 因此，蒙古帝国、元朝行省财政、多语文书与区域社会研究 只能扩展问题，不能代替未保存的细节。

\subsection{1260年}
\noindent\eventanchor{YUAN-1260-EVENT-168}{忽必烈即位，与阿里不哥展开争位冲突}\textbf{蒙古。}
〔忽必烈即位，与阿里不哥展开争位冲突〕见于 《元史》相关本纪；《元朝秘史》、波斯文史籍或《高丽史》对应叙事；《元典章》《通制条格》、多语碑刻与文书（据事核） 的 1260年 记录。蒙古 的选择因此有了可回查的时间位置；材料未将地点细化到城址，不能用中枢所在地替它补足。

它承接的条件是：窝阔台至蒙哥时期的扩张与继承中前任统治的结束与宫廷、部族或官僚的权力安排相互交织，促成“忽必烈即位，与阿里不哥展开争位冲突”这一继承节点。 随后可追踪的变化为：继承并不自动消除既有矛盾；“忽必烈即位，与阿里不哥展开争位冲突”后的任官、军权和对外策略需由后续条目追踪。

关于可说范围，须保留：《元史》明初速修，纪年与人名有异同；蒙古大汗政权不等同所有汗国，征服、赋役和身份分类须按地区及材料语言分列。 进一步讨论可参照 蒙古帝国、元朝行省财政、多语文书与区域社会研究

\subsection{1261年}
\noindent\eventanchor{YUAN-1261-EVENT-169}{李璮在山东反叛后被镇压}\textbf{蒙古。}
1261年，《元史》相关本纪；《元朝秘史》、波斯文史籍或《高丽史》对应叙事；《元典章》《通制条格》、多语碑刻与文书（据事核） 所见的〔李璮在山东反叛后被镇压〕把 蒙古 的处置留在可复核的年序中；题名保留的城、路、水道或机构名称，是目前可以落地的空间线索。

前一层压力可由以下线索辨认：忽必烈政权、征宋和元朝建制中既有的联盟、交通与补给压力，为“李璮在山东反叛后被镇压”提供了直接的军事或动员背景。 这一处置留下的后续条件是：“李璮在山东反叛后被镇压”改变了相关城镇、道路或政权间的军事选择；伤亡、俘获和控制范围仅可按各方材料分别确认。

证据的边界是：《元史》明初速修，纪年与人名有异同；蒙古大汗政权不等同所有汗国，征服、赋役和身份分类须按地区及材料语言分列。 蒙古帝国、元朝行省财政、多语文书与区域社会研究 可用于继续检验这一结论。

\subsection{1262 年（共享年序桥接）}
\noindent\eventanchor{SY-1262-CHRONOLOGY-BRIDGE}{【C级年序桥接】保留1262 年的多政权并行检索入口；本条不主张所列政权在当年发生直接互动，具体事件须另以单项材料入账。}\textbf{南宋、金、西夏与蒙古（年序桥接）。}
南宋、金、西夏与蒙古（年序桥接） 在 1262 年（共享年序桥接） 的材料痕迹是〔【C级年序桥接】保留1262 年的多政权并行检索入口；本条不主张所列政权在当年发生直接互动，具体事件须另以单项材料入账。〕，出处为 《宋史》；《金史》；《元史》；蒙古文、波斯文及地方材料（据事核）。题名保留的城、路、水道或机构名称，是目前可以落地的空间线索。

相邻已入账事件之间缺少该年度的共享年序入口，易使跨政权材料被单一政权叙事遮蔽。 记录所见的结果则是：保留该年度的检索与补账位置；后续仅在获得可核单项材料时拆分为具体政治、财政、军事、社会或文化事件。 日期相邻不等于因果已经证实。

本条只确认该年位于可追踪的共享年序中；正史、地方文书和外部材料的保存密度不一，不能由桥接标签推断行动、统属、疆界或社会影响。 因此，蒙古征服、区域政权转换与跨语种年序研究 只能扩展问题，不能代替未保存的细节。

\subsection{1262年（元中统3年；公元换算据《元史》本纪纪年）}
\noindent\eventanchor{SYD-1262-MONGOL-045}{丙午，詔特徵人員，宜令乘傳}\textbf{蒙古。}
从 《元史》卷005〈本纪第5卷本年条〉 的 1262年（元中统3年；公元换算据《元史》本纪纪年） 条目看，〔丙午，詔特徵人員，宜令乘傳〕使 蒙古 的一个具体行动进入叙事；材料未将地点细化到城址，不能用中枢所在地替它补足。

把本项放回事件链，能够看到的前提为：本纪年条记录政令或机构处置；实施背景须参照制度材料。 而它所打开或收紧的下一步是：可在同年及后续政令中追踪；条文不单独证明地方执行。

关于可说范围，须保留：《元史》为明初仓促官修，本纪保存大量实录条目但有编次与纪年问题；行省执行、数字、族称及对外关系须以条格、多语文书和对方史料互校。 进一步讨论可参照 蒙古征服、草原政治与跨区域文献比较研究

\noindent\eventanchor{SYD-1262-MONGOL-095}{壬申，遣使收輯諸路軍民官海青牌及驛券}\textbf{蒙古。}
1262年（元中统3年；公元换算据《元史》本纪纪年） 的记载将 蒙古 与〔壬申，遣使收輯諸路軍民官海青牌及驛券〕相连。此处可确认的是这项被记录的行动；题名保留的城、路、水道或机构名称，是目前可以落地的空间线索。

它承接的条件是：本纪记录使节、贡赐或往来，不等于稳定统属关系。 随后可追踪的变化为：可与对方编年和交通、文书材料核对其后续落实。

证据的边界是：《元史》为明初仓促官修，本纪保存大量实录条目但有编次与纪年问题；行省执行、数字、族称及对外关系须以条格、多语文书和对方史料互校。 蒙古征服、草原政治与跨区域文献比较研究 可用于继续检验这一结论。

\noindent\eventanchor{SYD-1262-MONGOL-142}{乙丑，復立息州城以安其民}\textbf{蒙古。}
〔乙丑，復立息州城以安其民〕见于 《元史》卷005〈本纪第5卷本年条〉 的 1262年（元中统3年；公元换算据《元史》本纪纪年） 记录。蒙古 的选择因此有了可回查的时间位置；题名保留的城、路、水道或机构名称，是目前可以落地的空间线索。

前一层压力可由以下线索辨认：本纪年条记录政令或机构处置；实施背景须参照制度材料。 这一处置留下的后续条件是：可在同年及后续政令中追踪；条文不单独证明地方执行。

《元史》为明初仓促官修，本纪保存大量实录条目但有编次与纪年问题；行省执行、数字、族称及对外关系须以条格、多语文书和对方史料互校。 因此，蒙古征服、草原政治与跨区域文献比较研究 只能扩展问题，不能代替未保存的细节。

\noindent\eventanchor{SYD-1262-MONGOL-184}{詔安輯徐、邳民，禁征戍軍士及勢官，毋縱畜牧傷其禾稼桑棗}\textbf{蒙古。}
1262年（元中统3年；公元换算据《元史》本纪纪年），《元史》卷005〈本纪第5卷本年条〉 所见的〔詔安輯徐、邳民，禁征戍軍士及勢官，毋縱畜牧傷其禾稼桑棗〕把 蒙古 的处置留在可复核的年序中；材料未将地点细化到城址，不能用中枢所在地替它补足。

本纪年条记录政令或机构处置；实施背景须参照制度材料。 记录所见的结果则是：可在同年及后续政令中追踪；条文不单独证明地方执行。 日期相邻不等于因果已经证实。

关于可说范围，须保留：《元史》为明初仓促官修，本纪保存大量实录条目但有编次与纪年问题；行省执行、数字、族称及对外关系须以条格、多语文书和对方史料互校。 进一步讨论可参照 蒙古征服、草原政治与跨区域文献比较研究

\noindent\eventanchor{SYD-1262-MONGOL-224}{作佛事於昊天寺七晝夜，賜銀萬五千兩}\textbf{蒙古。}
蒙古 在 1262年（元中统3年；公元换算据《元史》本纪纪年） 的材料痕迹是〔作佛事於昊天寺七晝夜，賜銀萬五千兩〕，出处为 《元史》卷005〈本纪第5卷本年条〉。题名保留的城、路、水道或机构名称，是目前可以落地的空间线索。

把本项放回事件链，能够看到的前提为：本纪年条提供日期和中枢可见行动；前因不据单条补定。 而它所打开或收紧的下一步是：可回查同年及后续条目，地方影响仍须另证。

证据的边界是：《元史》为明初仓促官修，本纪保存大量实录条目但有编次与纪年问题；行省执行、数字、族称及对外关系须以条格、多语文书和对方史料互校。 蒙古征服、草原政治与跨区域文献比较研究 可用于继续检验这一结论。

\noindent\eventanchor{SYD-1262-MONGOL-257}{安南國陳光昞遣使貢方物}\textbf{蒙古。}
从 《元史》卷005〈本纪第5卷本年条〉 的 1262年（元中统3年；公元换算据《元史》本纪纪年） 条目看，〔安南國陳光昞遣使貢方物〕使 蒙古 的一个具体行动进入叙事；材料未将地点细化到城址，不能用中枢所在地替它补足。

它承接的条件是：本纪记录使节、贡赐或往来，不等于稳定统属关系。 随后可追踪的变化为：可与对方编年和交通、文书材料核对其后续落实。

《元史》为明初仓促官修，本纪保存大量实录条目但有编次与纪年问题；行省执行、数字、族称及对外关系须以条格、多语文书和对方史料互校。 因此，蒙古征服、草原政治与跨区域文献比较研究 只能扩展问题，不能代替未保存的细节。

\noindent\eventanchor{SYD-1262-MONGOL-289}{八月己丑，郭守敬請開玉泉水以通漕運}\textbf{蒙古。}
1262年（元中统3年；公元换算据《元史》本纪纪年） 的记载将 蒙古 与〔八月己丑，郭守敬請開玉泉水以通漕運〕相连。此处可确认的是这项被记录的行动；材料未将地点细化到城址，不能用中枢所在地替它补足。

前一层压力可由以下线索辨认：本纪保留灾害或工程的报告地点与朝廷处置；灾情范围需另证。 这一处置留下的后续条件是：可与地方文书、河工和环境材料比较，不将单点记录外推为全域因果。

关于可说范围，须保留：《元史》为明初仓促官修，本纪保存大量实录条目但有编次与纪年问题；行省执行、数字、族称及对外关系须以条格、多语文书和对方史料互校。 进一步讨论可参照 蒙古征服、草原政治与跨区域文献比较研究

\subsection{1263 年（共享年序桥接）}
\noindent\eventanchor{SY-1263-CHRONOLOGY-BRIDGE}{【C级年序桥接】保留1263 年的多政权并行检索入口；本条不主张所列政权在当年发生直接互动，具体事件须另以单项材料入账。}\textbf{南宋、金、西夏与蒙古（年序桥接）。}
〔【C级年序桥接】保留1263 年的多政权并行检索入口；本条不主张所列政权在当年发生直接互动，具体事件须另以单项材料入账。〕见于 《宋史》；《金史》；《元史》；蒙古文、波斯文及地方材料（据事核） 的 1263 年（共享年序桥接） 记录。南宋、金、西夏与蒙古（年序桥接） 的选择因此有了可回查的时间位置；题名保留的城、路、水道或机构名称，是目前可以落地的空间线索。

相邻已入账事件之间缺少该年度的共享年序入口，易使跨政权材料被单一政权叙事遮蔽。 记录所见的结果则是：保留该年度的检索与补账位置；后续仅在获得可核单项材料时拆分为具体政治、财政、军事、社会或文化事件。 日期相邻不等于因果已经证实。

证据的边界是：本条只确认该年位于可追踪的共享年序中；正史、地方文书和外部材料的保存密度不一，不能由桥接标签推断行动、统属、疆界或社会影响。 蒙古征服、区域政权转换与跨语种年序研究 可用于继续检验这一结论。

\subsection{1263年（元中统4年；公元换算据《元史》本纪纪年）}
\noindent\eventanchor{SYD-1263-MONGOL-046}{:比者星芒示儆，雨澤愆常，皆闕政之所繇，顧斯民之何罪}\textbf{蒙古。}
1263年（元中统4年；公元换算据《元史》本纪纪年），《元史》卷005〈本纪第5卷本年条〉 所见的〔:比者星芒示儆，雨澤愆常，皆闕政之所繇，顧斯民之何罪〕把 蒙古 的处置留在可复核的年序中；材料未将地点细化到城址，不能用中枢所在地替它补足。

把本项放回事件链，能够看到的前提为：本纪保留灾害或工程的报告地点与朝廷处置；灾情范围需另证。 而它所打开或收紧的下一步是：可与地方文书、河工和环境材料比较，不将单点记录外推为全域因果。

《元史》为明初仓促官修，本纪保存大量实录条目但有编次与纪年问题；行省执行、数字、族称及对外关系须以条格、多语文书和对方史料互校。 因此，蒙古征服、草原政治与跨区域文献比较研究 只能扩展问题，不能代替未保存的细节。

\noindent\eventanchor{SYD-1263-MONGOL-096}{庚戌，命燕王署敕、諸王設僚屬及說書官}\textbf{蒙古。}
蒙古 在 1263年（元中统4年；公元换算据《元史》本纪纪年） 的材料痕迹是〔庚戌，命燕王署敕、諸王設僚屬及說書官〕，出处为 《元史》卷005〈本纪第5卷本年条〉。题名保留的城、路、水道或机构名称，是目前可以落地的空间线索。

它承接的条件是：本纪年条记录政令或机构处置；实施背景须参照制度材料。 随后可追踪的变化为：可在同年及后续政令中追踪；条文不单独证明地方执行。

关于可说范围，须保留：《元史》为明初仓促官修，本纪保存大量实录条目但有编次与纪年问题；行省执行、数字、族称及对外关系须以条格、多语文书和对方史料互校。 进一步讨论可参照 蒙古征服、草原政治与跨区域文献比较研究

\noindent\eventanchor{SYD-1263-MONGOL-143}{〔己〕丑，以曆日賜高麗國王王禃}\textbf{蒙古。}
从 《元史》卷005〈本纪第5卷本年条〉 的 1263年（元中统4年；公元换算据《元史》本纪纪年） 条目看，〔〔己〕丑，以曆日賜高麗國王王禃〕使 蒙古 的一个具体行动进入叙事；材料未将地点细化到城址，不能用中枢所在地替它补足。

前一层压力可由以下线索辨认：本纪年条提供日期和中枢可见行动；前因不据单条补定。 这一处置留下的后续条件是：可回查同年及后续条目，地方影响仍须另证。

证据的边界是：《元史》为明初仓促官修，本纪保存大量实录条目但有编次与纪年问题；行省执行、数字、族称及对外关系须以条格、多语文书和对方史料互校。 蒙古征服、草原政治与跨区域文献比较研究 可用于继续检验这一结论。

\noindent\eventanchor{SYD-1263-MONGOL-185}{〔乙〕卯，詔高麗國王王植來朝上都，修世見之禮}\textbf{蒙古。}
1263年（元中统4年；公元换算据《元史》本纪纪年） 的记载将 蒙古 与〔〔乙〕卯，詔高麗國王王植來朝上都，修世見之禮〕相连。此处可确认的是这项被记录的行动；题名保留的城、路、水道或机构名称，是目前可以落地的空间线索。

本纪年条记录政令或机构处置；实施背景须参照制度材料。 记录所见的结果则是：可在同年及后续政令中追踪；条文不单独证明地方执行。 日期相邻不等于因果已经证实。

《元史》为明初仓促官修，本纪保存大量实录条目但有编次与纪年问题；行省执行、数字、族称及对外关系须以条格、多语文书和对方史料互校。 因此，蒙古征服、草原政治与跨区域文献比较研究 只能扩展问题，不能代替未保存的细节。

\noindent\eventanchor{SYD-1263-MONGOL-225}{八月壬寅朔，陝西行省臣上言：「川蜀戍兵軍需，請令奧魯官徵入官庫，移文於近戍官司，依數取之}\textbf{蒙古。}
〔八月壬寅朔，陝西行省臣上言：「川蜀戍兵軍需，請令奧魯官徵入官庫，移文於近戍官司，依數取之〕见于 《元史》卷005〈本纪第5卷本年条〉 的 1263年（元中统4年；公元换算据《元史》本纪纪年） 记录。蒙古 的选择因此有了可回查的时间位置；题名保留的城、路、水道或机构名称，是目前可以落地的空间线索。

把本项放回事件链，能够看到的前提为：本纪年条记录政令或机构处置；实施背景须参照制度材料。 而它所打开或收紧的下一步是：可在同年及后续政令中追踪；条文不单独证明地方执行。

关于可说范围，须保留：《元史》为明初仓促官修，本纪保存大量实录条目但有编次与纪年问题；行省执行、数字、族称及对外关系须以条格、多语文书和对方史料互校。 进一步讨论可参照 蒙古征服、草原政治与跨区域文献比较研究

\noindent\eventanchor{SYD-1263-MONGOL-258}{八月戊申朔，詔霍木海總管諸路驛，佩金符}\textbf{蒙古。}
1263年（元中统4年；公元换算据《元史》本纪纪年），《元史》卷005〈本纪第5卷本年条〉 所见的〔八月戊申朔，詔霍木海總管諸路驛，佩金符〕把 蒙古 的处置留在可复核的年序中；题名保留的城、路、水道或机构名称，是目前可以落地的空间线索。

它承接的条件是：本纪年条记录政令或机构处置；实施背景须参照制度材料。 随后可追踪的变化为：可在同年及后续政令中追踪；条文不单独证明地方执行。

证据的边界是：《元史》为明初仓促官修，本纪保存大量实录条目但有编次与纪年问题；行省执行、数字、族称及对外关系须以条格、多语文书和对方史料互校。 蒙古征服、草原政治与跨区域文献比较研究 可用于继续检验这一结论。

\noindent\eventanchor{SYD-1263-MONGOL-290}{濱、棣二州蝗，真定路旱}\textbf{蒙古。}
蒙古 在 1263年（元中统4年；公元换算据《元史》本纪纪年） 的材料痕迹是〔濱、棣二州蝗，真定路旱〕，出处为 《元史》卷005〈本纪第5卷本年条〉。题名保留的城、路、水道或机构名称，是目前可以落地的空间线索。

前一层压力可由以下线索辨认：本纪保留灾害或工程的报告地点与朝廷处置；灾情范围需另证。 这一处置留下的后续条件是：可与地方文书、河工和环境材料比较，不将单点记录外推为全域因果。

《元史》为明初仓促官修，本纪保存大量实录条目但有编次与纪年问题；行省执行、数字、族称及对外关系须以条格、多语文书和对方史料互校。 因此，蒙古征服、草原政治与跨区域文献比较研究 只能扩展问题，不能代替未保存的细节。

\subsection{1264年（元至元1年；公元换算据《元史》本纪纪年）}
\noindent\eventanchor{SYD-1264-MONGOL-047}{:朕祗紹天明，入纂丕緒，于今三年，夙夜寅畏，罔敢怠荒}\textbf{蒙古。}
从 《元史》卷038〈至元元年〉 的 1264年（元至元1年；公元换算据《元史》本纪纪年） 条目看，〔:朕祗紹天明，入纂丕緒，于今三年，夙夜寅畏，罔敢怠荒〕使 蒙古 的一个具体行动进入叙事；材料未将地点细化到城址，不能用中枢所在地替它补足。

本纪年条提供日期和中枢可见行动；前因不据单条补定。 记录所见的结果则是：可回查同年及后续条目，地方影响仍须另证。 日期相邻不等于因果已经证实。

关于可说范围，须保留：《元史》为明初仓促官修，本纪保存大量实录条目但有编次与纪年问题；行省执行、数字、族称及对外关系须以条格、多语文书和对方史料互校。 进一步讨论可参照 蒙古征服、草原政治与跨区域文献比较研究

\noindent\eventanchor{SYD-1264-MONGOL-097}{命宣政院使末吉以司徒就第}\textbf{蒙古。}
1264年（元至元1年；公元换算据《元史》本纪纪年） 的记载将 蒙古 与〔命宣政院使末吉以司徒就第〕相连。此处可确认的是这项被记录的行动；材料未将地点细化到城址，不能用中枢所在地替它补足。

把本项放回事件链，能够看到的前提为：本纪年条记录政令或机构处置；实施背景须参照制度材料。 而它所打开或收紧的下一步是：可在同年及后续政令中追踪；条文不单独证明地方执行。

证据的边界是：《元史》为明初仓促官修，本纪保存大量实录条目但有编次与纪年问题；行省执行、数字、族称及对外关系须以条格、多语文书和对方史料互校。 蒙古征服、草原政治与跨区域文献比较研究 可用于继续检验这一结论。

\noindent\eventanchor{SYD-1264-MONGOL-144}{是月西和州、徽州雨雹，民饑，發米賑貸之}\textbf{蒙古。}
〔是月西和州、徽州雨雹，民饑，發米賑貸之〕见于 《元史》卷038〈至元元年〉 的 1264年（元至元1年；公元换算据《元史》本纪纪年） 记录。蒙古 的选择因此有了可回查的时间位置；题名保留的城、路、水道或机构名称，是目前可以落地的空间线索。

它承接的条件是：本纪所记财用、赈济或征发属于特定报告场景。 随后可追踪的变化为：可与食货志、文书和物质材料比较，不外推为全境总量。

《元史》为明初仓促官修，本纪保存大量实录条目但有编次与纪年问题；行省执行、数字、族称及对外关系须以条格、多语文书和对方史料互校。 因此，蒙古征服、草原政治与跨区域文献比较研究 只能扩展问题，不能代替未保存的细节。

\noindent\eventanchor{SYD-1264-MONGOL-186}{丙辰，以大司農塔失海牙為太尉，置僚屬，商議中書省事}\textbf{蒙古。}
1264年（元至元1年；公元换算据《元史》本纪纪年），《元史》卷038〈至元元年〉 所见的〔丙辰，以大司農塔失海牙為太尉，置僚屬，商議中書省事〕把 蒙古 的处置留在可复核的年序中；题名保留的城、路、水道或机构名称，是目前可以落地的空间线索。

前一层压力可由以下线索辨认：本纪年条记录政令或机构处置；实施背景须参照制度材料。 这一处置留下的后续条件是：可在同年及后续政令中追踪；条文不单独证明地方执行。

关于可说范围，须保留：《元史》为明初仓促官修，本纪保存大量实录条目但有编次与纪年问题；行省执行、数字、族称及对外关系须以条格、多语文书和对方史料互校。 进一步讨论可参照 蒙古征服、草原政治与跨区域文献比较研究

\noindent\eventanchor{SYD-1264-MONGOL-226}{丙辰，制省諸王、公主、駙馬飲饍之費}\textbf{蒙古。}
蒙古 在 1264年（元至元1年；公元换算据《元史》本纪纪年） 的材料痕迹是〔丙辰，制省諸王、公主、駙馬飲饍之費〕，出处为 《元史》卷038〈至元元年〉。材料未将地点细化到城址，不能用中枢所在地替它补足。

本纪年条记录政令或机构处置；实施背景须参照制度材料。 记录所见的结果则是：可在同年及后续政令中追踪；条文不单独证明地方执行。 日期相邻不等于因果已经证实。

证据的边界是：《元史》为明初仓促官修，本纪保存大量实录条目但有编次与纪年问题；行省执行、数字、族称及对外关系须以条格、多语文书和对方史料互校。 蒙古征服、草原政治与跨区域文献比较研究 可用于继续检验这一结论。

\noindent\eventanchor{SYD-1264-MONGOL-259}{丙申，中書省臣言：「甘肅甘州路十字寺奉安世祖皇帝母別吉太后於內，請定祭禮}\textbf{蒙古。}
从 《元史》卷038〈至元元年〉 的 1264年（元至元1年；公元换算据《元史》本纪纪年） 条目看，〔丙申，中書省臣言：「甘肅甘州路十字寺奉安世祖皇帝母別吉太后於內，請定祭禮〕使 蒙古 的一个具体行动进入叙事；题名保留的城、路、水道或机构名称，是目前可以落地的空间线索。

把本项放回事件链，能够看到的前提为：本纪年条记录政令或机构处置；实施背景须参照制度材料。 而它所打开或收紧的下一步是：可在同年及后续政令中追踪；条文不单独证明地方执行。

《元史》为明初仓促官修，本纪保存大量实录条目但有编次与纪年问题；行省执行、数字、族称及对外关系须以条格、多语文书和对方史料互校。 因此，蒙古征服、草原政治与跨区域文献比较研究 只能扩展问题，不能代替未保存的细节。

\subsection{1264年}
\noindent\eventanchor{YUAN-1264-EVENT-170}{阿里不哥失败，忽必烈结束主要继承战争}\textbf{蒙古。}
1264年 的记载将 蒙古 与〔阿里不哥失败，忽必烈结束主要继承战争〕相连。此处可确认的是这项被记录的行动；材料未将地点细化到城址，不能用中枢所在地替它补足。

它承接的条件是：忽必烈政权、征宋和元朝建制中既有的联盟、交通与补给压力，为“阿里不哥失败，忽必烈结束主要继承战争”提供了直接的军事或动员背景。 随后可追踪的变化为：“阿里不哥失败，忽必烈结束主要继承战争”改变了相关城镇、道路或政权间的军事选择；伤亡、俘获和控制范围仅可按各方材料分别确认。

关于可说范围，须保留：《元史》明初速修，纪年与人名有异同；蒙古大汗政权不等同所有汗国，征服、赋役和身份分类须按地区及材料语言分列。 进一步讨论可参照 蒙古帝国、元朝行省财政、多语文书与区域社会研究

\subsection{1265 年（共享年序桥接）}
\noindent\eventanchor{SY-1265-CHRONOLOGY-BRIDGE}{【C级年序桥接】保留1265 年的多政权并行检索入口；本条不主张所列政权在当年发生直接互动，具体事件须另以单项材料入账。}\textbf{南宋、金、西夏与蒙古（年序桥接）。}
〔【C级年序桥接】保留1265 年的多政权并行检索入口；本条不主张所列政权在当年发生直接互动，具体事件须另以单项材料入账。〕见于 《宋史》；《金史》；《元史》；蒙古文、波斯文及地方材料（据事核） 的 1265 年（共享年序桥接） 记录。南宋、金、西夏与蒙古（年序桥接） 的选择因此有了可回查的时间位置；题名保留的城、路、水道或机构名称，是目前可以落地的空间线索。

前一层压力可由以下线索辨认：相邻已入账事件之间缺少该年度的共享年序入口，易使跨政权材料被单一政权叙事遮蔽。 这一处置留下的后续条件是：保留该年度的检索与补账位置；后续仅在获得可核单项材料时拆分为具体政治、财政、军事、社会或文化事件。

证据的边界是：本条只确认该年位于可追踪的共享年序中；正史、地方文书和外部材料的保存密度不一，不能由桥接标签推断行动、统属、疆界或社会影响。 蒙古征服、区域政权转换与跨语种年序研究 可用于继续检验这一结论。

\subsection{1265年（元至元2年；公元换算据《元史》本纪纪年）}
\noindent\eventanchor{SYD-1265-MONGOL-048}{敕上都商稅、酒醋諸課毋徵，其榷鹽仍舊}\textbf{蒙古。}
1265年（元至元2年；公元换算据《元史》本纪纪年），《元史》卷006〈本纪第6卷本年条〉 所见的〔敕上都商稅、酒醋諸課毋徵，其榷鹽仍舊〕把 蒙古 的处置留在可复核的年序中；题名保留的城、路、水道或机构名称，是目前可以落地的空间线索。

本纪所记财用、赈济或征发属于特定报告场景。 记录所见的结果则是：可与食货志、文书和物质材料比较，不外推为全境总量。 日期相邻不等于因果已经证实。

《元史》为明初仓促官修，本纪保存大量实录条目但有编次与纪年问题；行省执行、数字、族称及对外关系须以条格、多语文书和对方史料互校。 因此，蒙古征服、草原政治与跨区域文献比较研究 只能扩展问题，不能代替未保存的细节。

\noindent\eventanchor{SYD-1265-MONGOL-098}{又請罷北京行中書省，別立宣慰司以控制東北州郡}\textbf{蒙古。}
蒙古 在 1265年（元至元2年；公元换算据《元史》本纪纪年） 的材料痕迹是〔又請罷北京行中書省，別立宣慰司以控制東北州郡〕，出处为 《元史》卷006〈本纪第6卷本年条〉。题名保留的城、路、水道或机构名称，是目前可以落地的空间线索。

把本项放回事件链，能够看到的前提为：本纪年条记录政令或机构处置；实施背景须参照制度材料。 而它所打开或收紧的下一步是：可在同年及后续政令中追踪；条文不单独证明地方执行。

关于可说范围，须保留：《元史》为明初仓促官修，本纪保存大量实录条目但有编次与纪年问题；行省执行、数字、族称及对外关系须以条格、多语文书和对方史料互校。 进一步讨论可参照 蒙古征服、草原政治与跨区域文献比较研究

\noindent\eventanchor{SYD-1265-MONGOL-145}{詔以宏嘗告李璮反，免宏死罪，罷其職，徵贓物償官}\textbf{蒙古。}
从 《元史》卷006〈本纪第6卷本年条〉 的 1265年（元至元2年；公元换算据《元史》本纪纪年） 条目看，〔詔以宏嘗告李璮反，免宏死罪，罷其職，徵贓物償官〕使 蒙古 的一个具体行动进入叙事；材料未将地点细化到城址，不能用中枢所在地替它补足。

它承接的条件是：本纪年条记录政令或机构处置；实施背景须参照制度材料。 随后可追踪的变化为：可在同年及后续政令中追踪；条文不单独证明地方执行。

证据的边界是：《元史》为明初仓促官修，本纪保存大量实录条目但有编次与纪年问题；行省执行、数字、族称及对外关系须以条格、多语文书和对方史料互校。 蒙古征服、草原政治与跨区域文献比较研究 可用于继续检验这一结论。

\noindent\eventanchor{SYD-1265-MONGOL-187}{丙辰，雅州碉門宣撫使請復碉門城邑，詔相度之}\textbf{蒙古。}
1265年（元至元2年；公元换算据《元史》本纪纪年） 的记载将 蒙古 与〔丙辰，雅州碉門宣撫使請復碉門城邑，詔相度之〕相连。此处可确认的是这项被记录的行动；题名保留的城、路、水道或机构名称，是目前可以落地的空间线索。

前一层压力可由以下线索辨认：本纪年条记录政令或机构处置；实施背景须参照制度材料。 这一处置留下的后续条件是：可在同年及后续政令中追踪；条文不单独证明地方执行。

《元史》为明初仓促官修，本纪保存大量实录条目但有编次与纪年问题；行省执行、数字、族称及对外关系须以条格、多语文书和对方史料互校。 因此，蒙古征服、草原政治与跨区域文献比较研究 只能扩展问题，不能代替未保存的细节。

\noindent\eventanchor{SYD-1265-MONGOL-227}{丙寅，命四川行院分兵屯田}\textbf{蒙古。}
〔丙寅，命四川行院分兵屯田〕见于 《元史》卷006〈本纪第6卷本年条〉 的 1265年（元至元2年；公元换算据《元史》本纪纪年） 记录。蒙古 的选择因此有了可回查的时间位置；题名保留的城、路、水道或机构名称，是目前可以落地的空间线索。

本纪年条记录政令或机构处置；实施背景须参照制度材料。 记录所见的结果则是：可在同年及后续政令中追踪；条文不单独证明地方执行。 日期相邻不等于因果已经证实。

关于可说范围，须保留：《元史》为明初仓促官修，本纪保存大量实录条目但有编次与纪年问题；行省执行、数字、族称及对外关系须以条格、多语文书和对方史料互校。 进一步讨论可参照 蒙古征服、草原政治与跨区域文献比较研究

\noindent\eventanchor{SYD-1265-MONGOL-260}{敕徙鎮海、百〔八里〕、謙謙州諸色匠戶於中都，給銀萬五千兩為行費}\textbf{蒙古。}
1265年（元至元2年；公元换算据《元史》本纪纪年），《元史》卷006〈本纪第6卷本年条〉 所见的〔敕徙鎮海、百〔八里〕、謙謙州諸色匠戶於中都，給銀萬五千兩為行費〕把 蒙古 的处置留在可复核的年序中；题名保留的城、路、水道或机构名称，是目前可以落地的空间线索。

把本项放回事件链，能够看到的前提为：本纪所记财用、赈济或征发属于特定报告场景。 而它所打开或收紧的下一步是：可与食货志、文书和物质材料比较，不外推为全境总量。

证据的边界是：《元史》为明初仓促官修，本纪保存大量实录条目但有编次与纪年问题；行省执行、数字、族称及对外关系须以条格、多语文书和对方史料互校。 蒙古征服、草原政治与跨区域文献比较研究 可用于继续检验这一结论。

\subsection{1266 年（共享年序桥接）}
\noindent\eventanchor{SY-1266-CHRONOLOGY-BRIDGE}{【C级年序桥接】保留1266 年的多政权并行检索入口；本条不主张所列政权在当年发生直接互动，具体事件须另以单项材料入账。}\textbf{南宋、金、西夏与蒙古（年序桥接）。}
南宋、金、西夏与蒙古（年序桥接） 在 1266 年（共享年序桥接） 的材料痕迹是〔【C级年序桥接】保留1266 年的多政权并行检索入口；本条不主张所列政权在当年发生直接互动，具体事件须另以单项材料入账。〕，出处为 《宋史》；《金史》；《元史》；蒙古文、波斯文及地方材料（据事核）。题名保留的城、路、水道或机构名称，是目前可以落地的空间线索。

它承接的条件是：相邻已入账事件之间缺少该年度的共享年序入口，易使跨政权材料被单一政权叙事遮蔽。 随后可追踪的变化为：保留该年度的检索与补账位置；后续仅在获得可核单项材料时拆分为具体政治、财政、军事、社会或文化事件。

本条只确认该年位于可追踪的共享年序中；正史、地方文书和外部材料的保存密度不一，不能由桥接标签推断行动、统属、疆界或社会影响。 因此，蒙古征服、区域政权转换与跨语种年序研究 只能扩展问题，不能代替未保存的细节。

\subsection{1266年（元至元3年；公元换算据《元史》本纪纪年）}
\noindent\eventanchor{SYD-1266-MONGOL-049}{命平章政事趙璧等集羣臣議，定為八室}\textbf{蒙古。}
从 《元史》卷006〈本纪第6卷本年条〉 的 1266年（元至元3年；公元换算据《元史》本纪纪年） 条目看，〔命平章政事趙璧等集羣臣議，定為八室〕使 蒙古 的一个具体行动进入叙事；材料未将地点细化到城址，不能用中枢所在地替它补足。

前一层压力可由以下线索辨认：本纪年条记录政令或机构处置；实施背景须参照制度材料。 这一处置留下的后续条件是：可在同年及后续政令中追踪；条文不单独证明地方执行。

关于可说范围，须保留：《元史》为明初仓促官修，本纪保存大量实录条目但有编次与纪年问题；行省执行、数字、族称及对外关系须以条格、多语文书和对方史料互校。 进一步讨论可参照 蒙古征服、草原政治与跨区域文献比较研究

\noindent\eventanchor{SYD-1266-MONGOL-099}{又詔高麗導去使至其國}\textbf{蒙古。}
1266年（元至元3年；公元换算据《元史》本纪纪年） 的记载将 蒙古 与〔又詔高麗導去使至其國〕相连。此处可确认的是这项被记录的行动；材料未将地点细化到城址，不能用中枢所在地替它补足。

本纪年条记录政令或机构处置；实施背景须参照制度材料。 记录所见的结果则是：可在同年及后续政令中追踪；条文不单独证明地方执行。 日期相邻不等于因果已经证实。

证据的边界是：《元史》为明初仓促官修，本纪保存大量实录条目但有编次与纪年问题；行省执行、数字、族称及对外关系须以条格、多语文书和对方史料互校。 蒙古征服、草原政治与跨区域文献比较研究 可用于继续检验这一结论。

\noindent\eventanchor{SYD-1266-MONGOL-146}{八月癸亥，賜丞相伯顏第一區}\textbf{蒙古。}
〔八月癸亥，賜丞相伯顏第一區〕见于 《元史》卷006〈本纪第6卷本年条〉 的 1266年（元至元3年；公元换算据《元史》本纪纪年） 记录。蒙古 的选择因此有了可回查的时间位置；材料未将地点细化到城址，不能用中枢所在地替它补足。

把本项放回事件链，能够看到的前提为：本纪年条提供日期和中枢可见行动；前因不据单条补定。 而它所打开或收紧的下一步是：可回查同年及后续条目，地方影响仍须另证。

《元史》为明初仓促官修，本纪保存大量实录条目但有编次与纪年问题；行省执行、数字、族称及对外关系须以条格、多语文书和对方史料互校。 因此，蒙古征服、草原政治与跨区域文献比较研究 只能扩展问题，不能代替未保存的细节。

\noindent\eventanchor{SYD-1266-MONGOL-188}{丙辰，千戶散竹帶以嗜酒失所守大良平，罪當死，錄其前功免死，令往東川軍前自效}\textbf{蒙古。}
1266年（元至元3年；公元换算据《元史》本纪纪年），《元史》卷006〈本纪第6卷本年条〉 所见的〔丙辰，千戶散竹帶以嗜酒失所守大良平，罪當死，錄其前功免死，令往東川軍前自效〕把 蒙古 的处置留在可复核的年序中；题名保留的城、路、水道或机构名称，是目前可以落地的空间线索。

它承接的条件是：本纪年条提供日期和中枢可见行动；前因不据单条补定。 随后可追踪的变化为：可回查同年及后续条目，地方影响仍须另证。

关于可说范围，须保留：《元史》为明初仓促官修，本纪保存大量实录条目但有编次与纪年问题；行省执行、数字、族称及对外关系须以条格、多语文书和对方史料互校。 进一步讨论可参照 蒙古征服、草原政治与跨区域文献比较研究

\noindent\eventanchor{SYD-1266-MONGOL-228}{丙午，浚西夏中興漢延、唐來等渠}\textbf{蒙古与西夏。}
蒙古与西夏 在 1266年（元至元3年；公元换算据《元史》本纪纪年） 的材料痕迹是〔丙午，浚西夏中興漢延、唐來等渠〕，出处为 《元史》卷006〈本纪第6卷本年条〉。材料未将地点细化到城址，不能用中枢所在地替它补足。

前一层压力可由以下线索辨认：本纪年条提供日期和中枢可见行动；前因不据单条补定。 这一处置留下的后续条件是：可回查同年及后续条目，地方影响仍须另证。

证据的边界是：《元史》为明初仓促官修，本纪保存大量实录条目但有编次与纪年问题；行省执行、数字、族称及对外关系须以条格、多语文书和对方史料互校。 蒙古征服、草原政治与跨区域文献比较研究 可用于继续检验这一结论。

\noindent\eventanchor{SYD-1266-MONGOL-261}{丙午，遣朵端、趙璧持詔撫諭四川將吏軍民}\textbf{蒙古。}
从 《元史》卷006〈本纪第6卷本年条〉 的 1266年（元至元3年；公元换算据《元史》本纪纪年） 条目看，〔丙午，遣朵端、趙璧持詔撫諭四川將吏軍民〕使 蒙古 的一个具体行动进入叙事；题名保留的城、路、水道或机构名称，是目前可以落地的空间线索。

本纪年条记录政令或机构处置；实施背景须参照制度材料。 记录所见的结果则是：可在同年及后续政令中追踪；条文不单独证明地方执行。 日期相邻不等于因果已经证实。

《元史》为明初仓促官修，本纪保存大量实录条目但有编次与纪年问题；行省执行、数字、族称及对外关系须以条格、多语文书和对方史料互校。 因此，蒙古征服、草原政治与跨区域文献比较研究 只能扩展问题，不能代替未保存的细节。

\subsection{1267年（元至元4年；公元换算据《元史》本纪纪年）}
\noindent\eventanchor{SYD-1267-MONGOL-050}{詔以安童為長，史天澤次之，其餘蒙古、漢人參用，勿令員數過多}\textbf{蒙古。}
1267年（元至元4年；公元换算据《元史》本纪纪年） 的记载将 蒙古 与〔詔以安童為長，史天澤次之，其餘蒙古、漢人參用，勿令員數過多〕相连。此处可确认的是这项被记录的行动；材料未将地点细化到城址，不能用中枢所在地替它补足。

把本项放回事件链，能够看到的前提为：本纪年条记录政令或机构处置；实施背景须参照制度材料。 而它所打开或收紧的下一步是：可在同年及后续政令中追踪；条文不单独证明地方执行。

关于可说范围，须保留：《元史》为明初仓促官修，本纪保存大量实录条目但有编次与纪年问题；行省执行、数字、族称及对外关系须以条格、多语文书和对方史料互校。 进一步讨论可参照 蒙古征服、草原政治与跨区域文献比较研究

\noindent\eventanchor{SYD-1267-MONGOL-100}{安南國王陳光昞遣使來貢，優詔答之}\textbf{蒙古。}
〔安南國王陳光昞遣使來貢，優詔答之〕见于 《元史》卷006〈本纪第6卷本年条〉 的 1267年（元至元4年；公元换算据《元史》本纪纪年） 记录。蒙古 的选择因此有了可回查的时间位置；材料未将地点细化到城址，不能用中枢所在地替它补足。

它承接的条件是：本纪年条记录政令或机构处置；实施背景须参照制度材料。 随后可追踪的变化为：可在同年及后续政令中追踪；条文不单独证明地方执行。

证据的边界是：《元史》为明初仓促官修，本纪保存大量实录条目但有编次与纪年问题；行省执行、数字、族称及对外关系须以条格、多语文书和对方史料互校。 蒙古征服、草原政治与跨区域文献比较研究 可用于继续检验这一结论。

\noindent\eventanchor{SYD-1267-MONGOL-147}{罷懷孟路安撫李宗傑，以其民隸本路}\textbf{蒙古。}
1267年（元至元4年；公元换算据《元史》本纪纪年），《元史》卷006〈本纪第6卷本年条〉 所见的〔罷懷孟路安撫李宗傑，以其民隸本路〕把 蒙古 的处置留在可复核的年序中；题名保留的城、路、水道或机构名称，是目前可以落地的空间线索。

前一层压力可由以下线索辨认：本纪年条记录政令或机构处置；实施背景须参照制度材料。 这一处置留下的后续条件是：可在同年及后续政令中追踪；条文不单独证明地方执行。

《元史》为明初仓促官修，本纪保存大量实录条目但有编次与纪年问题；行省执行、数字、族称及对外关系须以条格、多语文书和对方史料互校。 因此，蒙古征服、草原政治与跨区域文献比较研究 只能扩展问题，不能代替未保存的细节。

\noindent\eventanchor{SYD-1267-MONGOL-189}{丙申，威州山後大番弄麻等十一族來附，賜以璽書、金銀符}\textbf{蒙古。}
蒙古 在 1267年（元至元4年；公元换算据《元史》本纪纪年） 的材料痕迹是〔丙申，威州山後大番弄麻等十一族來附，賜以璽書、金銀符〕，出处为 《元史》卷006〈本纪第6卷本年条〉。题名保留的城、路、水道或机构名称，是目前可以落地的空间线索。

本纪年条提供日期和中枢可见行动；前因不据单条补定。 记录所见的结果则是：可回查同年及后续条目，地方影响仍须另证。 日期相邻不等于因果已经证实。

关于可说范围，须保留：《元史》为明初仓促官修，本纪保存大量实录条目但有编次与纪年问题；行省执行、数字、族称及对外关系须以条格、多语文书和对方史料互校。 进一步讨论可参照 蒙古征服、草原政治与跨区域文献比较研究

\noindent\eventanchor{SYD-1267-MONGOL-229}{丙寅，復立宣徽院，以前中書右丞相線真為使}\textbf{蒙古。}
从 《元史》卷006〈本纪第6卷本年条〉 的 1267年（元至元4年；公元换算据《元史》本纪纪年） 条目看，〔丙寅，復立宣徽院，以前中書右丞相線真為使〕使 蒙古 的一个具体行动进入叙事；材料未将地点细化到城址，不能用中枢所在地替它补足。

把本项放回事件链，能够看到的前提为：本纪年条记录政令或机构处置；实施背景须参照制度材料。 而它所打开或收紧的下一步是：可在同年及后续政令中追踪；条文不单独证明地方执行。

证据的边界是：《元史》为明初仓促官修，本纪保存大量实录条目但有编次与纪年问题；行省执行、数字、族称及对外关系须以条格、多语文书和对方史料互校。 蒙古征服、草原政治与跨区域文献比较研究 可用于继续检验这一结论。

\noindent\eventanchor{SYD-1267-MONGOL-262}{丙子，賑親王移相哥所部饑民}\textbf{蒙古。}
1267年（元至元4年；公元换算据《元史》本纪纪年） 的记载将 蒙古 与〔丙子，賑親王移相哥所部饑民〕相连。此处可确认的是这项被记录的行动；材料未将地点细化到城址，不能用中枢所在地替它补足。

它承接的条件是：本纪所记财用、赈济或征发属于特定报告场景。 随后可追踪的变化为：可与食货志、文书和物质材料比较，不外推为全境总量。

《元史》为明初仓促官修，本纪保存大量实录条目但有编次与纪年问题；行省执行、数字、族称及对外关系须以条格、多语文书和对方史料互校。 因此，蒙古征服、草原政治与跨区域文献比较研究 只能扩展问题，不能代替未保存的细节。

\subsection{1267年}
\noindent\eventanchor{YUAN-1267-EVENT-171}{忽必烈营建大都}\textbf{蒙古。}
〔忽必烈营建大都〕见于 《元史》相关本纪；《元朝秘史》、波斯文史籍或《高丽史》对应叙事；《元典章》《通制条格》、多语碑刻与文书（据事核） 的 1267年 记录。蒙古 的选择因此有了可回查的时间位置；题名保留的城、路、水道或机构名称，是目前可以落地的空间线索。

前一层压力可由以下线索辨认：忽必烈政权、征宋和元朝建制对中枢资源、交通或行政协同的要求，推动了“忽必烈营建大都”这一制度或空间安排。 这一处置留下的后续条件是：“忽必烈营建大都”提供新的治理框架或资源通道，但颁行、复申与不同地区的实际执行须分开核验。

关于可说范围，须保留：《元史》明初速修，纪年与人名有异同；蒙古大汗政权不等同所有汗国，征服、赋役和身份分类须按地区及材料语言分列。 进一步讨论可参照 蒙古帝国、元朝行省财政、多语文书与区域社会研究

\subsection{1268年（元至元5年；公元换算据《元史》本纪纪年）}
\noindent\eventanchor{SYD-1268-MONGOL-051}{阿朮言：「所領者蒙古軍，若遇山水、寨柵，非漢軍不可}\textbf{蒙古。}
1268年（元至元5年；公元换算据《元史》本纪纪年），《元史》卷006〈本纪第6卷本年条〉 所见的〔阿朮言：「所領者蒙古軍，若遇山水、寨柵，非漢軍不可〕把 蒙古 的处置留在可复核的年序中；题名保留的城、路、水道或机构名称，是目前可以落地的空间线索。

本纪保留灾害或工程的报告地点与朝廷处置；灾情范围需另证。 记录所见的结果则是：可与地方文书、河工和环境材料比较，不将单点记录外推为全域因果。 日期相邻不等于因果已经证实。

证据的边界是：《元史》为明初仓促官修，本纪保存大量实录条目但有编次与纪年问题；行省执行、数字、族称及对外关系须以条格、多语文书和对方史料互校。 蒙古征服、草原政治与跨区域文献比较研究 可用于继续检验这一结论。

\noindent\eventanchor{SYD-1268-MONGOL-101}{罷各路奧魯官，令管民官兼領}\textbf{蒙古。}
蒙古 在 1268年（元至元5年；公元换算据《元史》本纪纪年） 的材料痕迹是〔罷各路奧魯官，令管民官兼領〕，出处为 《元史》卷006〈本纪第6卷本年条〉。题名保留的城、路、水道或机构名称，是目前可以落地的空间线索。

把本项放回事件链，能够看到的前提为：本纪年条记录政令或机构处置；实施背景须参照制度材料。 而它所打开或收紧的下一步是：可在同年及后续政令中追踪；条文不单独证明地方执行。

《元史》为明初仓促官修，本纪保存大量实录条目但有编次与纪年问题；行省执行、数字、族称及对外关系须以条格、多语文书和对方史料互校。 因此，蒙古征服、草原政治与跨区域文献比较研究 只能扩展问题，不能代替未保存的细节。

\noindent\eventanchor{SYD-1268-MONGOL-148}{罷諸路女直、契丹、漢人為達魯花赤者，回回、畏兀、乃蠻、唐兀人仍舊}\textbf{蒙古。}
从 《元史》卷006〈本纪第6卷本年条〉 的 1268年（元至元5年；公元换算据《元史》本纪纪年） 条目看，〔罷諸路女直、契丹、漢人為達魯花赤者，回回、畏兀、乃蠻、唐兀人仍舊〕使 蒙古 的一个具体行动进入叙事；题名保留的城、路、水道或机构名称，是目前可以落地的空间线索。

它承接的条件是：本纪年条记录政令或机构处置；实施背景须参照制度材料。 随后可追踪的变化为：可在同年及后续政令中追踪；条文不单独证明地方执行。

关于可说范围，须保留：《元史》为明初仓促官修，本纪保存大量实录条目但有编次与纪年问题；行省执行、数字、族称及对外关系须以条格、多语文书和对方史料互校。 进一步讨论可参照 蒙古征服、草原政治与跨区域文献比较研究

\noindent\eventanchor{SYD-1268-MONGOL-190}{敕給黎、雅、嘉定新附民田}\textbf{蒙古。}
1268年（元至元5年；公元换算据《元史》本纪纪年） 的记载将 蒙古 与〔敕給黎、雅、嘉定新附民田〕相连。此处可确认的是这项被记录的行动；材料未将地点细化到城址，不能用中枢所在地替它补足。

前一层压力可由以下线索辨认：本纪所记财用、赈济或征发属于特定报告场景。 这一处置留下的后续条件是：可与食货志、文书和物质材料比较，不外推为全境总量。

证据的边界是：《元史》为明初仓促官修，本纪保存大量实录条目但有编次与纪年问题；行省执行、数字、族称及对外关系须以条格、多语文书和对方史料互校。 蒙古征服、草原政治与跨区域文献比较研究 可用于继续检验这一结论。

\noindent\eventanchor{SYD-1268-MONGOL-230}{敕長春宮修設金籙周天大醮七晝夜}\textbf{蒙古。}
〔敕長春宮修設金籙周天大醮七晝夜〕见于 《元史》卷006〈本纪第6卷本年条〉 的 1268年（元至元5年；公元换算据《元史》本纪纪年） 记录。蒙古 的选择因此有了可回查的时间位置；题名保留的城、路、水道或机构名称，是目前可以落地的空间线索。

本纪年条提供日期和中枢可见行动；前因不据单条补定。 记录所见的结果则是：可回查同年及后续条目，地方影响仍须另证。 日期相邻不等于因果已经证实。

《元史》为明初仓促官修，本纪保存大量实录条目但有编次与纪年问题；行省执行、数字、族称及对外关系须以条格、多语文书和对方史料互校。 因此，蒙古征服、草原政治与跨区域文献比较研究 只能扩展问题，不能代替未保存的细节。

\noindent\eventanchor{SYD-1268-MONGOL-263}{賜諸王金、銀、幣、帛如歲例}\textbf{蒙古。}
1268年（元至元5年；公元换算据《元史》本纪纪年），《元史》卷006〈本纪第6卷本年条〉 所见的〔賜諸王金、銀、幣、帛如歲例〕把 蒙古 的处置留在可复核的年序中；材料未将地点细化到城址，不能用中枢所在地替它补足。

把本项放回事件链，能够看到的前提为：本纪年条提供日期和中枢可见行动；前因不据单条补定。 而它所打开或收紧的下一步是：可回查同年及后续条目，地方影响仍须另证。

关于可说范围，须保留：《元史》为明初仓促官修，本纪保存大量实录条目但有编次与纪年问题；行省执行、数字、族称及对外关系须以条格、多语文书和对方史料互校。 进一步讨论可参照 蒙古征服、草原政治与跨区域文献比较研究

\subsection{1268年}
\noindent\eventanchor{YUAN-1268-EVENT-172}{蒙古军围攻襄阳、樊城}\textbf{蒙古与南宋。}
蒙古与南宋 在 1268年 的材料痕迹是〔蒙古军围攻襄阳、樊城〕，出处为 《元史》相关本纪；《元朝秘史》、波斯文史籍或《高丽史》对应叙事；《元典章》《通制条格》、多语碑刻与文书（据事核）。题名保留的城、路、水道或机构名称，是目前可以落地的空间线索。

它承接的条件是：忽必烈政权、征宋和元朝建制中既有的联盟、交通与补给压力，为“蒙古军围攻襄阳、樊城”提供了直接的军事或动员背景。 随后可追踪的变化为：“蒙古军围攻襄阳、樊城”改变了相关城镇、道路或政权间的军事选择；伤亡、俘获和控制范围仅可按各方材料分别确认。

证据的边界是：《元史》明初速修，纪年与人名有异同；蒙古大汗政权不等同所有汗国，征服、赋役和身份分类须按地区及材料语言分列。 蒙古帝国、元朝行省财政、多语文书与区域社会研究 可用于继续检验这一结论。

\subsection{1269 年（共享年序桥接）}
\noindent\eventanchor{SY-1269-CHRONOLOGY-BRIDGE}{【C级年序桥接】保留1269 年的多政权并行检索入口；本条不主张所列政权在当年发生直接互动，具体事件须另以单项材料入账。}\textbf{南宋、金、西夏与蒙古（年序桥接）。}
从 《宋史》；《金史》；《元史》；蒙古文、波斯文及地方材料（据事核） 的 1269 年（共享年序桥接） 条目看，〔【C级年序桥接】保留1269 年的多政权并行检索入口；本条不主张所列政权在当年发生直接互动，具体事件须另以单项材料入账。〕使 南宋、金、西夏与蒙古（年序桥接） 的一个具体行动进入叙事；题名保留的城、路、水道或机构名称，是目前可以落地的空间线索。

前一层压力可由以下线索辨认：相邻已入账事件之间缺少该年度的共享年序入口，易使跨政权材料被单一政权叙事遮蔽。 这一处置留下的后续条件是：保留该年度的检索与补账位置；后续仅在获得可核单项材料时拆分为具体政治、财政、军事、社会或文化事件。

本条只确认该年位于可追踪的共享年序中；正史、地方文书和外部材料的保存密度不一，不能由桥接标签推断行动、统属、疆界或社会影响。 因此，蒙古征服、区域政权转换与跨语种年序研究 只能扩展问题，不能代替未保存的细节。

\subsection{1269年（元至元6年；公元换算据《元史》本纪纪年）}
\noindent\eventanchor{SYD-1269-MONGOL-052}{安南國王陳光昞遣使來貢}\textbf{蒙古。}
1269年（元至元6年；公元换算据《元史》本纪纪年） 的记载将 蒙古 与〔安南國王陳光昞遣使來貢〕相连。此处可确认的是这项被记录的行动；材料未将地点细化到城址，不能用中枢所在地替它补足。

本纪记录使节、贡赐或往来，不等于稳定统属关系。 记录所见的结果则是：可与对方编年和交通、文书材料核对其后续落实。 日期相邻不等于因果已经证实。

关于可说范围，须保留：《元史》为明初仓促官修，本纪保存大量实录条目但有编次与纪年问题；行省执行、数字、族称及对外关系须以条格、多语文书和对方史料互校。 进一步讨论可参照 蒙古征服、草原政治与跨区域文献比较研究

\noindent\eventanchor{SYD-1269-MONGOL-102}{丙申，高麗國王王禃遣其世子愖來朝，賜禃玉帶一，愖金五十兩，從官銀幣有差}\textbf{蒙古。}
〔丙申，高麗國王王禃遣其世子愖來朝，賜禃玉帶一，愖金五十兩，從官銀幣有差〕见于 《元史》卷006〈本纪第6卷本年条〉 的 1269年（元至元6年；公元换算据《元史》本纪纪年） 记录。蒙古 的选择因此有了可回查的时间位置；材料未将地点细化到城址，不能用中枢所在地替它补足。

把本项放回事件链，能够看到的前提为：本纪记录使节、贡赐或往来，不等于稳定统属关系。 而它所打开或收紧的下一步是：可与对方编年和交通、文书材料核对其后续落实。

证据的边界是：《元史》为明初仓促官修，本纪保存大量实录条目但有编次与纪年问题；行省执行、数字、族称及对外关系须以条格、多语文书和对方史料互校。 蒙古征服、草原政治与跨区域文献比较研究 可用于继续检验这一结论。

\noindent\eventanchor{SYD-1269-MONGOL-149}{開元等路饑，減戶賦布二匹，秋稅減其半，水達達戶減青鼠二，其租稅被災者免徵}\textbf{蒙古。}
1269年（元至元6年；公元换算据《元史》本纪纪年），《元史》卷006〈本纪第6卷本年条〉 所见的〔開元等路饑，減戶賦布二匹，秋稅減其半，水達達戶減青鼠二，其租稅被災者免徵〕把 蒙古 的处置留在可复核的年序中；题名保留的城、路、水道或机构名称，是目前可以落地的空间线索。

它承接的条件是：本纪所记财用、赈济或征发属于特定报告场景。 随后可追踪的变化为：可与食货志、文书和物质材料比较，不外推为全境总量。

《元史》为明初仓促官修，本纪保存大量实录条目但有编次与纪年问题；行省执行、数字、族称及对外关系须以条格、多语文书和对方史料互校。 因此，蒙古征服、草原政治与跨区域文献比较研究 只能扩展问题，不能代替未保存的细节。

\noindent\eventanchor{SYD-1269-MONGOL-191}{八月己卯，立金州招討司}\textbf{蒙古。}
蒙古 在 1269年（元至元6年；公元换算据《元史》本纪纪年） 的材料痕迹是〔八月己卯，立金州招討司〕，出处为 《元史》卷006〈本纪第6卷本年条〉。题名保留的城、路、水道或机构名称，是目前可以落地的空间线索。

前一层压力可由以下线索辨认：本纪年条提供日期和中枢可见行动；前因不据单条补定。 这一处置留下的后续条件是：可回查同年及后续条目，地方影响仍须另证。

关于可说范围，须保留：《元史》为明初仓促官修，本纪保存大量实录条目但有编次与纪年问题；行省执行、数字、族称及对外关系须以条格、多语文书和对方史料互校。 进一步讨论可参照 蒙古征服、草原政治与跨区域文献比较研究

\noindent\eventanchor{SYD-1269-MONGOL-231}{丙申，罷宣德府稅課所，以上都轉運司兼領}\textbf{蒙古。}
从 《元史》卷006〈本纪第6卷本年条〉 的 1269年（元至元6年；公元换算据《元史》本纪纪年） 条目看，〔丙申，罷宣德府稅課所，以上都轉運司兼領〕使 蒙古 的一个具体行动进入叙事；题名保留的城、路、水道或机构名称，是目前可以落地的空间线索。

本纪年条记录政令或机构处置；实施背景须参照制度材料。 记录所见的结果则是：可在同年及后续政令中追踪；条文不单独证明地方执行。 日期相邻不等于因果已经证实。

证据的边界是：《元史》为明初仓促官修，本纪保存大量实录条目但有编次与纪年问题；行省执行、数字、族称及对外关系须以条格、多语文书和对方史料互校。 蒙古征服、草原政治与跨区域文献比较研究 可用于继续检验这一结论。

\noindent\eventanchor{SYD-1269-MONGOL-264}{丙申，以沙、肅州鈔法未行，降詔諭之}\textbf{蒙古。}
1269年（元至元6年；公元换算据《元史》本纪纪年） 的记载将 蒙古 与〔丙申，以沙、肅州鈔法未行，降詔諭之〕相连。此处可确认的是这项被记录的行动；题名保留的城、路、水道或机构名称，是目前可以落地的空间线索。

把本项放回事件链，能够看到的前提为：本纪年条记录政令或机构处置；实施背景须参照制度材料。 而它所打开或收紧的下一步是：可在同年及后续政令中追踪；条文不单独证明地方执行。

《元史》为明初仓促官修，本纪保存大量实录条目但有编次与纪年问题；行省执行、数字、族称及对外关系须以条格、多语文书和对方史料互校。 因此，蒙古征服、草原政治与跨区域文献比较研究 只能扩展问题，不能代替未保存的细节。

\subsection{1270 年（共享年序桥接）}
\noindent\eventanchor{SY-1270-CHRONOLOGY-BRIDGE}{【C级年序桥接】保留1270 年的多政权并行检索入口；本条不主张所列政权在当年发生直接互动，具体事件须另以单项材料入账。}\textbf{南宋、金、西夏与蒙古（年序桥接）。}
〔【C级年序桥接】保留1270 年的多政权并行检索入口；本条不主张所列政权在当年发生直接互动，具体事件须另以单项材料入账。〕见于 《宋史》；《金史》；《元史》；蒙古文、波斯文及地方材料（据事核） 的 1270 年（共享年序桥接） 记录。南宋、金、西夏与蒙古（年序桥接） 的选择因此有了可回查的时间位置；题名保留的城、路、水道或机构名称，是目前可以落地的空间线索。

它承接的条件是：相邻已入账事件之间缺少该年度的共享年序入口，易使跨政权材料被单一政权叙事遮蔽。 随后可追踪的变化为：保留该年度的检索与补账位置；后续仅在获得可核单项材料时拆分为具体政治、财政、军事、社会或文化事件。

关于可说范围，须保留：本条只确认该年位于可追踪的共享年序中；正史、地方文书和外部材料的保存密度不一，不能由桥接标签推断行动、统属、疆界或社会影响。 进一步讨论可参照 蒙古征服、区域政权转换与跨语种年序研究

\subsection{1272 年（共享年序桥接）}
\noindent\eventanchor{SY-1272-CHRONOLOGY-BRIDGE}{【C级年序桥接】保留1272 年的多政权并行检索入口；本条不主张所列政权在当年发生直接互动，具体事件须另以单项材料入账。}\textbf{南宋、金、西夏与蒙古（年序桥接）。}
1272 年（共享年序桥接），《宋史》；《金史》；《元史》；蒙古文、波斯文及地方材料（据事核） 所见的〔【C级年序桥接】保留1272 年的多政权并行检索入口；本条不主张所列政权在当年发生直接互动，具体事件须另以单项材料入账。〕把 南宋、金、西夏与蒙古（年序桥接） 的处置留在可复核的年序中；题名保留的城、路、水道或机构名称，是目前可以落地的空间线索。

前一层压力可由以下线索辨认：相邻已入账事件之间缺少该年度的共享年序入口，易使跨政权材料被单一政权叙事遮蔽。 这一处置留下的后续条件是：保留该年度的检索与补账位置；后续仅在获得可核单项材料时拆分为具体政治、财政、军事、社会或文化事件。

证据的边界是：本条只确认该年位于可追踪的共享年序中；正史、地方文书和外部材料的保存密度不一，不能由桥接标签推断行动、统属、疆界或社会影响。 蒙古征服、区域政权转换与跨语种年序研究 可用于继续检验这一结论。

\subsection{1278 年（共享年序桥接）}
\noindent\eventanchor{SY-1278-CHRONOLOGY-BRIDGE}{【C级年序桥接】保留1278 年的多政权并行检索入口；本条不主张所列政权在当年发生直接互动，具体事件须另以单项材料入账。}\textbf{南宋、金、西夏与蒙古（年序桥接）。}
南宋、金、西夏与蒙古（年序桥接） 在 1278 年（共享年序桥接） 的材料痕迹是〔【C级年序桥接】保留1278 年的多政权并行检索入口；本条不主张所列政权在当年发生直接互动，具体事件须另以单项材料入账。〕，出处为 《宋史》；《金史》；《元史》；蒙古文、波斯文及地方材料（据事核）。题名保留的城、路、水道或机构名称，是目前可以落地的空间线索。

相邻已入账事件之间缺少该年度的共享年序入口，易使跨政权材料被单一政权叙事遮蔽。 记录所见的结果则是：保留该年度的检索与补账位置；后续仅在获得可核单项材料时拆分为具体政治、财政、军事、社会或文化事件。 日期相邻不等于因果已经证实。

本条只确认该年位于可追踪的共享年序中；正史、地方文书和外部材料的保存密度不一，不能由桥接标签推断行动、统属、疆界或社会影响。 因此，蒙古征服、区域政权转换与跨语种年序研究 只能扩展问题，不能代替未保存的细节。

